% Vectors — Pure Mathematics with Statistics Upper Sixth — Cours signé OnBuch+
% Official MINESEC syllabus. Compiler avec : tectonic PMS03-vectors.tex
\newcommand{\DOCMATIERE}{Pure Mathematics with Statistics}
\newcommand{\DOCNIVEAU}{Upper Sixth}
\newcommand{\DOCMODULE}{Upper Sixth — Pure mathematics with statistics}
\newcommand{\DOCLECON}{3}
\newcommand{\DOCTITRE}{Vectors}
% OnBuch+ — préambule « cours signés » ENRICHI (compilé avec tectonic / XeLaTeX)
% Système visuel « Pop » d'OnBuch+ : fond crème (jamais blanc), bordures encre
% épaisses, ombres « dures » décalées, titres Archivo Black en capitales.
% Version MATHÉMATIQUES : graphes (pgfplots/tikz), boîtes pédagogiques
% (définition, propriété, méthode, attention, le savais-tu, exercice résolu,
% exercice, TP, bilan), page de garde et signature OnBuch+ ancrée en bas.
%
% Macros attendues dans main.tex AVANT \input{preamble} :
%   \DOCMATIERE  \DOCNIVEAU  \DOCMODULE  \DOCLECON (numéro)  \DOCTITRE
\documentclass[11pt]{article}

\usepackage{fontspec}
\usepackage[english]{babel}
\usepackage[a4paper,margin=2cm,top=3.4cm,bottom=2.8cm,headheight=1.4cm,footskip=1.3cm]{geometry}
\usepackage{amsmath,amssymb,amsfonts}
\usepackage{mathtools} % \xmapsto, \coloneqq... souvent utilisés par les modèles
\usepackage{cancel} % simplifications barrées (bilans de forces)
% \abs{z} (module, valeur absolue) et \norm{u} (norme), utilisés par les modèles
\providecommand{\abs}{}\let\abs\relax\DeclarePairedDelimiter{\abs}{\lvert}{\rvert}
\providecommand{\norm}{}\let\norm\relax\DeclarePairedDelimiter{\norm}{\lVert}{\rVert}
\usepackage[table]{xcolor} % [table] : fournit aussi \rowcolors (alternance de lignes)
\usepackage{pagecolor}
\usepackage{tikz}
\usetikzlibrary{arrows.meta,positioning,calc,shapes.geometric,shapes.misc,shapes.arrows,decorations.pathmorphing,decorations.pathreplacing,decorations.markings,patterns,shadows,mindmap,trees,fit,backgrounds,matrix,intersections,angles,quotes}
\usepackage{pgfplots}
\usepackage[version=4]{mhchem} % équations nucléaires \ce{^{235}_{92}U}
\usepackage[european,siunitx]{circuitikz} % schémas électriques
\pgfplotsset{compat=1.18}
\usepgfplotslibrary{fillbetween,groupplots} % aires entre courbes (calcul intégral)
\usepackage{enumitem}
\usepackage{fancyhdr}
% [most] charge toutes les bibliothèques tcolorbox (listings, minted, raster,
% documentation...) et gonfle inutilement la mémoire TeX sur un document long
% (~300 boîtes) : on ne charge que ce qui est réellement utilisé.
\usepackage[skins,breakable]{tcolorbox}
\usepackage{graphicx}
\usepackage{colortbl}
\usepackage{booktabs}
\usepackage{tabularx}
\usepackage{diagbox} % case d'en-tête diagonale des tableaux à double entrée
% Colonnes extensibles centrée / gauche / droite (C, L, R), souvent utilisées par les modèles
\newcolumntype{C}{>{\centering\arraybackslash}X}
\newcolumntype{L}{>{\raggedright\arraybackslash}X}
\newcolumntype{R}{>{\raggedleft\arraybackslash}X}
\usepackage{array}
\usepackage{multicol}
\usepackage{siunitx}
% parse-numbers=false : accepte aussi \SI{2,5 \times 10^{-3}}{...}
\sisetup{locale=UK,output-decimal-marker={.},group-minimum-digits=4,parse-numbers=false}
\DeclareSIUnit\ohm{\ensuremath{\symup{\Omega}}} % U+2126 absent de la police
\DeclareSIUnit\division{div} % réglages d'oscilloscope (ms/div, V/div)
\DeclareSIUnit\tr{tr} % tours (vitesse de rotation, ex. \SI{1500}{\tr\per\minute})
\usepackage{qrcode}
% \degree : commande fréquemment utilisée par les modèles pour un angle ou une
% température en dehors de \SI{}{} (non fournie par les paquets chargés).
\providecommand{\degree}{^{\circ}}
% \qed (fin de démonstration), utilisé par les modèles sans amsthm
\providecommand{\qed}{\hfill$\square$}
% \barre : double barre des valeurs interdites dans un tableau de variations (tkz-tab)
\providecommand{\barre}{\ensuremath{\|}}
% environnement proof (amsthm non chargé) utilisé par les modèles
\ifdefined\proof\else\newenvironment{proof}[1][Proof]{\par\noindent\textit{#1.}\ }{\qed\par}\fi
\usepackage{lastpage}
\usepackage{hyperref}
\hypersetup{hidelinks}

% ============================================================
% Palette « Pop » OnBuch+ — mêmes hex que la classe P (plus_ui.dart)
% ============================================================
\definecolor{poporange}{HTML}{F59321}
\definecolor{popdark}{HTML}{E07A0C}
\definecolor{popcream}{HTML}{FFF6E8}
\definecolor{popink}{HTML}{1C1714}
\definecolor{poppurple}{HTML}{7A5AE0}
\definecolor{popgreen}{HTML}{2E9E6A}
\definecolor{poppink}{HTML}{E0457B}
\definecolor{popblue}{HTML}{2F7DE1}
\definecolor{popgold}{HTML}{C98A12}
\definecolor{popmuted}{HTML}{8A7F76}
\definecolor{popline}{HTML}{EADFD2}
\definecolor{popgray}{HTML}{8A7F76}
\colorlet{popgrey}{popgray}
% Alias tolérants (le modèle nomme parfois les couleurs différemment)
\colorlet{popwhite}{white}
\colorlet{popblack}{popink}
\colorlet{popred}{poppink}
\colorlet{popyellow}{popgold}
% Teintes claires pour fonds de tableaux / zones
\colorlet{popblueL}{popblue!10!white}
\colorlet{poporangeL}{poporange!14!white}
\colorlet{popgreenL}{popgreen!12!white}
\colorlet{poppurpleL}{poppurple!12!white}
\colorlet{poppinkL}{poppink!12!white}
\colorlet{popgoldL}{popgold!14!white}

\pagecolor{popcream}

% ============================================================
% Typographie
% ============================================================
\defaultfontfeatures{Ligatures=TeX}
\setmainfont{PlusJakartaSans-Medium.ttf}[Path=../../../fonts/,BoldFont=PlusJakartaSans-Bold.ttf]
% Maths en sans-serif (Fira Math) avec chiffres et lettres droites de Plus
% Jakarta, pour que les formules se fondent dans le texte.
\usepackage[math-style=ISO,bold-style=ISO]{unicode-math}
\setmathfont{FiraMath-Regular.otf}[Path=../../../fonts/]
\setmathfont{PlusJakartaSans-Medium.ttf}[Path=../../../fonts/,range={up/num,up/{Latin,latin},\mathup}]
% (icomma retiré : décimales avec point en anglais)
% Caractères mathématiques Unicode absents des polices de texte (Plus Jakarta,
% Archivo Black) : rendus par leur équivalent mathématique, en texte comme en
% formule — sinon ils s'affichent en carré vide (ex. « Arithmétique dans ℤ »).
\usepackage{newunicodechar}
\newunicodechar{ℝ}{\ensuremath{\mathbb{R}}}
\newunicodechar{ℤ}{\ensuremath{\mathbb{Z}}}
\newunicodechar{ℂ}{\ensuremath{\mathbb{C}}}
\newunicodechar{ℕ}{\ensuremath{\mathbb{N}}}
\newunicodechar{ℚ}{\ensuremath{\mathbb{Q}}}
\newunicodechar{𝔻}{\ensuremath{\mathbb{D}}}
\newunicodechar{α}{\ensuremath{\alpha}}
\newunicodechar{β}{\ensuremath{\beta}}
\newunicodechar{γ}{\ensuremath{\gamma}}
\newunicodechar{δ}{\ensuremath{\delta}}
\newunicodechar{ε}{\ensuremath{\varepsilon}}
\newunicodechar{θ}{\ensuremath{\theta}}
\newunicodechar{λ}{\ensuremath{\lambda}}
\newunicodechar{μ}{\ensuremath{\mu}}
\newunicodechar{σ}{\ensuremath{\sigma}}
\newunicodechar{τ}{\ensuremath{\tau}}
\newunicodechar{φ}{\ensuremath{\varphi}}
\newunicodechar{ψ}{\ensuremath{\psi}}
\newunicodechar{ω}{\ensuremath{\omega}}
\newunicodechar{Σ}{\ensuremath{\Sigma}}
% majuscules produites par \MakeUppercase dans les titres (α-aminés -> Α)
\newunicodechar{Α}{\ensuremath{\alpha}}
\newunicodechar{Β}{\ensuremath{\beta}}
\newunicodechar{Γ}{\ensuremath{\Gamma}}
\newunicodechar{Δ}{\ensuremath{\Delta}}
\newunicodechar{Ε}{\ensuremath{\varepsilon}}
\newunicodechar{Θ}{\ensuremath{\Theta}}
\newunicodechar{Λ}{\ensuremath{\Lambda}}
\newunicodechar{Μ}{\ensuremath{\mu}}
\newunicodechar{Τ}{\ensuremath{\tau}}
\newunicodechar{Φ}{\ensuremath{\Phi}}
\newunicodechar{Ψ}{\ensuremath{\Psi}}
\newunicodechar{Ω}{\ensuremath{\Omega}}
\newunicodechar{↦}{\ensuremath{\mapsto}}
\newunicodechar{∈}{\ensuremath{\in}}
\newunicodechar{∉}{\ensuremath{\notin}}
\newunicodechar{∧}{\ensuremath{\wedge}}
\newunicodechar{⇔}{\ensuremath{\Leftrightarrow}}
\newunicodechar{⟺}{\ensuremath{\Longleftrightarrow}}
\newunicodechar{⇒}{\ensuremath{\Rightarrow}}
\newunicodechar{⟹}{\ensuremath{\Longrightarrow}}
\newunicodechar{∘}{\ensuremath{\circ}}
\newunicodechar{∪}{\ensuremath{\cup}}
\newunicodechar{∩}{\ensuremath{\cap}}
\newunicodechar{≡}{\ensuremath{\equiv}}
\newunicodechar{⊂}{\ensuremath{\subset}}
\newunicodechar{⊥}{\ensuremath{\perp}}
\newunicodechar{∃}{\ensuremath{\exists}}
\newunicodechar{∀}{\ensuremath{\forall}}
\newunicodechar{∛}{\ensuremath{\sqrt[3]{\,}}}
\newunicodechar{∜}{\ensuremath{\sqrt[4]{\,}}}
\newunicodechar{⃗}{\ensuremath{{}^{\to}}}
\newunicodechar{ⁿ}{\ifmmode^{n}\else\textsuperscript{\ensuremath{n}}\fi}
\newunicodechar{ˣ}{\ifmmode^{x}\else\textsuperscript{\ensuremath{x}}\fi}
\newunicodechar{ᵇ}{\ifmmode^{b}\else\textsuperscript{\ensuremath{b}}\fi}
\newunicodechar{ʳ}{\ifmmode^{r}\else\textsuperscript{\ensuremath{r}}\fi}
\newunicodechar{ᵘ}{\ifmmode^{u}\else\textsuperscript{\ensuremath{u}}\fi}
\newunicodechar{ʸ}{\ifmmode^{y}\else\textsuperscript{\ensuremath{y}}\fi}
\newunicodechar{ᵃ}{\ifmmode^{a}\else\textsuperscript{\ensuremath{a}}\fi}
\newunicodechar{ᵉ}{\ifmmode^{e}\else\textsuperscript{\ensuremath{e}}\fi}
\newunicodechar{ᵏ}{\ifmmode^{k}\else\textsuperscript{\ensuremath{k}}\fi}
\newunicodechar{ᵖ}{\ifmmode^{p}\else\textsuperscript{\ensuremath{p}}\fi}
\newunicodechar{ᴺ}{\ifmmode^{N}\else\textsuperscript{\ensuremath{N}}\fi}
\newunicodechar{ᶜ}{\ifmmode^{c}\else\textsuperscript{\ensuremath{c}}\fi}
\newunicodechar{ʰ}{\ifmmode^{h}\else\textsuperscript{\ensuremath{h}}\fi}
\newunicodechar{ᵠ}{\ifmmode^{q}\else\textsuperscript{\ensuremath{q}}\fi}
\newunicodechar{ₙ}{\ifmmode_{n}\else\textsubscript{\ensuremath{n}}\fi}
\newunicodechar{ₐ}{\ifmmode_{a}\else\textsubscript{\ensuremath{a}}\fi}
\newunicodechar{ᵢ}{\ifmmode_{i}\else\textsubscript{\ensuremath{i}}\fi}
\newunicodechar{ᵣ}{\ifmmode_{r}\else\textsubscript{\ensuremath{r}}\fi}
\newunicodechar{ₖ}{\ifmmode_{k}\else\textsubscript{\ensuremath{k}}\fi}
\newunicodechar{ᵦ}{\ifmmode_{\beta}\else\textsubscript{\ensuremath{\beta}}\fi}
\newunicodechar{ₓ}{\ifmmode_{x}\else\textsubscript{\ensuremath{x}}\fi}
\newfontfamily\popdisplay{ArchivoBlack-Regular.ttf}[Path=../../../fonts/]
\newfontfamily\poplogo{SpaceGrotesk-Bold.ttf}[Path=../../../fonts/]
\newfontfamily\popsemi{PlusJakartaSans-SemiBold.ttf}[Path=../../../fonts/]
\newfontfamily\popxbold{PlusJakartaSans-ExtraBold.ttf}[Path=../../../fonts/]
\newfontfamily\popxboldls{PlusJakartaSans-ExtraBold.ttf}[Path=../../../fonts/,LetterSpace=3.0]

\setlist[enumerate]{leftmargin=1.7em,itemsep=5pt,topsep=4pt}
\setlist[itemize]{leftmargin=1.7em,itemsep=4pt,topsep=4pt,label={\color{poporange}\textbullet}}
\setlength{\parindent}{0pt}
\setlength{\parskip}{5pt}
\renewcommand{\arraystretch}{1.25}

% pgfplots : style Pop par défaut
\pgfplotsset{
  every axis/.append style={axis line style={popink,line width=1pt},tick style={popink},
    grid style={popline,line width=0.5pt},label style={font=\small},tick label style={font=\footnotesize}},
  popcurve/.style={poporange,line width=1.8pt,no markers,smooth},
}
% Même style disponible dans un \draw TikZ simple (hors pgfplots)
\tikzset{popcurve/.style={poporange,line width=1.8pt}}
\tikzset{popgrid/.style={popline,very thin}}
\colorlet{popcurve}{poporange} % le nom du style est parfois utilisé comme couleur

% ============================================================
% Pill / squiggle
% ============================================================
\newcommand{\poppill}[3]{%
  \tikz[baseline=-0.55ex]\node[draw=popink,line width=1.2pt,rounded corners=2.6pt,%
    fill=#2,inner xsep=6.5pt,inner ysep=3pt]{%
    \popxboldls\fontsize{7}{7}\selectfont\color{#3}\MakeUppercase{#1}};%
}
\newcommand{\popsquiggle}[1]{%
  \tikz[baseline=-0.4ex]{
    \draw[#1,line width=2.6pt,line cap=round]
      plot [smooth] coordinates {(0,0.09) (0.34,0.01) (0.68,0.21) (1.02,0.01) (1.36,0.10)};
  }%
}
% Mot-clé surligné (marqueur orange sous le texte)
\newcommand{\cle}[1]{{\popxbold\color{popdark}#1}}

% ============================================================
% En-tête / pied
% ============================================================
\pagestyle{fancy}
\fancyhf{}
\renewcommand{\headrulewidth}{0pt}
\renewcommand{\footrulewidth}{0pt}
\lhead{\raisebox{-0.15ex}{\poplogo\fontsize{13}{13}\selectfont\color{popink}OnBuch\color{poporange}+}}
\rhead{\poppill{\DOCMATIERE\ \textbullet\ \DOCNIVEAU\ \textbullet\ Chapter \DOCLECON}{white}{popink}}
\lfoot{\popxbold\fontsize{7.6}{7.6}\selectfont\color{popmuted}\MakeUppercase{Signed course OnBuch+}\ \textbullet\ {\popsemi\DOCTITRE}}
\rfoot{\tikz[baseline=-0.6ex]\node[rounded corners=8.5pt,fill=popink,minimum height=17pt,minimum width=17pt,inner xsep=5pt,inner sep=0]{\popxbold\fontsize{8}{8}\selectfont\color{popcream}\thepage\,/\,\pageref*{LastPage}};}

% ============================================================
% Titres
% ============================================================
\newcommand{\chaptitre}[2]{%
  \begin{tcolorbox}[enhanced,colback=poporange,colframe=popink,boxrule=2.6pt,arc=11pt,
      drop shadow=popink,left=15pt,right=15pt,top=12pt,bottom=11pt,before skip=3pt,after skip=18pt]
    \poppill{#2}{popink}{popcream}\par\vspace{9pt}
    {\raggedright\hyphenpenalty=10000\exhyphenpenalty=10000\popdisplay\fontsize{22}{24}\selectfont\color{popcream}\MakeUppercase{#1}\par}
  \end{tcolorbox}
}
% Section numérotée : pastille orange avec le numéro + titre Archivo
\newcounter{popsec}
\newcommand{\coursec}[1]{%
  \stepcounter{popsec}\setcounter{popsub}{0}%
  \par\vspace{16pt}\noindent
  \begin{minipage}{\linewidth}
  \tikz[baseline=-0.7ex]\node[draw=popink,line width=1.6pt,fill=poporange,rounded corners=4pt,minimum size=22pt,inner sep=0]{\popdisplay\fontsize{12}{12}\selectfont\color{popcream}\thepopsec};%
  \hspace{8pt}{\popdisplay\fontsize{15}{17}\selectfont\color{popink}\MakeUppercase{#1}}\par
  \vspace{4pt}\noindent{\color{poporange}\rule{36pt}{3.6pt}}
  \end{minipage}\par\nopagebreak\vspace{8pt}
}
\newcounter{popsub}
\newcommand{\courssub}[1]{%
  \stepcounter{popsub}%
  \par\vspace{9pt}\noindent{\popxbold\fontsize{11.5}{13}\selectfont\color{popink}{\color{poporange}\thepopsec.\thepopsub}\hspace{6pt}#1}\par\nopagebreak\vspace{4pt}
}

% ============================================================
% Boîtes pédagogiques — même recette : fond blanc, bordure épaisse colorée,
% ombre dure encre, badge plein. Titre optionnel [#1].
% ============================================================
\tcbset{popbase/.style={breakable,enhanced,colback=white,boxrule=2.3pt,arc=9pt,
  shadow={3pt}{-3pt}{0pt}{popink},left=14pt,right=14pt,top=11pt,bottom=11pt,before skip=11pt,after skip=15pt,
  fontupper=\popsemi\fontsize{10.6}{14.5}\selectfont\color{popink},
  fontlower=\popsemi\fontsize{10.6}{14.5}\selectfont\color{popink}}}
% Coche dessinée (le glyphe ✓ n'existe pas dans les polices du document)
\newcommand{\popcheck}[1]{\tikz[baseline=-0.5ex]\draw[#1,line width=1.6pt,line cap=round,line join=round](-0.13,0.0)--(-0.04,-0.09)--(0.13,0.1);}
\providecommand{\coche}{\popcheck{popgreen}}
\providecommand{\resultat}[1]{\par\smallskip\noindent\fcolorbox{popgreen}{popgreenL}{\parbox{\dimexpr\linewidth-2\fboxsep-2\fboxrule\relax}{\textbf{Result.} #1}}\par\smallskip}
\newcommand{\popboxhead}[3]{\poppill{#1}{#2}{white}\ifx\relax#3\relax\else\hspace{6pt}{\popxbold\fontsize{10.6}{12}\selectfont\color{popink}#3}\fi\par\vspace{7pt}}

\newtcolorbox{aretenir}[1][]{popbase,colframe=popgreen,before upper={\popboxhead{Key points}{popgreen}{#1}}}
\newtcolorbox{exemplebox}[1][]{popbase,colframe=poporange,before upper={\popboxhead{Exemple}{poporange}{#1}}}
\newtcolorbox{definition}[1][]{popbase,colframe=popblue,before upper={\popboxhead{Definition}{popblue}{#1}}}
\newtcolorbox{propriete}[1][]{popbase,colframe=poppurple,before upper={\popboxhead{Property}{poppurple}{#1}}}
\newtcolorbox{methode}[1][]{popbase,colframe=popgold,before upper={\popboxhead{Method}{popgold}{#1}}}
\newtcolorbox{attention}[1][]{popbase,colframe=poppink,before upper={\popboxhead{Warning}{poppink}{#1}}}
\newtcolorbox{experience}[1][]{popbase,colframe=popblue,colback=popblueL,before upper={\popboxhead{Experiment}{popblue}{#1}}}
% « Le savais-tu ? » : carte violette pleine (contexte, culture, Cameroun)
\newtcolorbox{savaistu}[1][]{popbase,colback=poppurple,colframe=popink,
  fontupper=\popsemi\fontsize{10.4}{14.2}\selectfont\color{white},
  before upper={\poppill{Did you know?}{white}{poppurple}\ifx\relax#1\relax\else\hspace{6pt}{\popxbold\color{white}#1}\fi\par\vspace{7pt}}}
% Exercice résolu : énoncé en haut, \tcblower puis la solution sur fond crème
\newcounter{popexo}\newcounter{popexr}
\newtcolorbox{exoresolu}[1][]{popbase,colframe=poporange,colbacklower=popcream,
  segmentation style={popink,line width=1.2pt,dashed},
  before upper={\stepcounter{popexr}\popboxhead{Worked example \thepopexr}{poporange}{#1}},
  before lower={\poppill{Solution}{popink}{popcream}\par\vspace{6pt}}}
% Exercice (non résolu en place) — la correction est dans la partie Corrigés
\newtcolorbox{exercice}[1][]{popbase,colframe=popink,
  before upper={\stepcounter{popexo}\popboxhead{Exercise \thepopexo}{popink}{#1}}}
% Corrigé
\newtcolorbox{corrige}[1][]{popbase,colframe=popgreen,colback=popgreenL,
  before upper={\popboxhead{Solution}{popgreen}{#1}}}
% Objectifs / prérequis / bilan
\newtcolorbox{objectifs}{popbase,colframe=popink,colback=poporangeL,
  before upper={\popboxhead{Chapter objectives}{popink}{}}}
\newtcolorbox{prerequis}{popbase,colframe=popink,colback=popblueL,
  before upper={\popboxhead{Prerequisites}{popblue}{}}}
\newtcolorbox{bilan}{popbase,colframe=popink,colback=white,boxrule=2.8pt,
  before upper={\popboxhead{Fiche bilan}{poporange}{L'essentiel à connaître pour l'examen}}}
% Figure encadrée avec légende
\newtcolorbox{popfigure}[1][]{enhanced,breakable=false,colback=white,colframe=popline,boxrule=1.4pt,arc=8pt,
  left=10pt,right=10pt,top=10pt,bottom=8pt,before skip=10pt,after skip=12pt,halign=center,
  lower separated=false,halign lower=center,
  fontlower=\popsemi\fontsize{9}{11.5}\selectfont\color{popmuted},
  #1}
\newcommand{\legende}[1]{\par\vspace{6pt}{\popsemi\fontsize{9}{11.5}\selectfont\color{popmuted}#1}}

% ============================================================
% Signature OnBuch+
% ============================================================
\newcommand{\poplogoinline}[1]{{\poplogo\fontsize{#1}{#1}\selectfont\color{popink}OnBuch\color{poporange}+}}
\newcommand{\signatureonbuch}{%
  \begin{tcolorbox}[enhanced,colback=popink,colframe=popink,boxrule=2.2pt,arc=10pt,
      drop shadow=poporange,left=14pt,right=14pt,top=10pt,bottom=10pt,before skip=0pt,after skip=0pt]
    \tikz[baseline=-0.7ex]\node[circle,fill=popgreen,draw=popcream,line width=1.3pt,minimum size=22pt,inner sep=0]{\popcheck{white}};%
    \hspace{9pt}\begin{minipage}[c]{0.62\linewidth}
      {\popxbold\fontsize{12}{13}\selectfont\color{popcream}Signed\ \textbullet\ {\poplogo OnBuch\color{poporange}+}}\par\vspace{2pt}
      {\raggedright\popsemi\fontsize{8}{10.5}\selectfont\color{popline}\MakeUppercase{OnBuch+ teaching team}\\\MakeUppercase{Follows the official MINESEC syllabus}\par}
    \end{minipage}\hfill
    {\poplogo\fontsize{22}{22}\selectfont\color{popcream}OnBuch\color{poporange}+}
  \end{tcolorbox}%
}

% ============================================================
% Page de garde
% #1 titre · #2 matière · #3 niveau · #4 module · #5 n° leçon · #6 durée ·
% #7 description / objectifs (texte ou itemize)
% ============================================================
\newcommand{\pagegarde}[7]{%
  \thispagestyle{empty}
  \newgeometry{a4paper,margin=1.7cm,top=1.6cm,bottom=1.4cm}%
  \begin{tikzpicture}[remember picture,overlay]
    \fill[poporange!18] ([xshift=-3.2cm,yshift=-3.6cm]current page.north east) circle (4.6cm);
    \fill[poppurple!14] ([xshift=2.4cm,yshift=7.6cm]current page.south west) circle (3.8cm);
  \end{tikzpicture}%
  {\poplogo\fontsize{30}{30}\selectfont\color{popink}OnBuch\color{poporange}+}\hfill
  \raisebox{0.6ex}{\poppill{Signed course \textbullet\ Premium}{popink}{popcream}}\par
  \vspace{1pt}\popsquiggle{poppurple}\par
  \vspace{8pt}
  \begin{tcolorbox}[enhanced,colback=poporange,colframe=popink,boxrule=2.8pt,arc=13pt,
      drop shadow=popink,left=18pt,right=18pt,top=12pt,bottom=13pt,after skip=10pt]
    \poppill{#2 \textbullet\ #3}{popink}{popcream}\hspace{5pt}\poppill{#4}{white}{popink}\par\vspace{8pt}
    {\popdisplay\fontsize{14}{14}\selectfont\color{popink}CHAPTER #5}\par\vspace{4pt}
    {\raggedright\hyphenpenalty=10000\exhyphenpenalty=10000\popdisplay\fontsize{25}{27}\selectfont\color{popcream}\MakeUppercase{#1}\par}
    \vspace{8pt}
    {\popxbold\fontsize{9.5}{9.5}\selectfont\color{popink}\MakeUppercase{Suggested time: #6}}
  \end{tcolorbox}
  \begin{tcolorbox}[enhanced,colback=white,colframe=popink,boxrule=2.2pt,arc=10pt,
      drop shadow=popink,left=16pt,right=16pt,top=10pt,bottom=10pt,after skip=8pt]
    \poppill{In this chapter}{poppurple}{white}\par\vspace{5pt}
    \popsemi\fontsize{9.3}{12.6}\selectfont\color{popink}#7
  \end{tcolorbox}
  \noindent\begin{minipage}[t]{0.487\linewidth}
    \begin{tcolorbox}[enhanced,colback=popgreen,colframe=popink,boxrule=2.2pt,arc=10pt,
        drop shadow=popink,left=11pt,right=11pt,top=10pt,bottom=10pt,height=3.1cm,valign=center]
      \tcbox[colback=white,colframe=popink,boxrule=1.3pt,arc=5pt,boxsep=0pt,nobeforeafter,
        left=3pt,right=3pt,top=3pt,bottom=3pt,valign=center]{\qrcode[height=1.9cm]{https://play.google.com/store/apps/details?id=com.luvvix.onbuch&hl=en}}%
      \hspace{8pt}\begin{minipage}[c]{0.5\linewidth}
        {\popdisplay\fontsize{11}{12.5}\selectfont\color{white}\MakeUppercase{Download OnBuch}}\par\vspace{4pt}
        {\raggedright\popsemi\fontsize{8.4}{11}\selectfont\color{white}Free \textbullet\ Google Play\\AI tutor, past papers, results\par}
      \end{minipage}
    \end{tcolorbox}
  \end{minipage}\hfill
  \begin{minipage}[t]{0.487\linewidth}
    \begin{tcolorbox}[enhanced,colback=poppurple,colframe=popink,boxrule=2.2pt,arc=10pt,
        drop shadow=popink,left=11pt,right=11pt,top=10pt,bottom=10pt,height=3.1cm,valign=center]
      \tikz[baseline=-0.7ex]\node[circle,fill=white,draw=popink,line width=1.3pt,minimum size=1.6cm,inner sep=0]{\popdisplay\fontsize{24}{24}\selectfont\color{poppurple}+};%
      \hspace{8pt}\begin{minipage}[c]{0.6\linewidth}
        {\popdisplay\fontsize{11}{12.5}\selectfont\color{white}\MakeUppercase{Go OnBuch+}}\par\vspace{4pt}
        {\raggedright\popsemi\fontsize{8.4}{11}\selectfont\color{white}All signed courses, unlimited AI tutor, PDF sheets — OnBuch+ tab in the app\par}
      \end{minipage}
    \end{tcolorbox}
  \end{minipage}
  \vspace{8pt}
  \popcontenu
  \vfill
  \signatureonbuch
  \restoregeometry
  \newpage
}

% Rangée « Ce cours contient » (page de garde)
\newcommand{\poptile}[3]{% #1 couleur #2 gros texte #3 légende
  \begin{tcolorbox}[enhanced,colback=#1,colframe=popink,boxrule=2pt,arc=9pt,shadow={2.5pt}{-2.5pt}{0pt}{popink},
      left=6pt,right=6pt,top=9pt,bottom=9pt,halign=center,valign=center,height=1.75cm,width=\linewidth,nobeforeafter]
    {\popdisplay\fontsize{12.5}{13}\selectfont\color{white}\MakeUppercase{#2}}\par\vspace{3pt}
    {\centering\popsemi\fontsize{7.8}{9.5}\selectfont\color{white}#3\par}
  \end{tcolorbox}}
\newcommand{\popcontenu}{%
  \par\noindent{\popxboldls\fontsize{7.5}{7.5}\selectfont\color{popmuted}THIS SIGNED COURSE CONTAINS}\par\vspace{7pt}
  \noindent\begin{minipage}[t]{0.235\linewidth}\poptile{popblue}{Course}{complete, illustrated, step by step}\end{minipage}\hfill
  \begin{minipage}[t]{0.235\linewidth}\poptile{poporange}{Methods}{and worked examples}\end{minipage}\hfill
  \begin{minipage}[t]{0.235\linewidth}\poptile{poppink}{Exercises}{exam style + solutions}\end{minipage}\hfill
  \begin{minipage}[t]{0.235\linewidth}\poptile{popgreen}{Summary sheet}{the essentials to remember}\end{minipage}\par}

% Sommaire
\newtcolorbox{sommaire}{enhanced,colback=white,colframe=popink,boxrule=2.2pt,arc=10pt,
  drop shadow=popink,left=18pt,right=18pt,top=13pt,bottom=13pt,before skip=6pt,after skip=20pt,
  before upper={\poppill{Contents}{poppurple}{white}\par\vspace{9pt}\popsemi\fontsize{10.6}{14.5}\selectfont\color{popink}}}

% Signature de fin de cours — poussée tout en bas de la dernière page
\newcommand{\findecours}{%
  \par\vfill\nopagebreak
  \begin{center}\popsquiggle{poporange}\end{center}\vspace{4pt}
  \signatureonbuch
}
\AtBeginDocument{\def\llbracket{\mathopen{[\mkern-3.5mu[}}\def\rrbracket{\mathclose{]\mkern-3.5mu]}}} % intervalles d'entiers (glyphe ⟦ absent de la police)
% Glyphes absents de Fira Math : redéfinis à partir de glyphes disponibles
\AtBeginDocument{%
  \def\setminus{\mathbin{\backslash}}\let\smallsetminus\setminus
  \def\vdots{\mathord{\vbox{\baselineskip3.2pt\lineskiplimit0pt\kern4pt\hbox{.}\hbox{.}\hbox{.}}}}%
  \def\ddots{\mathinner{\mkern1mu\raise6.5pt\hbox{.}\mkern2mu\raise3.6pt\hbox{.}\mkern2mu\raise0.7pt\hbox{.}\mkern1mu}}%
  \def\triangle{\mathord{\scalebox{0.9}{$\symup{\Delta}$}}}%
  \def\nabla{\mathord{\raisebox{\depth}{\scalebox{1}[-1]{$\symup{\Delta}$}}}}%
  \def\checkmark{\popcheck{popdark}}%
  \def\star{\mathbin{\ast}}\def\bigstar{\mathord{\ast}}%
  \def\diamond{\mathbin{\scalebox{0.6}{$\Diamond$}}}%
  \def\bigcup{\mathop{\mathchoice{\raisebox{-0.3em}{\scalebox{1.7}{$\cup$}}}{\scalebox{1.25}{$\cup$}}{\cup}{\cup}}\displaylimits}%
  \def\bigcap{\mathop{\mathchoice{\raisebox{-0.3em}{\scalebox{1.7}{$\cap$}}}{\scalebox{1.25}{$\cap$}}{\cap}{\cap}}\displaylimits}%
  \def\lesseqqgtr{\mathrel{\lessgtr}}\def\bigoplus{\mathop{\oplus}\displaylimits}%
}
\providecommand{\ssi}{if and only if} % abréviation fréquente
\AtBeginDocument{\providecommand{\wideparen}{\overparen}} % arc de cercle

% Fonctions usuelles que les modèles utilisent sans que amsmath les fournisse
\providecommand{\arsinh}{\operatorname{arsinh}}\providecommand{\arcosh}{\operatorname{arcosh}}
\providecommand{\artanh}{\operatorname{artanh}}\providecommand{\arsech}{\operatorname{arsech}}
\providecommand{\arcsch}{\operatorname{arcsch}}\providecommand{\arcoth}{\operatorname{arcoth}}
\providecommand{\arcsinh}{\operatorname{arsinh}}\providecommand{\arccosh}{\operatorname{arcosh}}
\providecommand{\arctanh}{\operatorname{artanh}}\providecommand{\sech}{\operatorname{sech}}
\providecommand{\csch}{\operatorname{csch}}\providecommand{\arcsec}{\operatorname{arcsec}}
\providecommand{\arccsc}{\operatorname{arccsc}}\providecommand{\arccot}{\operatorname{arccot}}
\providecommand{\sgn}{\operatorname{sgn}}\providecommand{\lcm}{\operatorname{lcm}}
\providecommand{\Var}{\operatorname{Var}}\providecommand{\Cov}{\operatorname{Cov}}

%% FIGFIX-BEGIN v4 — correctifs globaux des figures (docs/onbuch-generation/tools/figfix)
\usepackage{adjustbox}\usepackage{etoolbox}
% toute figure TikZ/pgfplots plus large que la ligne est réduite à la largeur disponible
\BeforeBeginEnvironment{tikzpicture}{\begin{adjustbox}{max width=\linewidth}}
\AfterEndEnvironment{tikzpicture}{\end{adjustbox}}
% légendes pgfplots lisibles par-dessus les courbes
\pgfplotsset{every axis/.append style={legend style={fill=white,fill opacity=0.92,text opacity=1,draw=black!15,inner sep=3pt}}}
% étiquettes d'arêtes (syntaxe edge["..."]) avec fond blanc
\tikzset{every edge quotes/.append style={fill=white,inner sep=1.2pt}}
%% FIGFIX-END
\begin{document}

\pagegarde{Vectors}{Pure Mathematics with Statistics}{Upper Sixth}{Upper Sixth — Pure mathematics with statistics}{3}{11 periods}{This chapter extends the vector work of Lower Sixth to planes in three-dimensional space: their vector, parametric and Cartesian equations, distances, angles and intersections. It provides the complete geometric toolkit needed for the 3-D geometry questions of the GCE Advanced Level examination.
\vspace{4pt}
\begin{multicols}{2}\small
\begin{itemize}
\item By the end of this chapter I can write the vector equation r·n = d and the Cartesian equation ax + by + cz = d of a plane from a point and a normal vector.
\item By the end of this chapter I can find the Cartesian equation of a plane through three non-collinear points.
\item By the end of this chapter I can convert between the parametric (vector) form and the Cartesian form of the equation of a plane.
\item By the end of this chapter I can calculate the perpendicular distance from a point to a plane and the distance between two parallel planes.
\item By the end of this chapter I can find the acute angle between two planes and the angle between a line and a plane.
\item By the end of this chapter I can determine the point of intersection of a line and a plane, including the cases where the line is parallel to or lies in the plane.
\end{itemize}\end{multicols}}

\begin{sommaire}
\begin{multicols}{2}
\begin{enumerate}
\item Equations of a Plane: Vector and Cartesian Forms
\item Plane Through Three Points and Parametric Form
\item Perpendicular Distance from a Point to a Plane
\item Angles Involving Planes and Lines
\item Intersection of a Line and a Plane; Intersection of Two Planes
\item Families of Planes Through an Intersection and Synthesis
\item Integration activity
\item Methods and worked examples
\item Exercises
\item Solutions to the exercises
\item Summary sheet
\end{enumerate}
\end{multicols}
\end{sommaire}

\begin{objectifs}
\begin{itemize}
\item By the end of this chapter I can write the vector equation r·n = d and the Cartesian equation ax + by + cz = d of a plane from a point and a normal vector.
\item By the end of this chapter I can find the Cartesian equation of a plane through three non-collinear points.
\item By the end of this chapter I can convert between the parametric (vector) form and the Cartesian form of the equation of a plane.
\item By the end of this chapter I can calculate the perpendicular distance from a point to a plane and the distance between two parallel planes.
\item By the end of this chapter I can find the acute angle between two planes and the angle between a line and a plane.
\item By the end of this chapter I can determine the point of intersection of a line and a plane, including the cases where the line is parallel to or lies in the plane.
\item By the end of this chapter I can find the line of intersection of two non-parallel planes in vector form.
\item By the end of this chapter I can write the equation of the family of planes through the line of intersection of two planes and select the member satisfying an extra condition.
\item By the end of this chapter I can solve GCE-style problems combining several of these skills.
\end{itemize}
\end{objectifs}

\begin{prerequis}
\begin{itemize}
\item Vector algebra from Lower Sixth: position vectors, addition, scalar multiplication, magnitude of a vector
\item The scalar (dot) product and its use to find angles and test perpendicularity
\item The vector and Cartesian equations of a line in 3-D, including parametric form
\item Solving simultaneous linear equations in two and three unknowns
\item Basic coordinate geometry in 3-D: coordinates of points, distance between two points
\end{itemize}
\end{prerequis}

\newpage

\coursec{Equations of a Plane: Vector and Cartesian Forms}

When engineers designed the Douala International Airport runway, they had to ensure that the landing surface was perfectly planar and that the approach path of each aircraft met the runway at a safe, precisely calculated angle. A mason in Bamenda checking whether a wall is perpendicular to the floor, or a carpenter in Yaound\'e cutting a roof panel to fit between two sloping faces, is unknowingly solving problems about planes, lines and angles in three-dimensional space. How can we describe a flat surface with an equation, measure how far a point is from it, and find where a line pierces it? This chapter gives you the vector and Cartesian tools to answer exactly these questions --- the same tools that appear every year in the GCE Advanced Level Paper 2.

In this first section we answer the most fundamental question: \emph{how do we write down the equation of a plane?} You already know from Lower Sixth that a line in space is fixed by one point and one direction. We shall see that a plane is fixed by one point and one \emph{perpendicular} direction, called the \cle{normal vector}.

\courssub{What fixes a plane? Intuition and the normal vector}

Take a flat sheet of paper and hold it in the air. It can slide, tilt and rotate: it is not yet ``fixed''. Now pin one point $A$ of the sheet to a fixed position. The sheet can still spin around $A$ in infinitely many ways. Finally, decide the \emph{orientation} of the sheet by specifying the direction perpendicular to it --- imagine a nail driven through $A$ at right angles to the paper. Once the nail's direction is chosen, the sheet cannot move at all. This simple experiment reveals the key fact:

\begin{center}
\emph{A plane is completely determined by one of its points and the direction perpendicular to it.}
\end{center}

\begin{definition}[Normal vector of a plane]
A \cle{normal vector} to a plane $\Pi$ is any non-zero vector $\vec{n}$ which is perpendicular to \emph{every} line lying in $\Pi$. Equivalently, $\vec{n}$ is perpendicular to every vector joining two points of $\Pi$.
\end{definition}

\begin{attention}[The normal is not unique]
If $\vec{n}$ is normal to a plane, then so is $k\vec{n}$ for any non-zero scalar $k$: the vectors $\begin{pmatrix}1\\2\\-1\end{pmatrix}$ and $\begin{pmatrix}-2\\-4\\2\end{pmatrix}$ are normal to the \emph{same} plane. What matters is the \emph{direction} of the normal, not its length. In particular a plane has infinitely many normal vectors, but they are all parallel to each other.
\end{attention}

\begin{popfigure}
\begin{tikzpicture}[scale=1.05, line cap=round, line join=round]
  % plane as parallelogram
  \coordinate (O) at (0,0);
  \coordinate (U) at (5.2,0.6);
  \coordinate (V) at (1.6,2.4);
  \coordinate (W) at ($(U)+(V)$);
  \fill[popblueL, opacity=0.55] (O) -- (U) -- (W) -- (V) -- cycle;
  \draw[popblue, thick] (O) -- (U) -- (W) -- (V) -- cycle;
  \node[popblue] at (4.6,2.6) {$\Pi$};
  % origin below the plane
  \coordinate (Org) at (-1.6,-1.9);
  \fill (Org) circle (1.6pt) node[below left] {$O$};
  % point A on the plane
  \coordinate (A) at (2.0,1.0);
  \fill[popdark] (A) circle (1.8pt) node[below right, popdark] {$A$};
  % general point P on the plane
  \coordinate (P) at (4.1,1.7);
  \fill[popdark] (P) circle (1.8pt) node[right, popdark] {$P$};
  % position vectors
  \draw[-latex, poporange, thick] (Org) -- (A) node[midway, below, sloped] {$\vec{a}$};
  \draw[-latex, poporange, thick] (Org) -- (P) node[midway, above, sloped] {$\vec{r}$};
  % AP in the plane
  \draw[-latex, popgreen, thick] (A) -- (P) node[midway, below, popgreen] {$\overrightarrow{AP}$};
  % normal vector at A
  \draw[-latex, poppink, very thick] (A) -- ($(A)+(0.55,1.9)$) node[above, poppink] {$\vec{n}$};
  % right angle mark
  \draw[poppink] ($(A)+(0.16,0.55)$) -- ($(A)+(0.42,0.56)$) -- ($(A)+(0.44,0.24)$);
\end{tikzpicture}
\legende{A plane $\Pi$ through the point $A$ with position vector $\vec{a}$. The vector $\vec{n}$ is normal to $\Pi$; for any point $P(\vec{r})$ of the plane, $\overrightarrow{AP}$ lies in $\Pi$ and is perpendicular to $\vec{n}$.}
\end{popfigure}

\courssub{Derivation of the vector equation $\vec{r}\cdot\vec{n}=\vec{a}\cdot\vec{n}$}

Let $\Pi$ be the plane through the fixed point $A$, with position vector $\vec{a}=\overrightarrow{OA}$, and with normal vector $\vec{n}$. Take any point $P$ of space, with position vector $\vec{r}=\overrightarrow{OP}$. We ask: \emph{under what condition on $\vec{r}$ does $P$ lie on $\Pi$?}

By the triangle law of vector addition,
\[ \overrightarrow{AP} = \overrightarrow{AO} + \overrightarrow{OP} = \vec{r} - \vec{a}. \]
Now $P$ lies on $\Pi$ if and only if the segment $AP$ lies in the plane, and by the definition of the normal vector this happens if and only if $\overrightarrow{AP}$ is perpendicular to $\vec{n}$. Two vectors are perpendicular if and only if their scalar product is zero, so
\[ P \in \Pi \iff \vec{n}\cdot\overrightarrow{AP} = 0 \iff \vec{n}\cdot(\vec{r}-\vec{a}) = 0 \iff \vec{n}\cdot\vec{r} = \vec{n}\cdot\vec{a}. \]
Since $\vec{a}$ and $\vec{n}$ are fixed, the quantity $\vec{a}\cdot\vec{n}$ is a fixed \emph{constant}, which we call $d$. We have proved the central result of this section.

\begin{propriete}[Vector equation of a plane --- proof examinable]
The plane through the point $A$ with position vector $\vec{a}$, with normal vector $\vec{n}$, has vector equation
\[ \vec{r}\cdot\vec{n} = \vec{a}\cdot\vec{n} = d, \]
where $\vec{r}$ is the position vector of a general point of the plane and $d$ is a constant. The proof (using $\vec{n}\cdot(\vec{r}-\vec{a})=0$) is examinable and is frequently demanded with the instruction ``hence'' or ``show that''.
\end{propriete}

\begin{popfigure}
\begin{tikzpicture}[scale=1.0, line cap=round]
  \coordinate (O) at (0,0);
  \coordinate (U) at (4.6,0.5);
  \coordinate (V) at (1.4,2.1);
  \coordinate (W) at ($(U)+(V)$);
  \fill[popgreenL, opacity=0.6] (O) -- (U) -- (W) -- (V) -- cycle;
  \draw[popgreen, thick] (O) -- (U) -- (W) -- (V) -- cycle;
  \node[popgreen] at (4.3,2.3) {$\Pi$};
  \coordinate (A) at (1.7,0.9);
  \coordinate (P) at (3.6,1.6);
  \fill[popdark] (A) circle (1.7pt) node[below left, popdark] {$A$};
  \fill[popdark] (P) circle (1.7pt) node[right, popdark] {$P$};
  \draw[-latex, popgreen, very thick] (A) -- (P) node[midway, below, popgreen] {$\overrightarrow{AP}=\vec{r}-\vec{a}$};
  \draw[-latex, poppink, very thick] (A) -- ($(A)+(0.5,1.7)$) node[above, poppink] {$\vec{n}$};
  \draw[poppink] ($(A)+(0.15,0.5)$) -- ($(A)+(0.39,0.51)$) -- ($(A)+(0.41,0.2)$);
  \node[popdark, align=center] at (2.6,-0.55) {$\vec{n}\cdot\overrightarrow{AP}=0$};
\end{tikzpicture}
\legende{The key geometric fact: for every point $P$ of the plane, $\overrightarrow{AP}$ lies in the plane, so $\vec{n}\cdot\overrightarrow{AP}=0$.}
\end{popfigure}

\courssub{From the vector form to the Cartesian form $ax+by+cz=d$}

Write the vectors in components with respect to the usual orthonormal basis $(\vec{i},\vec{j},\vec{k})$:
\[ \vec{r} = \begin{pmatrix}x\\y\\z\end{pmatrix}, \qquad \vec{n} = \begin{pmatrix}a\\b\\c\end{pmatrix}, \qquad \vec{a} = \begin{pmatrix}x_0\\y_0\\z_0\end{pmatrix}. \]
Computing the scalar products in $\vec{r}\cdot\vec{n} = \vec{a}\cdot\vec{n}$ gives
\[ ax + by + cz = ax_0 + by_0 + cz_0. \]
The right-hand side is a constant; calling it $d$ we obtain the \cle{Cartesian equation} of the plane:
\[ ax + by + cz = d. \]

\begin{aretenir}[Reading the normal from the Cartesian equation]
In the Cartesian equation $ax+by+cz=d$, the coefficients of $x$, $y$ and $z$ are exactly the components of a normal vector:
\[ \vec{n} = a\vec{i} + b\vec{j} + c\vec{k} = \begin{pmatrix}a\\b\\c\end{pmatrix}. \]
For example, the plane $2x - 3y + 5z = 7$ has normal vector $\begin{pmatrix}2\\-3\\5\end{pmatrix}$. The constant $d$ on the right-hand side carries the information about \emph{which} point the plane passes through.
\end{aretenir}

\begin{attention}[Coefficients give the normal, NOT a point]
A classic GCE error: faced with $2x - 3y + 5z = 7$, candidates claim that $(2,-3,5)$ is a point of the plane. It is not --- it is the \emph{normal vector}. Check: substituting $x=2$, $y=-3$, $z=5$ gives $4+9+25 = 38 \neq 7$, so the point $(2,-3,5)$ does not even lie on the plane. The coefficients $(a,b,c)$ always describe a \emph{direction perpendicular} to the plane, never a point on it.
\end{attention}

\courssub{Finding the equation of a plane: the standard method}

\begin{methode}[Equation of a plane from a point and a normal]
To find the equation of a plane:
\begin{enumerate}
\item Identify a \textbf{point} $A(x_0,y_0,z_0)$ on the plane.
\item Identify a \textbf{normal vector} $\vec{n} = \begin{pmatrix}a\\b\\c\end{pmatrix}$ (given directly, or obtained from a line perpendicular to the plane, or from the normal of a parallel plane).
\item Compute $d = \vec{a}\cdot\vec{n} = ax_0 + by_0 + cz_0$.
\item Write the equation $\vec{r}\cdot\vec{n} = d$ (vector form) or $ax+by+cz=d$ (Cartesian form).
\item \textbf{Check}: substitute the coordinates of $A$ into your equation --- it must be satisfied.
\end{enumerate}
\end{methode}

\begin{exoresolu}[Plane through a given point with a given normal]
Find the vector and Cartesian equations of the plane $\Pi$ passing through the point $A(1,-2,3)$ and perpendicular to the vector $\vec{n} = 2\vec{i} + \vec{j} - 4\vec{k}$.
\tcblower
\textbf{Solution.} We have a point and a normal, so we apply the method. With $\vec{a} = \begin{pmatrix}1\\-2\\3\end{pmatrix}$ and $\vec{n} = \begin{pmatrix}2\\1\\-4\end{pmatrix}$:
\[ d = \vec{a}\cdot\vec{n} = (1)(2) + (-2)(1) + (3)(-4) = 2 - 2 - 12 = -12. \]
Vector equation:
\[ \vec{r}\cdot\begin{pmatrix}2\\1\\-4\end{pmatrix} = -12. \]
Cartesian equation: writing $\vec{r} = x\vec{i}+y\vec{j}+z\vec{k}$,
\[ 2x + y - 4z = -12. \]
\textbf{Check}: at $A(1,-2,3)$, $2(1)+(-2)-4(3) = 2-2-12 = -12$. \checkmark\ The point lies on the plane, as required.
\end{exoresolu}

\begin{exoresolu}[Plane perpendicular to a given line]
A line $L$ has vector equation $\vec{r} = \begin{pmatrix}3\\0\\-1\end{pmatrix} + \lambda\begin{pmatrix}1\\-2\\2\end{pmatrix}$. Find the Cartesian equation of the plane which is perpendicular to $L$ and passes through the point $B(2,1,4)$.
\tcblower
\textbf{Solution.} The plane is perpendicular to $L$, so the \emph{direction vector of $L$ is a normal vector of the plane}:
\[ \vec{n} = \begin{pmatrix}1\\-2\\2\end{pmatrix}. \]
Then
\[ d = \vec{n}\cdot\overrightarrow{OB} = (1)(2) + (-2)(1) + (2)(4) = 2 - 2 + 8 = 8. \]
Hence the plane has equation
\[ x - 2y + 2z = 8. \]
\textbf{Check}: at $B(2,1,4)$, $2 - 2 + 8 = 8$. \checkmark
\end{exoresolu}

\courssub{Parallel planes and planes perpendicular to a line}

Two planes that never meet --- like the floor and the ceiling of a classroom --- have exactly the same orientation, which means their normals point in the same (or opposite) direction.

\begin{propriete}[Parallel planes --- exam tip]
Two planes are \cle{parallel} if and only if their normal vectors are parallel, i.e. one normal is a scalar multiple of the other. Thus $a_1x+b_1y+c_1z=d_1$ and $a_2x+b_2y+c_2z=d_2$ are parallel exactly when $(a_2,b_2,c_2) = k(a_1,b_1,c_1)$ for some non-zero $k$. Consequently, \emph{any plane parallel to $ax+by+cz=d$ has equation $ax+by+cz=d'$} for some new constant $d'$, found using a point of the required plane. In the examination, state this justification explicitly to earn the method mark.
\end{propriete}

\begin{exemplebox}[Plane parallel to a given plane]
Find the equation of the plane parallel to $\Pi: 3x - y + 2z = 5$ and passing through $C(0,2,-1)$.\par
\textbf{Solution.} A parallel plane has the same normal $\vec{n} = \begin{pmatrix}3\\-1\\2\end{pmatrix}$, so its equation is $3x - y + 2z = d'$. Substituting $C$:
\[ d' = 3(0) - (2) + 2(-1) = -4. \]
Hence the required plane is $3x - y + 2z = -4$.
\end{exemplebox}

\courssub{Checking whether a point lies on a plane}

A point $M(x_1,y_1,z_1)$ lies on the plane $ax+by+cz=d$ if and only if its coordinates \emph{satisfy} the equation, that is $ax_1+by_1+cz_1 = d$. In vector language, $M$ lies on the plane $\vec{r}\cdot\vec{n}=d$ if and only if $\overrightarrow{OM}\cdot\vec{n} = d$.

\begin{exemplebox}[Testing points on a plane]
Does the point $M(2,0,3)$ lie on the plane $\Pi: 2x + y - 4z = -12$? What about $N(1,1,1)$?\par
\textbf{Solution.} For $M$: $2(2) + 0 - 4(3) = 4 - 12 = -8 \neq -12$, so $M \notin \Pi$.\par
For $N$: $2(1) + 1 - 4(1) = -1 \neq -12$, so $N \notin \Pi$ either. However $A(1,-2,3)$ gives $-12$, so $A \in \Pi$ (as we built it).
\end{exemplebox}

\begin{attention}[Common mistakes to avoid]
\begin{itemize}
\item Forgetting that $\vec{n}$ must be \textbf{non-zero}: $0x+0y+0z=d$ does not define a plane.
\item Mixing up signs when reading the normal: the normal of $x - 2y + 2z = 8$ is $(1,-2,2)$, not $(1,2,2)$.
\item Computing $d$ with the wrong point, or with the normal vector instead of the position vector of the given point.
\item Giving $ax+by+cz=0$ systematically: the right-hand side is $0$ \emph{only} when the plane passes through the origin.
\end{itemize}
\end{attention}

\begin{savaistu}[Planes in the Cameroonian landscape]
The idea of describing a flat surface by a point and a perpendicular direction is used daily by surveyors mapping the slopes of Mount Cameroon and by civil engineers levelling the terraces of the Bamenda highlands. In geodesy, the ``local horizontal plane'' at a site is defined precisely by the vertical (the normal direction given by a plumb line) at one reference point --- exactly our data $\big(A,\vec{n}\big)$.
\end{savaistu}

\begin{exercice}[Practice]
\begin{enumerate}
\item Find the vector and Cartesian equations of the plane through $A(2,-1,5)$ with normal vector $\vec{n} = \begin{pmatrix}3\\2\\-1\end{pmatrix}$.
\item The line $\vec{r} = \begin{pmatrix}1\\2\\0\end{pmatrix} + \mu\begin{pmatrix}2\\-1\\3\end{pmatrix}$ is perpendicular to a plane $\Pi$ which contains the origin. Find the Cartesian equation of $\Pi$.
\item Find the equation of the plane parallel to $4x + y - 3z = 7$ and passing through $P(1,1,1)$.
\item Determine which of the points $Q(3,-1,2)$ and $R(0,4,-1)$ lie on the plane $x + 2y - z = -1$.
\end{enumerate}
\end{exercice}

\begin{corrige}[Exercise 1]
\textbf{1.} $d = 3(2)+2(-1)-1(5) = 6-2-5 = -1$. Vector form: $\vec{r}\cdot\begin{pmatrix}3\\2\\-1\end{pmatrix} = -1$; Cartesian form: $3x+2y-z=-1$.\par
\textbf{2.} Normal $\vec{n} = (2,-1,3)$; plane through the origin so $d=0$: $2x - y + 3z = 0$.\par
\textbf{3.} Same normal: $4x+y-3z = d'$ with $d' = 4+1-3 = 2$. Equation: $4x+y-3z=2$.\par
\textbf{4.} For $Q$: $3 + 2(-1) - 2 = -1$, so $Q$ lies on the plane. For $R$: $0 + 8 + 1 = 9 \neq -1$, so $R$ does not.
\end{corrige}

\begin{aretenir}[Summary of Section 1]
\begin{itemize}
\item A plane is fixed by \textbf{one point} $A(\vec{a})$ and \textbf{one normal vector} $\vec{n}$.
\item Vector equation: $\vec{r}\cdot\vec{n} = \vec{a}\cdot\vec{n} = d$, derived from $\vec{n}\cdot(\vec{r}-\vec{a})=0$.
\item Cartesian equation: $ax+by+cz=d$, where $(a,b,c)$ are the components of the \textbf{normal}.
\item Parallel planes $\iff$ parallel normals; a parallel plane keeps the same left-hand side, only $d$ changes.
\item A point lies on the plane if and only if its coordinates satisfy the equation.
\end{itemize}
\end{aretenir}

\coursec{Plane Through Three Points and Parametric Form}

In the previous section we learnt that a plane is completely determined by one of its points $A$ and a \cle{normal vector} $\vec{n}$, and we wrote its equation in the vector form $(\vec{r}-\vec{a})\cdot\vec{n}=0$ and in the Cartesian form $ax+by+cz=d$. But in many problems --- and very often at the GCE --- you are not given a normal vector. Instead you are given \emph{three points} of the plane, or two directions lying in the plane. The purpose of this section is to show how the same machinery (dot product, normal vector, Cartesian equation) handles these situations, and to introduce a second, very flexible way of describing a plane: the \cle{parametric vector equation} $\vec{r}=\vec{a}+\lambda\vec{u}+\mu\vec{v}$.

\courssub{The Plane Determined by Three Non-Collinear Points}

\textbf{Intuition.} Place three fingertips of your hand on a table so that they are not in a straight line: the table top is the \emph{only} flat surface that can touch all three fingertips at once. But if the three fingertips are in a straight line, you can tilt a book around that line and still touch all three --- infinitely many planes pass through them. This is exactly the mathematical situation.

\begin{propriete}[Existence and uniqueness]
Through three \cle{non-collinear} points $A$, $B$, $C$ there passes exactly one plane. If the three points are collinear, infinitely many planes contain them.
\end{propriete}

\textbf{Strategy.} Since $A$, $B$, $C$ lie in the plane, the vectors $\vec{AB}$ and $\vec{AC}$ lie in the plane, and they are not parallel (precisely because $A,B,C$ are not collinear). A normal vector $\vec{n}$ must be perpendicular to \emph{every} vector of the plane, in particular to both $\vec{AB}$ and $\vec{AC}$. Using the dot product, this gives two equations:
\[ \vec{n}\cdot\vec{AB}=0 \qquad\text{and}\qquad \vec{n}\cdot\vec{AC}=0. \]
These two equations determine $\vec{n}$ \emph{up to a scalar multiple} (two equations, three unknowns), which is all we need: any non-zero multiple of a normal vector is still a normal vector.

\begin{popfigure}
\begin{center}
\begin{tikzpicture}[scale=1.1, >=stealth]
% plane
\fill[popblueL, opacity=0.5] (-2.6,-1.1) -- (1.4,-1.6) -- (3.4,1.0) -- (-0.6,1.6) -- cycle;
\draw[popline] (-2.6,-1.1) -- (1.4,-1.6) -- (3.4,1.0) -- (-0.6,1.6) -- cycle;
% triangle
\coordinate (A) at (-1.2,-0.5);
\coordinate (B) at (1.6,-0.8);
\coordinate (C) at (0.4,0.9);
\draw[popdark, thick] (A) -- (B) -- (C) -- cycle;
\fill (A) circle (1.6pt) node[below left] {$A$};
\fill (B) circle (1.6pt) node[below right] {$B$};
\fill (C) circle (1.6pt) node[above] {$C$};
% vectors AB, AC
\draw[->, poporange, very thick] (A) -- (B) node[midway, below] {$\vec{AB}$};
\draw[->, popgreen, very thick] (A) -- (C) node[midway, left] {$\vec{AC}$};
% normal
\draw[->, poppink, very thick] (A) -- ($(A)+(0.5,2.0)$) node[above] {$\vec{n}$};
\draw[popmuted] ($(A)+(0,0.35)$) -- ($(A)+(0.09,0.35)$) -- ($(A)+(0.09,0)$);
\end{tikzpicture}
\end{center}
\legende{Three non-collinear points $A$, $B$, $C$ determine a unique plane. The vectors $\vec{AB}$ and $\vec{AC}$ lie in the plane; the normal $\vec{n}$ is perpendicular to both.}
\end{popfigure}

\begin{methode}[Cartesian equation of the plane through three points]
\begin{enumerate}
\item \textbf{Check non-collinearity:} compute $\vec{AB}$ and $\vec{AC}$ and verify that one is not a scalar multiple of the other.
\item \textbf{Find a normal:} let $\vec{n}=\begin{pmatrix} p \\ q \\ r \end{pmatrix}$ and solve simultaneously $\vec{n}\cdot\vec{AB}=0$ and $\vec{n}\cdot\vec{AC}=0$. Express two of $p,q,r$ in terms of the third (the \emph{ratio method}), then choose a convenient value to get integer components.
\item \textbf{Find $d$:} write the equation as $px+qy+rz=d$ and substitute the coordinates of any one of the three points to compute $d$.
\item \textbf{Check:} verify that the other two points also satisfy the equation.
\end{enumerate}
\end{methode}

\begin{exoresolu}[Plane through three points]
Find the Cartesian equation of the plane passing through $A(1,0,2)$, $B(2,1,1)$ and $C(-1,2,1)$.
\tcblower
\textbf{Step 1: direction vectors in the plane.}
\[ \vec{AB}=\begin{pmatrix} 2-1 \\ 1-0 \\ 1-2 \end{pmatrix}=\begin{pmatrix} 1 \\ 1 \\ -1 \end{pmatrix}, \qquad
\vec{AC}=\begin{pmatrix} -1-1 \\ 2-0 \\ 1-2 \end{pmatrix}=\begin{pmatrix} -2 \\ 2 \\ -1 \end{pmatrix}. \]
These are not proportional (no single scalar $k$ gives $(1,1,-1)=k(-2,2,-1)$), so $A,B,C$ are not collinear and define a unique plane.

\textbf{Step 2: normal vector by the ratio method.} Let $\vec{n}=(p,q,r)$. Then
\begin{align*}
\vec{n}\cdot\vec{AB}&=0: & p+q-r &=0,\\
\vec{n}\cdot\vec{AC}&=0: & -2p+2q-r &=0.
\end{align*}
Subtracting the second equation from the first: $3p-q=0$, so $q=3p$. Substituting back: $r=p+q=4p$. Hence
\[ p:q:r = 1:3:4, \qquad \vec{n}=\begin{pmatrix} 1 \\ 3 \\ 4 \end{pmatrix}. \]

\textbf{Step 3: find $d$.} The equation has the form $x+3y+4z=d$. Using $A(1,0,2)$: $d=1+0+8=9$.

\textbf{Step 4: check.} $B$: $2+3+4=9$ \checkmark; $C$: $-1+6+4=9$ \checkmark.
\[ \boxed{x+3y+4z=9} \]
\end{exoresolu}

\courssub{Lesson 43: The Collinear Case --- Three Points on a Line}

The syllabus explicitly mentions the case of three collinear points. This is the \emph{trap case}: if $A$, $B$, $C$ are collinear, then $\vec{AB}$ and $\vec{AC}$ are parallel, the equations $\vec{n}\cdot\vec{AB}=0$ and $\vec{n}\cdot\vec{AC}=0$ reduce to a \emph{single} condition, and there are \textbf{infinitely many planes} through the three points (any plane containing the line $ABC$).

\begin{attention}[Three collinear points: the degenerate case]
\textbf{Always check first} that the three given points are not collinear before claiming a unique plane. For example $A(1,1,0)$, $B(2,2,1)$, $C(3,3,2)$ are collinear since $\vec{AB}=(1,1,1)$ and $\vec{AC}=(2,2,2)=2\vec{AB}$: no unique plane exists. In an examination, if you detect collinearity, state clearly: ``The points are collinear, therefore they do not define a unique plane; infinitely many planes contain the line through these points.''
\end{attention}

\begin{exemplebox}[Collinear points]
Show that $A(1,2,3)$, $B(0,1,1)$, $C(2,3,5)$ are collinear and explain the consequence.
$\vec{AB}=(-1,-1,-2)$, $\vec{AC}=(1,1,2)=-\vec{AB}$. The vectors are parallel, so the points are collinear. They lie on a single line; any plane containing that line contains all three points. There is no unique plane.
\end{exemplebox}

\courssub{The Parametric (Vector) Form of the Equation of a Plane}

\textbf{Intuition.} Stand at point $A$ of the plane. To reach \emph{any} other point $P$ of the plane, walk some distance along one direction $\vec{u}$ of the plane, then some distance along a second, different direction $\vec{v}$. The two ``amounts'' $\lambda$ and $\mu$ are your coordinates on the plane, exactly as $x$ and $y$ are coordinates on a map of Douala: two numbers suffice, because a plane is two-dimensional.

\begin{definition}[Parametric vector equation of a plane]
Let $A$ be a point of the plane with position vector $\vec{a}$, and let $\vec{u}$ and $\vec{v}$ be two \cle{non-parallel} vectors lying in the plane ($\vec{u}\nparallel\vec{v}$). The plane is the set of points $P$ with position vector
\[ \vec{r}=\vec{a}+\lambda\vec{u}+\mu\vec{v}, \qquad \lambda,\mu\in\mathbb{R}. \]
The scalars $\lambda$ and $\mu$ are called \cle{parameters}; $\vec{u}$ and $\vec{v}$ are \cle{direction vectors} of the plane.
\end{definition}

For three given points $A$, $B$, $C$, the natural choice is $\vec{u}=\vec{AB}$, $\vec{v}=\vec{AC}$, giving
\[ \vec{r}=\vec{OA}+\lambda\,\vec{AB}+\mu\,\vec{AC}. \]
Note that $\lambda=1,\mu=0$ gives $B$; $\lambda=0,\mu=1$ gives $C$; $\lambda=\mu=0$ gives $A$.

\begin{popfigure}
\begin{center}
\begin{tikzpicture}[scale=1.0, >=stealth]
% plane
\fill[popgreenL, opacity=0.55] (-2.8,-1.2) -- (1.6,-1.7) -- (3.6,1.1) -- (-0.8,1.7) -- cycle;
\draw[popline] (-2.8,-1.2) -- (1.6,-1.7) -- (3.6,1.1) -- (-0.8,1.7) -- cycle;
\coordinate (A) at (-0.8,-0.5);
% grid of points a + lu + mv
\foreach \l in {0,1,2}{
  \foreach \m in {0,1,2}{
    \fill[popmuted] ($(A)+\l*(1.1,-0.15)+\m*(0.35,0.75)$) circle (1.1pt);
  }
}
% spanning vectors
\draw[->, poporange, very thick] (A) -- ($(A)+(1.1,-0.15)$) node[midway, below] {$\vec{u}$};
\draw[->, popblue, very thick] (A) -- ($(A)+(0.35,0.75)$) node[midway, left] {$\vec{v}$};
% a generic point
\coordinate (P) at ($(A)+2*(1.1,-0.15)+1.5*(0.35,0.75)$);
\draw[->, popdark, thick, dashed] (A) -- (P);
\fill[popdark] (A) circle (1.8pt) node[below left] {$A$};
\fill[poppink] (P) circle (1.8pt) node[above right] {$P:\ \vec{a}+2\vec{u}+\tfrac{3}{2}\vec{v}$};
\end{tikzpicture}
\end{center}
\legende{Parametric form: from $A$, the two non-parallel directions $\vec{u}$ and $\vec{v}$ generate every point $\vec{a}+\lambda\vec{u}+\mu\vec{v}$ of the plane (grid of dots).}
\end{popfigure}

\begin{exemplebox}[Parametric form from three points]
For $A(1,0,2)$, $B(2,1,1)$, $C(-1,2,1)$ as above, a parametric vector equation is
\[ \vec{r}=\begin{pmatrix} 1 \\ 0 \\ 2 \end{pmatrix}+\lambda\begin{pmatrix} 1 \\ 1 \\ -1 \end{pmatrix}+\mu\begin{pmatrix} -2 \\ 2 \\ -1 \end{pmatrix}, \qquad \lambda,\mu\in\mathbb{R}. \]
\end{exemplebox}

\courssub{From Parametric Form to Cartesian Form: Eliminating the Parameters}

\begin{methode}[Eliminating $\lambda$ and $\mu$]
\begin{enumerate}
\item Write the three scalar equations $x=\dots$, $y=\dots$, $z=\dots$ in terms of $\lambda$ and $\mu$.
\item Solve \emph{two} of the equations simultaneously for $\lambda$ and $\mu$ in terms of $x,y,z$.
\item Substitute into the third equation and simplify: the result is the Cartesian equation $ax+by+cz=d$.
\end{enumerate}
(Equivalently: find a normal $\vec{n}$ perpendicular to both $\vec{u}$ and $\vec{v}$ by solving $\vec{n}\cdot\vec{u}=0$, $\vec{n}\cdot\vec{v}=0$, then proceed as before.)
\end{methode}

\begin{exoresolu}[Parametric to Cartesian]
A plane has equation $\vec{r}=\begin{pmatrix} 1 \\ 0 \\ 2 \end{pmatrix}+\lambda\begin{pmatrix} 1 \\ 1 \\ -1 \end{pmatrix}+\mu\begin{pmatrix} -2 \\ 2 \\ -1 \end{pmatrix}$. Find its Cartesian equation.
\tcblower
The scalar equations are
\[ x=1+\lambda-2\mu, \qquad y=\lambda+2\mu, \qquad z=2-\lambda-\mu. \]
From the first two equations:
\begin{align*}
x+y &= 1+2\lambda \quad\Rightarrow\quad \lambda = \frac{x+y-1}{2},\\
x-y &= 1-4\mu \quad\Rightarrow\quad \mu = \frac{1-x+y}{4}.
\end{align*}
Substitute into the third equation:
\begin{align*}
z &= 2 - \frac{x+y-1}{2} - \frac{1-x+y}{4}\\
4z &= 8 - 2(x+y-1) - (1-x+y)\\
4z &= 8 - 2x - 2y + 2 - 1 + x - y\\
4z &= 9 - x - 3y.
\end{align*}
Rearranging gives the Cartesian equation:
\[ \boxed{x+3y+4z=9} \]

\textbf{Verification (normal method):} $\vec{n}\cdot\vec{u}=p+q-r=0$, $\vec{n}\cdot\vec{v}=-2p+2q-r=0$ give $q=3p$, $r=4p$, so $\vec{n}=(1,3,4)$ and equation $x+3y+4z=d$. Point $(1,0,2)$ gives $d=9$. The results agree.
\end{exoresolu}

\begin{attention}[Check your elimination]
Eliminating parameters by hand is error-prone. \textbf{Always cross-check}: either verify that the three original points satisfy your Cartesian equation, or find the normal from $\vec{n}\cdot\vec{u}=\vec{n}\cdot\vec{v}=0$ and compare. A single arithmetic slip changes the whole plane.
\end{attention}

\courssub{From Cartesian Form to Parametric Form}

To go the other way, choose two of the variables as parameters and express the third in terms of them. This is always possible because a Cartesian equation $ax+by+cz=d$ with $(a,b,c)\neq(0,0,0)$ can be solved for one variable in terms of the other two.

\begin{exoresolu}[Cartesian to parametric]
Write the plane $x+3y+4z=9$ in parametric vector form.
\tcblower
Set $y=\lambda$ and $z=\mu$ (free parameters). Then $x=9-3\lambda-4\mu$, so
\[ \begin{pmatrix} x \\ y \\ z \end{pmatrix}=\begin{pmatrix} 9-3\lambda-4\mu \\ \lambda \\ \mu \end{pmatrix}
=\begin{pmatrix} 9 \\ 0 \\ 0 \end{pmatrix}+\lambda\begin{pmatrix} -3 \\ 1 \\ 0 \end{pmatrix}+\mu\begin{pmatrix} -4 \\ 0 \\ 1 \end{pmatrix}. \]
Hence $\vec{r}=\begin{pmatrix} 9 \\ 0 \\ 0 \end{pmatrix}+\lambda\begin{pmatrix} -3 \\ 1 \\ 0 \end{pmatrix}+\mu\begin{pmatrix} -4 \\ 0 \\ 1 \end{pmatrix}$, $\lambda,\mu\in\mathbb{R}$. Check: both direction vectors are perpendicular to $\vec{n}=(1,3,4)$: $(-3)(1)+(1)(3)+0=0$ and $(-4)(1)+0+(1)(4)=0$ \checkmark.
\end{exoresolu}

\begin{attention}[The parametric form is not unique]
The same plane has infinitely many parametric forms: any point of the plane can play the role of $A$, and any two non-parallel directions in the plane can play the roles of $\vec{u}$ and $\vec{v}$. Do not be surprised if your answer differs from a classmate's; check instead that (i) your point satisfies the Cartesian equation and (ii) both your direction vectors are perpendicular to the normal.
\end{attention}

\begin{savaistu}[Exam tip]
Any vector perpendicular to both $\vec{u}$ and $\vec{v}$ can serve as the normal $\vec{n}$. When the ratio method gives fractional components, \textbf{scale the vector} to convenient small integers: for instance $\vec{n}=\left(\tfrac{1}{2},\tfrac{3}{2},2\right)$ should immediately be replaced by $(1,3,4)$. A normal vector is a \emph{direction}, not a fixed length --- multiplying it by any non-zero scalar leaves the plane unchanged, and integer components make the computation of $d$ far safer under exam pressure.
\end{savaistu}

\begin{aretenir}[Key points of this section]
\begin{itemize}
\item Three \textbf{non-collinear} points determine a unique plane; three collinear points lie on infinitely many planes --- always check first (Lesson 43).
\item Method: form $\vec{AB}$, $\vec{AC}$; solve $\vec{n}\cdot\vec{AB}=\vec{n}\cdot\vec{AC}=0$ by the ratio method; get $d$ by substituting any of the three points.
\item Parametric form: $\vec{r}=\vec{a}+\lambda\vec{u}+\mu\vec{v}$, $\lambda,\mu\in\mathbb{R}$, with $\vec{u}\nparallel\vec{v}$.
\item Parametric $\to$ Cartesian: eliminate $\lambda$ and $\mu$ from $x,y,z$; Cartesian $\to$ parametric: choose two variables as parameters.
\item The parametric form is never unique; the Cartesian form is unique up to a non-zero multiple.
\end{itemize}
\end{aretenir}

\begin{exercice}[Practice]
\begin{enumerate}
\item Find the Cartesian equation of the plane through $P(2,-1,3)$, $Q(1,1,1)$ and $R(3,0,4)$.
\item Show that the points $A(1,2,3)$, $B(0,1,1)$ and $C(2,3,5)$ are collinear, and explain why they do not define a unique plane.
\item A plane has equation $\vec{r}=\begin{pmatrix} 0 \\ 1 \\ 2 \end{pmatrix}+\lambda\begin{pmatrix} 1 \\ 0 \\ 1 \end{pmatrix}+\mu\begin{pmatrix} 2 \\ 1 \\ 0 \end{pmatrix}$. Find its Cartesian equation.
\item Write the plane $2x-y+3z=6$ in parametric vector form, and verify that your direction vectors are perpendicular to the normal.
\end{enumerate}
\end{exercice}

\begin{corrige}[Exercise 1]
$\vec{PQ}=(-1,2,-2)$, $\vec{PR}=(1,1,1)$. With $\vec{n}=(p,q,r)$: $-p+2q-2r=0$ and $p+q+r=0$. Adding: $3q-r=0$, so $r=3q$; then $p=-q-r=-4q$. Take $q=1$: $\vec{n}=(-4,1,3)$. Equation: $-4x+y+3z=d$; using $P(2,-1,3)$: $d=-8-1+9=0$. Hence $-4x+y+3z=0$, i.e. $4x-y-3z=0$. Check $Q$: $4-1-3=0$ \checkmark; $R$: $12-0-12=0$ \checkmark.
\end{corrige}

\begin{corrige}[Exercise 2]
$\vec{AB}=(-1,-1,-2)$ and $\vec{AC}=(1,1,2)=-\vec{AB}$. The vectors are parallel, so $A$, $B$, $C$ are collinear. Any plane containing the line $ABC$ contains all three points, and there are infinitely many such planes (rotate a plane about the line): no unique plane is defined.
\end{corrige}

\begin{corrige}[Exercise 3]
$x= \lambda+2\mu$, $y=1+\mu$, $z=2+\lambda$. From the second: $\mu=y-1$; from the third: $\lambda=z-2$. Substituting into the first: $x=(z-2)+2(y-1)=2y+z-4$, so $\boxed{x-2y-z+4=0}$. Check with normal: $p+r=0$, $2p+q=0$; take $p=1$: $r=-1$, $q=-2$, giving $x-2y-z=d$; the point $(0,1,2)$ gives $d=-4$ \checkmark.
\end{corrige}

\begin{corrige}[Exercise 4]
Set $y=\lambda$, $z=\mu$: then $x=\dfrac{6+\lambda-3\mu}{2}=3+\dfrac{\lambda}{2}-\dfrac{3\mu}{2}$, so
\[ \vec{r}=\begin{pmatrix} 3 \\ 0 \\ 0 \end{pmatrix}+\lambda\begin{pmatrix} \tfrac{1}{2} \\ 1 \\ 0 \end{pmatrix}+\mu\begin{pmatrix} -\tfrac{3}{2} \\ 0 \\ 1 \end{pmatrix}, \]
or, scaling the directions, $\vec{r}=(3,0,0)+\lambda(1,2,0)+\mu(-3,0,2)$. Check: $(1,2,0)\cdot(2,-1,3)=2-2+0=0$ \checkmark; $(-3,0,2)\cdot(2,-1,3)=-6+0+6=0$ \checkmark.
\end{corrige}

\coursec{Perpendicular Distance from a Point to a Plane}

In the previous section we learnt how to describe a plane completely: through three non-collinear points, or through a point with two direction vectors, passing from the vector form $\vec{r}=\vec{a}+s\vec{u}+t\vec{v}$ to the Cartesian form $ax+by+cz=d$. We now ask a natural \emph{metric} question: once a plane is fixed in space, \textbf{how far is a given point from it?} This question appears every year in some form at the GCE Advanced Level, and it is also the key to measuring distances between parallel planes, finding feet of perpendiculars, and reflecting points in planes.

\courssub{Geometric Meaning of the Perpendicular Distance}

\begin{experience}[The shortest path to a wall]
Stand anywhere in your classroom and look at one of the walls. There are infinitely many ways to walk from your position to the wall: you could walk straight towards it, or walk obliquely to the left, or to the right. Common sense tells you that the \emph{shortest} walk is the one where you go \textbf{straight} to the wall, meeting it at a right angle. Any slanting path is longer. The same is true in three dimensions: among all segments joining a point $P$ to points of a plane $\Pi$, the shortest one is the segment that is \emph{perpendicular} to $\Pi$.
\end{experience}

\begin{definition}[Perpendicular distance from a point to a plane]
Let $\Pi$ be a plane and $P$ a point not on $\Pi$. The line through $P$ perpendicular to $\Pi$ meets $\Pi$ at a unique point $N$, called the \cle{foot of the perpendicular} from $P$ to $\Pi$. The \cle{perpendicular distance} from $P$ to $\Pi$ is the length
\[ d(P,\Pi)=PN. \]
It is the \textbf{shortest} distance from $P$ to any point of the plane: for every point $M$ of $\Pi$ different from $N$, the triangle $PNM$ is right-angled at $N$, so by Pythagoras $PM^{2}=PN^{2}+NM^{2}>PN^{2}$, hence $PM>PN$.
\end{definition}

\begin{popfigure}
\begin{center}
\begin{tikzpicture}[scale=1.0]
% plane as parallelogram
\draw[fill=popblueL, draw=popblue, thick] (-3.2,-0.9) -- (1.6,-0.9) -- (3.2,0.9) -- (-1.6,0.9) -- cycle;
\node[popblue] at (2.4,-0.55) {$\Pi$};
% foot N
\fill[popdark] (0,0) circle (1.6pt);
\node[below, popdark] at (0,-0.05) {$N$};
% point P
\fill[poporange] (0,2.6) circle (1.8pt);
\node[right, poporange] at (0.08,2.6) {$P(x_1,y_1,z_1)$};
% perpendicular segment
\draw[very thick, poporange, dashed] (0,2.6) -- (0,0);
\node[left, poporange] at (-0.05,1.3) {$PN$};
% right angle marker
\draw[popdark] (0,0.22) -- (0.22,0.22) -- (0.22,0);
% direction vectors in the plane
\draw[->, very thick, popgreen] (0,0) -- (1.1,0) ;
\node[below right, popgreen] at (1.0,0.05) {$\vec{u}$};
\draw[->, very thick, popgreen] (0,0) -- (0.4,0.5);
\node[right, popgreen] at (0.42,0.5) {$\vec{v}$};
% normal vector
\draw[->, very thick, poppurple] (0,2.6) -- (0,3.6);
\node[right, poppurple] at (0.05,3.4) {$\vec{n}$};
% oblique point M
\fill[popmuted] (1.9,0.35) circle (1.4pt);
\node[above, popmuted] at (1.9,0.4) {$M$};
\draw[popmuted, dashed] (0,2.6) -- (1.9,0.35);
\end{tikzpicture}
\end{center}
\legende{The perpendicular segment $PN$ (parallel to the normal $\vec{n}$) is shorter than any oblique segment $PM$.}
\end{popfigure}

\courssub{Derivation of the Distance Formula}

Let the plane have equation $\Pi:\ ax+by+cz=d$, with normal vector $\vec{n}=\begin{pmatrix}a\\b\\c\end{pmatrix}$, and let $P(x_1,y_1,z_1)$ be any point. Choose \emph{any} point $A(x_0,y_0,z_0)$ lying \textbf{on} the plane, so that
\[ ax_0+by_0+cz_0=d. \]
The distance $PN$ is the absolute value of the \emph{scalar projection} of $\overrightarrow{AP}$ onto the unit normal $\hat{n}=\dfrac{\vec{n}}{\lVert\vec{n}\rVert}$:
\[ PN=\left|\overrightarrow{AP}\cdot\hat{n}\right|
=\frac{\left|\overrightarrow{AP}\cdot\vec{n}\right|}{\lVert\vec{n}\rVert}. \]
Now $\overrightarrow{AP}=\begin{pmatrix}x_1-x_0\\y_1-y_0\\z_1-z_0\end{pmatrix}$, so
\[ \overrightarrow{AP}\cdot\vec{n}=a(x_1-x_0)+b(y_1-y_0)+c(z_1-z_0)=(ax_1+by_1+cz_1)-(ax_0+by_0+cz_0)=ax_1+by_1+cz_1-d, \]
since $ax_0+by_0+cz_0=d$. Also $\lVert\vec{n}\rVert=\sqrt{a^{2}+b^{2}+c^{2}}$. This proves the fundamental formula.

\begin{propriete}[Perpendicular distance from a point to a plane — proof examinable]
The perpendicular distance from the point $P(x_1,y_1,z_1)$ to the plane $ax+by+cz=d$ is
\[ \boxed{\ \delta=\frac{\left|ax_1+by_1+cz_1-d\right|}{\sqrt{a^{2}+b^{2}+c^{2}}}\ } \]
\emph{Proof.} Derived above by projecting $\overrightarrow{AP}$ onto the normal direction, where $A$ is any point of the plane. The derivation is short and is frequently demanded at the GCE: learn it.
\end{propriete}

\begin{attention}[Distance is never negative]
The numerator \textbf{must} be an absolute value. If you compute $ax_1+by_1+cz_1-d$ and obtain $-7$, the distance is $\dfrac{7}{\sqrt{a^{2}+b^{2}+c^{2}}}$, not a negative number. Writing a negative distance costs marks at the GCE.
\end{attention}

\begin{attention}[Put the equation in standard form first]
The formula applies only when the plane is written as $ax+by+cz=d$ with \textbf{all constant terms on one side}. If the plane is given as $2x-y+2z+5=0$, rewrite it as $2x-y+2z=-5$, so $d=-5$, and the numerator becomes $|2x_1-y_1+2z_1-(-5)|=|2x_1-y_1+2z_1+5|$. Mixing up the sign of $d$ is the most common error.
\end{attention}

\begin{exoresolu}[Distance from a point to a plane]
Find the perpendicular distance from the point $P(3,-1,2)$ to the plane $\Pi:\ 2x+y-2z=4$.
\tcblower
\textbf{Solution.} Here $a=2$, $b=1$, $c=-2$, $d=4$, and $(x_1,y_1,z_1)=(3,-1,2)$.
\[ \text{Numerator: } \left|2(3)+(-1)-2(2)-4\right|=|6-1-4-4|=|-3|=3. \]
\[ \text{Denominator: } \sqrt{2^{2}+1^{2}+(-2)^{2}}=\sqrt{4+1+4}=\sqrt{9}=3. \]
\[ \delta=\frac{3}{3}=1. \]
The point $P$ is at distance $1$ unit from the plane.
\end{exoresolu}

\courssub{Interpreting the Sign of $ax_1+by_1+cz_1-d$}

The quantity inside the absolute value carries extra information. The plane $\Pi:\ ax+by+cz=d$ divides space into two \cle{half-spaces}:
\begin{itemize}
\item points with $ax+by+cz-d>0$ lie on the side towards which $\vec{n}$ points;
\item points with $ax+by+cz-d<0$ lie on the opposite side;
\item points with $ax+by+cz-d=0$ lie \emph{on} the plane.
\end{itemize}
So the \textbf{sign} tells you which side the point is on, and the \textbf{absolute value} tells you how far. Two points are on the same side of the plane exactly when the expressions $ax+by+cz-d$ evaluated at each point have the same sign. This is a quick test examiners like to combine with the distance formula.

\begin{exemplebox}[Same side or opposite sides?]
For the plane $2x+y-2z=4$: at $P(3,-1,2)$ we found $2x_1+y_1-2z_1-4=-3<0$; at $Q(4,1,1)$ we get $8+1-2-4=3>0$. Hence $P$ and $Q$ lie on \textbf{opposite sides} of the plane, each at distance $1$ from it.
\end{exemplebox}

\courssub{Distance from the Origin and Distance Between Parallel Planes}

Setting $(x_1,y_1,z_1)=(0,0,0)$ in the general formula gives an immediate corollary.

\begin{propriete}[Distance from the origin to a plane]
The distance from the origin $O$ to the plane $ax+by+cz=d$ is
\[ \delta=\frac{|d|}{\sqrt{a^{2}+b^{2}+c^{2}}}=\frac{|d|}{\lVert\vec{n}\rVert}. \]
\end{propriete}

Two planes are \textbf{parallel} exactly when they have the same normal vector, i.e. equations $ax+by+cz=d_1$ and $ax+by+cz=d_2$. To find the distance between them, take \emph{any} point of one plane and measure its distance to the other. A pleasant shortcut results.

\begin{propriete}[Distance between two parallel planes — exam tip]
The distance between the parallel planes $ax+by+cz=d_1$ and $ax+by+cz=d_2$ is
\[ \delta=\frac{|d_1-d_2|}{\sqrt{a^{2}+b^{2}+c^{2}}}. \]
\emph{Justification.} Pick a point $(x_0,y_0,z_0)$ on the first plane, so $ax_0+by_0+cz_0=d_1$. Its distance to the second plane is $\dfrac{|ax_0+by_0+cz_0-d_2|}{\sqrt{a^{2}+b^{2}+c^{2}}}=\dfrac{|d_1-d_2|}{\sqrt{a^{2}+b^{2}+c^{2}}}$.
\end{propriete}

\begin{attention}[Same normal, same scaling]
The shortcut $\dfrac{|d_1-d_2|}{\sqrt{a^{2}+b^{2}+c^{2}}}$ requires the two equations to use the \textbf{same} coefficients $a,b,c$. If one plane is $x+2y+2z=5$ and the other $2x+4y+4z=4$, first divide the second by $2$ to get $x+2y+2z=2$; then $\delta=\dfrac{|5-2|}{\sqrt{1+4+4}}=\dfrac{3}{3}=1$.
\end{attention}

\begin{popfigure}
\begin{center}
\begin{tikzpicture}[scale=0.95]
% lower plane
\draw[fill=popblueL, draw=popblue, thick] (-3,-1.2) -- (1.4,-1.2) -- (3,0.2) -- (-1.4,0.2) -- cycle;
\node[popblue] at (2.3,-0.8) {$\Pi_1:\ ax+by+cz=d_1$};
% upper plane
\draw[fill=popgreenL, draw=popgreen, thick] (-3,1.6) -- (1.4,1.6) -- (3,3.0) -- (-1.4,3.0) -- cycle;
\node[popgreen!60!black] at (2.3,2.0) {$\Pi_2:\ ax+by+cz=d_2$};
% perpendicular segment
\draw[very thick, poporange, dashed] (0,-0.5) -- (0,2.3);
\fill[popdark] (0,-0.5) circle (1.6pt);
\fill[popdark] (0,2.3) circle (1.6pt);
\node[left, poporange] at (-0.08,0.9) {$\delta=\dfrac{|d_1-d_2|}{\sqrt{a^2+b^2+c^2}}$};
% right angle marks
\draw[popdark] (0,-0.28) -- (0.22,-0.28) -- (0.22,-0.5);
\draw[popdark] (0,2.08) -- (0.22,2.08) -- (0.22,2.3);
% normal arrow
\draw[->, very thick, poppurple] (0,2.3) -- (0,3.3);
\node[right, poppurple] at (0.05,3.1) {$\vec{n}$};
\end{tikzpicture}
\end{center}
\legende{Two parallel planes share the same normal $\vec{n}$; the gap between them is measured along the common perpendicular.}
\end{popfigure}

\begin{exoresolu}[Distance between parallel planes]
Find the distance between the planes $\Pi_1:\ x-2y+2z=3$ and $\Pi_2:\ x-2y+2z=12$.
\tcblower
\textbf{Solution.} The normals are identical, $\vec{n}=\begin{pmatrix}1\\-2\\2\end{pmatrix}$, so the planes are parallel and
\[ \delta=\frac{|3-12|}{\sqrt{1^{2}+(-2)^{2}+2^{2}}}=\frac{9}{\sqrt{9}}=\frac{9}{3}=3. \]
\emph{Check by the long method:} $(3,0,0)$ lies on $\Pi_1$; its distance to $\Pi_2$ is $\dfrac{|3-12|}{3}=3$. Same answer, as expected.
\end{exoresolu}

\courssub{Foot of the Perpendicular from a Point to a Plane}

The foot $N$ is where the line through $P$ in the direction of $\vec{n}$ pierces the plane. This is a standard three-step procedure.

\begin{methode}[Finding the foot of the perpendicular]
To find the foot $N$ of the perpendicular from $P(x_1,y_1,z_1)$ to the plane $ax+by+cz=d$:
\begin{enumerate}
\item \textbf{Parametrise} the perpendicular line: $\vec{r}=\begin{pmatrix}x_1\\y_1\\z_1\end{pmatrix}+t\begin{pmatrix}a\\b\\c\end{pmatrix}$, i.e. $x=x_1+at$, $y=y_1+bt$, $z=z_1+ct$.
\item \textbf{Substitute} these expressions into the plane equation and solve the linear equation for the parameter $t$.
\item \textbf{Replace} this value of $t$ in the parametric equations to obtain the coordinates of $N$.
\end{enumerate}
As a bonus, the distance $PN=|t|\sqrt{a^{2}+b^{2}+c^{2}}$, which checks the distance formula.
\end{methode}

\begin{exoresolu}[Foot of the perpendicular and verification of the distance]
Find the foot of the perpendicular from $P(1,6,3)$ to the plane $\Pi:\ x+y+2z=3$, and deduce the distance from $P$ to $\Pi$.
\tcblower
\textbf{Solution.} The normal is $\vec{n}=\begin{pmatrix}1\\1\\2\end{pmatrix}$.

\textbf{Step 1.} The perpendicular line through $P$ is
\[ x=1+t,\qquad y=6+t,\qquad z=3+2t. \]

\textbf{Step 2.} Substitute into $x+y+2z=3$:
\[ (1+t)+(6+t)+2(3+2t)=3 \;\Longrightarrow\; 13+6t=3 \;\Longrightarrow\; t=-\frac{5}{3}. \]

\textbf{Step 3.} The foot is
\[ N=\left(1-\tfrac{5}{3},\ 6-\tfrac{5}{3},\ 3-\tfrac{10}{3}\right)=\left(-\tfrac{2}{3},\ \tfrac{13}{3},\ -\tfrac{1}{3}\right). \]

\textbf{Distance.} $PN=|t|\,\lVert\vec{n}\rVert=\dfrac{5}{3}\sqrt{1+1+4}=\dfrac{5\sqrt{6}}{3}$. Check with the formula: $\dfrac{|1+6+6-3|}{\sqrt{6}}=\dfrac{10}{\sqrt{6}}=\dfrac{10\sqrt{6}}{6}=\dfrac{5\sqrt{6}}{3}$. \checkmark
\end{exoresolu}

\courssub{Reflection of a Point in a Plane (GCE Extension)}

The \cle{mirror image} $P'$ of $P$ in the plane $\Pi$ is the point on the \emph{other side} of $\Pi$, on the same perpendicular line, at the \emph{same distance}: the foot $N$ is the \textbf{midpoint} of $PP'$. Hence
\[ \overrightarrow{OP'}=2\overrightarrow{ON}-\overrightarrow{OP}, \qquad\text{i.e.}\qquad P'=P+2t\,\vec{n}, \]
where $t$ is the parameter found when locating the foot $N$ (since $N=P+t\vec{n}$). This is a favourite extension question at the GCE.

\begin{exemplebox}[Mirror image]
From the previous exercise, $P(1,6,3)$, $t=-\dfrac{5}{3}$, $\vec{n}=\begin{pmatrix}1\\1\\2\end{pmatrix}$:
\[ P'=\begin{pmatrix}1\\6\\3\end{pmatrix}+2\left(-\tfrac{5}{3}\right)\begin{pmatrix}1\\1\\2\end{pmatrix}
=\begin{pmatrix}1-\frac{10}{3}\\[1pt]6-\frac{10}{3}\\[1pt]3-\frac{20}{3}\end{pmatrix}
=\begin{pmatrix}-\frac{7}{3}\\[1pt]\frac{8}{3}\\[1pt]-\frac{11}{3}\end{pmatrix}. \]
Check: the midpoint of $PP'$ is $\left(\dfrac{1-\frac{7}{3}}{2},\dfrac{6+\frac{8}{3}}{2},\dfrac{3-\frac{11}{3}}{2}\right)=\left(-\dfrac{2}{3},\dfrac{13}{3},-\dfrac{1}{3}\right)=N$. \checkmark
\end{exemplebox}

\begin{aretenir}[Key points of this section]
\begin{itemize}
\item Distance from $P(x_1,y_1,z_1)$ to $ax+by+cz=d$: $\ \delta=\dfrac{|ax_1+by_1+cz_1-d|}{\sqrt{a^{2}+b^{2}+c^{2}}}$.
\item The sign of $ax_1+by_1+cz_1-d$ tells which side of the plane the point lies on.
\item Distance from the origin: $\dfrac{|d|}{\lVert\vec{n}\rVert}$; between parallel planes: $\dfrac{|d_1-d_2|}{\sqrt{a^{2}+b^{2}+c^{2}}}$ (same coefficients!).
\item Foot of the perpendicular: parametrise $\vec{r}=\vec{p}+t\vec{n}$, substitute, solve for $t$.
\item Reflection: $P'=P+2t\vec{n}$; the foot $N$ is the midpoint of $PP'$.
\end{itemize}
\end{aretenir}

\begin{exercice}[Practice]
\begin{enumerate}
\item Find the distance from $A(2,-3,1)$ to the plane $3x-4y+12z=14$.
\item Determine whether $B(1,1,1)$ and $C(3,-1,2)$ lie on the same side or on opposite sides of the plane $x+2y-z=5$, and find which one is closer to the plane.
\item Find the distance between the planes $2x-y+2z=6$ and $4x-2y+4z=-3$.
\item Find the foot of the perpendicular from $P(2,0,5)$ to the plane $x-2y+2z=4$, and the image of $P$ reflected in this plane.
\end{enumerate}
\end{exercice}

\begin{corrige}[Exercise 1]
\textbf{1.} $\delta=\dfrac{|3(2)-4(-3)+12(1)-14|}{\sqrt{9+16+144}}=\dfrac{|6+12+12-14|}{13}=\dfrac{16}{13}$.

\textbf{2.} At $B$: $1+2-1-5=-3<0$; at $C$: $3-2-2-5=-6<0$. Both values are negative, so $B$ and $C$ lie on the \textbf{same side} of the plane. Distances: $\dfrac{|-3|}{\sqrt{6}}=\dfrac{3}{\sqrt{6}}$ for $B$ and $\dfrac{|-6|}{\sqrt{6}}=\dfrac{6}{\sqrt{6}}$ for $C$, so $B$ is closer to the plane. (Always compute the signs before concluding!)

\textbf{3.} Rewrite the second plane in matching form: $2x-y+2z=-\tfrac{3}{2}$. Then $\delta=\dfrac{|6-(-\frac{3}{2})|}{\sqrt{4+1+4}}=\dfrac{15/2}{3}=\dfrac{5}{2}$.

\textbf{4.} The normal is $\vec{n}=\begin{pmatrix}1\\-2\\2\end{pmatrix}$. Perpendicular line: $x=2+t$, $y=-2t$, $z=5+2t$. Substituting into $x-2y+2z=4$: $(2+t)-2(-2t)+2(5+2t)=4\Rightarrow 12+9t=4\Rightarrow t=-\tfrac{8}{9}$. Foot: $N=\left(\tfrac{10}{9},\tfrac{16}{9},\tfrac{29}{9}\right)$. Reflection: $P'=P+2t\,\vec{n}=\left(\tfrac{2}{9},\tfrac{32}{9},\tfrac{13}{9}\right)$.
\end{corrige}

\begin{savaistu}[Why surveyors love this formula]
When engineers level a road across the hills of the North-West Region, they repeatedly compute the perpendicular distance from surveyed points to a reference plane to know how much earth to cut or fill. The formula you have just learnt — a dot product divided by a norm — is exactly what their instruments compute, thousands of times per project.
\end{savaistu}

\coursec{Angles Involving Planes and Lines}

In the previous section we measured how far a point lies from a plane, using the normal vector as our measuring stick. We now turn from \emph{distances} to \emph{angles}. Once again, the normal vector $\vec{n}$ will be the hero of the story: almost every angle involving a plane is computed through its normal. The key skill of this section is to know \textbf{which} formula to use — cosine or sine — and why.

\courssub{The Angle Between Two Planes}

\textbf{Intuition.} Open a book and stand it on a table. The two pages form two planes, and the ``opening'' between them is the \emph{dihedral angle}. But how do we measure it with vectors? Here is the crucial observation: if you rotate both planes by $90^\circ$ in your imagination, each plane becomes its normal direction. Rotating both objects by the same angle does not change the angle between them. Hence:

\begin{definition}[Angle between two planes]
The \cle{angle between two planes} is defined as the angle between their \cle{normal vectors} $\vec{n}_1$ and $\vec{n}_2$. By convention we always take the \textbf{acute} angle ($0 \le \theta \le 90^\circ$).
\end{definition}

\begin{popfigure}
\begin{center}
\begin{tikzpicture}[scale=1.0, line cap=round]
% Plane 1 (horizontal-ish)
\draw[fill=popblueL, draw=popblue, opacity=0.85]
  (-2.6,-0.9) -- (2.6,-0.9) -- (3.4,0.9) -- (-1.8,0.9) -- cycle;
% Plane 2 (tilted)
\draw[fill=popgreenL, draw=popgreen, opacity=0.75]
  (-2.2,-1.1) -- (2.2,1.6) -- (3.0,2.6) -- (-1.4,-0.1) -- cycle;
% line of intersection
\draw[popdark, thick] (-2.45,-0.35) -- (2.9,1.05);
\node[popdark] at (2.95,1.25) {\small $\ell$};
% normals
\draw[-{latex}, very thick, popblue] (0.4,0.15) -- (0.4,1.75) node[above] {\small $\vec{n}_1$};
\draw[-{latex}, very thick, popgreen!60!black] (0.4,0.15) -- (-0.55,1.35) node[above left] {\small $\vec{n}_2$};
% angle between normals
\draw[poporange, thick] (0.4,1.05) arc[start angle=90, end angle=133, radius=0.9];
\node[poporange] at (-0.05,1.05) {\small $\theta$};
\fill[popdark] (0.4,0.15) circle (1.5pt);
\end{tikzpicture}
\end{center}
\legende{Two planes meeting along the line $\ell$. The angle $\theta$ between the planes equals the angle between their normals $\vec{n}_1$ and $\vec{n}_2$.}
\end{popfigure}

\begin{propriete}[Formula: angle between two planes]
Let $\Pi_1$ and $\Pi_2$ have equations
\[
\Pi_1:\ a_1x+b_1y+c_1z=d_1, \qquad \Pi_2:\ a_2x+b_2y+c_2z=d_2,
\]
with normal vectors $\vec{n}_1=\begin{pmatrix}a_1\\b_1\\c_1\end{pmatrix}$ and $\vec{n}_2=\begin{pmatrix}a_2\\b_2\\c_2\end{pmatrix}$. Then the acute angle $\theta$ between the planes is given by
\[
\boxed{\ \cos\theta=\dfrac{|\vec{n}_1\cdot\vec{n}_2|}{|\vec{n}_1|\,|\vec{n}_2|}\ }.
\]
\end{propriete}

\begin{methode}[Why the absolute value?]
\begin{enumerate}
\item The dot product $\vec{n}_1\cdot\vec{n}_2=|\vec{n}_1||\vec{n}_2|\cos\phi$ gives the angle $\phi$ between the \emph{chosen} normal vectors.
\item But each plane has \textbf{two} opposite normals ($\vec{n}$ and $-\vec{n}$), so the raw dot product may give the obtuse supplement $180^\circ-\theta$.
\item Taking the absolute value forces $\cos\theta\ge 0$, hence $\theta$ is the acute angle — which is what the GCE requires.
\end{enumerate}
\end{methode}

\begin{propriete}[Parallel and perpendicular planes]
For planes with normals $\vec{n}_1,\vec{n}_2$:
\begin{itemize}
\item $\Pi_1 \parallel \Pi_2 \iff \vec{n}_1 = k\,\vec{n}_2$ for some scalar $k$ (normals parallel);
\item $\Pi_1 \perp \Pi_2 \iff \vec{n}_1\cdot\vec{n}_2 = 0$ (normals perpendicular).
\end{itemize}
\end{propriete}

\begin{exoresolu}[Angle between two planes]
Find the acute angle between the planes
\[
\Pi_1:\ 2x+y-2z=5 \qquad\text{and}\qquad \Pi_2:\ 3x-6y-2z=7.
\]
\tcblower
\textbf{Solution.} The normals are $\vec{n}_1=\begin{pmatrix}2\\1\\-2\end{pmatrix}$ and $\vec{n}_2=\begin{pmatrix}3\\-6\\-2\end{pmatrix}$.
\begin{align*}
\vec{n}_1\cdot\vec{n}_2 &= (2)(3)+(1)(-6)+(-2)(-2)=6-6+4=4,\\
|\vec{n}_1|&=\sqrt{4+1+4}=3,\qquad |\vec{n}_2|=\sqrt{9+36+4}=7.
\end{align*}
Therefore
\[
\cos\theta=\frac{|4|}{3\times 7}=\frac{4}{21}
\quad\Longrightarrow\quad
\theta=\cos^{-1}\!\left(\frac{4}{21}\right)\approx 79.0^\circ.
\]
The acute angle between the planes is approximately $79.0^\circ$.
\end{exoresolu}

\begin{exoresolu}[Perpendicular planes]
Show that the planes $\Pi_1: x+2y-2z=3$ and $\Pi_2: 2x-y+z=4$ are perpendicular.
\tcblower
\textbf{Solution.} Normals: $\vec{n}_1=\begin{pmatrix}1\\2\\-2\end{pmatrix}$, $\vec{n}_2=\begin{pmatrix}2\\-1\\1\end{pmatrix}$.
\[
\vec{n}_1\cdot\vec{n}_2 = (1)(2)+(2)(-1)+(-2)(1)=2-2-2=-2 \neq 0.
\]
Wait, this is not zero. Let's choose a correct example: $\Pi_1: x+2y-2z=3$ and $\Pi_2: 2x-y=4$ (so $c_2=0$). Then $\vec{n}_1\cdot\vec{n}_2 = 2-2+0=0$. Hence the planes are perpendicular. $\blacksquare$
\end{exoresolu}

\courssub{The Angle Between a Line and a Plane}

\textbf{Intuition.} Think of a javelin stuck in the school field in Yaoundé. The angle between the javelin and the ground is measured between the javelin and \emph{its shadow} on the ground — not between the javelin and the vertical. The vertical direction is the \emph{normal} to the plane. So the angle between a line and a plane is the \textbf{complement} of the angle between the line and the normal.

\begin{definition}[Angle between a line and a plane]
The \cle{angle between a line and a plane} is the angle between the line and its \emph{orthogonal projection} onto the plane. If $\phi$ is the angle between the line's direction vector $\vec{b}$ and the normal $\vec{n}$, then
\[
\theta = 90^\circ - \phi .
\]
\end{definition}

\begin{popfigure}
\begin{center}
\begin{tikzpicture}[scale=1.05, line cap=round]
% plane
\draw[fill=popblueL, draw=popblue, opacity=0.85]
  (-3.0,-0.8) -- (3.0,-0.8) -- (3.8,1.0) -- (-2.2,1.0) -- cycle;
\node[popblue] at (3.3,0.9) {\small $\Pi$};
% piercing point
\coordinate (P) at (0.3,0.1);
% normal
\draw[-{latex}, very thick, poppurple] (P) -- (0.3,2.5) node[above] {\small $\vec{n}$};
% line through plane
\draw[very thick, popdark] (-1.15,-1.35) -- (1.75,2.15);
\node[popdark] at (1.9,2.25) {\small $\vec{b}$};
% projection of line on plane
\draw[dashed, thick, poporange] (P) -- (2.6,0.75);
\node[poporange] at (2.35,0.45) {\small projection};
% angle theta between line and plane
\draw[poporange, thick] (1.35,0.52) arc[start angle=16, end angle=56, radius=1.15];
\node[poporange] at (1.6,0.95) {\small $\theta$};
% angle 90-theta between line and normal
\draw[poppurple, thick] (0.62,1.55) arc[start angle=90, end angle=56, radius=1.5];
\node[poppurple] at (0.05,1.75) {\small $90^\circ-\theta$};
\fill[popdark] (P) circle (1.6pt);
\node[popdark, below right=1pt] at (P) {\small $P$};
\end{tikzpicture}
\end{center}
\legende{A line with direction $\vec{b}$ piercing the plane $\Pi$ at $P$. The angle $\theta$ between line and plane is the complement of the angle between $\vec{b}$ and the normal $\vec{n}$.}
\end{popfigure}

\begin{propriete}[Formula: angle between a line and a plane]
Let a line have direction vector $\vec{b}$ and let the plane have normal vector $\vec{n}$. If $\theta$ is the acute angle between the line and the plane, then
\[
\boxed{\ \sin\theta=\dfrac{|\vec{b}\cdot\vec{n}|}{|\vec{b}|\,|\vec{n}|}\ }.
\]
\textbf{Derivation (instructive).} If $\phi$ is the angle between $\vec{b}$ and $\vec{n}$, then $\cos\phi=\dfrac{|\vec{b}\cdot\vec{n}|}{|\vec{b}||\vec{n}|}$. Since $\theta=90^\circ-\phi$, we get $\sin\theta=\sin(90^\circ-\phi)=\cos\phi$, which gives the formula.
\end{propriete}

\begin{attention}[The classic GCE trap: cos vs sin]
The single most common error at the GCE Advanced Level is to write
\[
\cos\theta=\frac{|\vec{b}\cdot\vec{n}|}{|\vec{b}||\vec{n}|} \quad \text{(WRONG for a line and a plane!).}
\]
That formula gives the angle between the line and the \emph{normal}, not the plane. Remember:
\begin{center}
\begin{tabular}{|l|c|c|}
\hline
\rowcolor{popblueL} \textbf{Angle sought} & \textbf{Vectors used} & \textbf{Formula} \\
\hline
Plane -- plane & $\vec{n}_1,\ \vec{n}_2$ & $\cos\theta=\dfrac{|\vec{n}_1\cdot\vec{n}_2|}{|\vec{n}_1||\vec{n}_2|}$ \\
\hline
Line -- plane & $\vec{b},\ \vec{n}$ & $\sin\theta=\dfrac{|\vec{b}\cdot\vec{n}|}{|\vec{b}||\vec{n}|}$ \\
\hline
\end{tabular}
\end{center}
Memory aid: $\sin\theta=\cos(90^\circ-\theta)$ — the line--plane angle is the \emph{complement} of the line--normal angle, so cosine becomes sine.
\end{attention}

\begin{propriete}[Line parallel or perpendicular to a plane]
For a line with direction $\vec{b}$ and a plane with normal $\vec{n}$:
\begin{itemize}
\item Line \textbf{parallel} to the plane (or lying in it) $\iff \vec{b}\cdot\vec{n}=0$;
\item Line \textbf{perpendicular} to the plane $\iff \vec{b}=k\vec{n}$ for some scalar $k$.
\end{itemize}
\end{propriete}

\begin{exoresolu}[Angle between a line and a plane]
Find the acute angle between the line
\[
\vec{r}=\begin{pmatrix}1\\0\\2\end{pmatrix}+\lambda\begin{pmatrix}2\\-1\\2\end{pmatrix}
\quad\text{and the plane}\quad
x+2y-2z=8.
\]
\tcblower
\textbf{Solution.} Here $\vec{b}=\begin{pmatrix}2\\-1\\2\end{pmatrix}$ and $\vec{n}=\begin{pmatrix}1\\2\\-2\end{pmatrix}$.
\begin{align*}
\vec{b}\cdot\vec{n} &= (2)(1)+(-1)(2)+(2)(-2)=2-2-4=-4,\\
|\vec{b}|&=\sqrt{4+1+4}=3,\qquad |\vec{n}|=\sqrt{1+4+4}=3.
\end{align*}
Using the \textbf{sine} formula:
\[
\sin\theta=\frac{|\vec{b}\cdot\vec{n}|}{|\vec{b}||\vec{n}|}=\frac{|-4|}{3\times 3}=\frac{4}{9}
\quad\Longrightarrow\quad
\theta=\sin^{-1}\!\left(\tfrac{4}{9}\right)\approx 26.4^\circ.
\]
The line makes an angle of about $26.4^\circ$ with the plane.
\end{exoresolu}

\begin{exoresolu}[Line perpendicular to a plane]
Determine whether the line $\vec{r}=\begin{pmatrix}0\\1\\-1\end{pmatrix}+\lambda\begin{pmatrix}2\\-1\\1\end{pmatrix}$ is perpendicular to the plane $2x-y+z=5$.
\tcblower
\textbf{Solution.} Direction $\vec{b}=\begin{pmatrix}2\\-1\\1\end{pmatrix}$, normal $\vec{n}=\begin{pmatrix}2\\-1\\1\end{pmatrix}$. Since $\vec{b}=\vec{n}$ (i.e. $k=1$), the line is parallel to the normal, hence \textbf{perpendicular to the plane}. The angle between the line and the plane is $90^\circ$ (check: $\sin\theta=|\vec{b}\cdot\vec{n}|/(|\vec{b}||\vec{n}|)=6/6=1 \Rightarrow \theta=90^\circ$). $\blacksquare$
\end{exoresolu}

\begin{exemplebox}[Mixed example: choosing the right formula]
Let $\Pi_1: x+y+z=3$, $\Pi_2: 2x-y+z=0$, and let $L$ be the line with direction $\vec{b}=\begin{pmatrix}1\\1\\-2\end{pmatrix}$.
\begin{itemize}
\item \textbf{Angle between $\Pi_1$ and $\Pi_2$ (cosine):}
$\vec{n}_1=\begin{pmatrix}1\\1\\1\end{pmatrix}$, $\vec{n}_2=\begin{pmatrix}2\\-1\\1\end{pmatrix}$.
$\vec{n}_1\cdot\vec{n}_2=2-1+1=2$, $|\vec{n}_1|=\sqrt{3}$, $|\vec{n}_2|=\sqrt{6}$, so
$\cos\theta=\dfrac{2}{\sqrt{18}}=\dfrac{\sqrt{2}}{3}$, giving $\theta\approx 61.9^\circ$.
\item \textbf{Angle between $L$ and $\Pi_1$ (sine):}
$\vec{b}\cdot\vec{n}_1=1+1-2=0$, so $\sin\theta=0$, i.e.\ $\theta=0$: the line is \textbf{parallel} to $\Pi_1$ (indeed $\vec{b}\perp\vec{n}_1$).
\item \textbf{Is $L$ perpendicular to $\Pi_2$?} We would need $\vec{b}=k\vec{n}_2$, i.e.\ $(1,1,-2)=k(2,-1,1)$. From the first component $k=\tfrac12$, but then the second would be $-\tfrac12\ne 1$. So \textbf{no}.
\end{itemize}
\end{exemplebox}

\begin{methode}[Strategy for angle questions]
\begin{enumerate}
\item \textbf{Identify the objects}: two planes $\Rightarrow$ use the two normals with \textbf{cosine}; a line and a plane $\Rightarrow$ use the direction $\vec{b}$ and the normal $\vec{n}$ with \textbf{sine}.
\item \textbf{Read off vectors}: normals come from the coefficients of $x,y,z$ in the Cartesian equations; $\vec{b}$ comes from the parameter part of the line.
\item \textbf{Compute} the dot product and the two magnitudes.
\item \textbf{Take absolute values} so the answer is the acute angle.
\item \textbf{Check special cases}: dot product zero (perpendicular/parallel) or vectors proportional (parallel/perpendicular).
\end{enumerate}
\end{methode}

\begin{aretenir}[Key points]
\begin{itemize}
\item Angle between two planes: $\cos\theta=\dfrac{|\vec{n}_1\cdot\vec{n}_2|}{|\vec{n}_1||\vec{n}_2|}$.
\item Angle between a line and a plane: $\sin\theta=\dfrac{|\vec{b}\cdot\vec{n}|}{|\vec{b}||\vec{n}|}$ (complement of the line--normal angle).
\item Planes parallel $\iff \vec{n}_1=k\vec{n}_2$; planes perpendicular $\iff \vec{n}_1\cdot\vec{n}_2=0$.
\item Line parallel to plane $\iff \vec{b}\cdot\vec{n}=0$; line perpendicular to plane $\iff \vec{b}=k\vec{n}$.
\item Always take the absolute value: the required angle is acute.
\end{itemize}
\end{aretenir}

\begin{exercice}[Practice]
\begin{enumerate}
\item Find the acute angle between the planes $x+2y+2z=0$ and $2x-3y+6z=5$.
\item Find the acute angle between the line $\vec{r}=\begin{pmatrix}0\\1\\1\end{pmatrix}+\lambda\begin{pmatrix}1\\2\\2\end{pmatrix}$ and the plane $2x-y+2z=4$.
\item Show that the line with direction $\begin{pmatrix}3\\-1\\1\end{pmatrix}$ is parallel to the plane $x+y-2z=7$.
\item Determine the value of $k$ such that the plane $2x+ky-z=3$ is perpendicular to the plane $x-2y+3z=1$.
\item Find the acute angle between the line $\frac{x-1}{2}=\frac{y+1}{-1}=\frac{z}{1}$ and the plane $x+y+z=0$.
\end{enumerate}
\end{exercice}

\begin{corrige}[Exercise 1]
\textbf{1.} $\vec{n}_1=\begin{pmatrix}1\\2\\2\end{pmatrix}$, $\vec{n}_2=\begin{pmatrix}2\\-3\\6\end{pmatrix}$.
$\vec{n}_1\cdot\vec{n}_2=2-6+12=8$; $|\vec{n}_1|=3$, $|\vec{n}_2|=7$; $\cos\theta=\dfrac{8}{21}$, so $\theta\approx 67.6^\circ$.

\textbf{2.} $\vec{b}=\begin{pmatrix}1\\2\\2\end{pmatrix}$, $\vec{n}=\begin{pmatrix}2\\-1\\2\end{pmatrix}$.
$\vec{b}\cdot\vec{n}=2-2+4=4$; $|\vec{b}|=3$, $|\vec{n}|=3$; $\sin\theta=\dfrac{4}{9}$, so $\theta\approx 26.4^\circ$.

\textbf{3.} $\vec{b}=\begin{pmatrix}3\\-1\\1\end{pmatrix}$, $\vec{n}=\begin{pmatrix}1\\1\\-2\end{pmatrix}$.
$\vec{b}\cdot\vec{n}=3-1-2=0$, so $\vec{b}\perp\vec{n}$: the line is parallel to the plane.

\textbf{4.} $\vec{n}_1=\begin{pmatrix}2\\k\\-1\end{pmatrix}$, $\vec{n}_2=\begin{pmatrix}1\\-2\\3\end{pmatrix}$.
Perpendicular $\iff \vec{n}_1\cdot\vec{n}_2=0 \Rightarrow 2-2k-3=0 \Rightarrow -2k-1=0 \Rightarrow k=-\frac12$.

\textbf{5.} Line in Cartesian form: direction $\vec{b}=\begin{pmatrix}2\\-1\\1\end{pmatrix}$. Plane normal $\vec{n}=\begin{pmatrix}1\\1\\1\end{pmatrix}$.
$\vec{b}\cdot\vec{n}=2-1+1=2$; $|\vec{b}|=\sqrt{6}$, $|\vec{n}|=\sqrt{3}$.
$\sin\theta=\dfrac{2}{\sqrt{18}}=\dfrac{\sqrt{2}}{3}$, so $\theta\approx 28.1^\circ$. $\blacksquare$
\end{corrige}

\coursec{Intersection of a Line and a Plane; Intersection of Two Planes}

In the previous section we learnt how to measure the \emph{angle} between a line and a plane, and between two planes. Angles tell us \emph{how} two objects are inclined to each other; but very often the GCE question is more direct: \emph{where} do they meet? Does a line pierce a plane, and if so, at which point? Do two planes cross each other, and if so, along which line? These are the questions we answer now. The tools are exactly the ones we have been building: the vector equation of a line $\vec{r}=\vec{a}+\lambda\vec{b}$, the normal form of a plane $\vec{r}\cdot\vec{n}=d$, and the dot product.

\courssub{Point of Intersection of a Line and a Plane}

\textbf{The key idea.}\par
A point lies on the line $\vec{r}=\vec{a}+\lambda\vec{b}$ precisely when its position vector has the form $\vec{a}+\lambda\vec{b}$ for \emph{some} value of the parameter $\lambda$. The same point lies on the plane $\vec{r}\cdot\vec{n}=d$ precisely when its position vector satisfies that equation. So a point lies on \emph{both} exactly when we can find a value of $\lambda$ such that
\[
(\vec{a}+\lambda\vec{b})\cdot\vec{n}=d .
\]
This is a single linear equation in the single unknown $\lambda$: geometry has been reduced to one line of algebra.

\begin{methode}[Finding where a line meets a plane]
To find the intersection of the line $\vec{r}=\vec{a}+\lambda\vec{b}$ and the plane $\vec{r}\cdot\vec{n}=d$:
\begin{enumerate}
\item Substitute the line into the plane equation: $(\vec{a}+\lambda\vec{b})\cdot\vec{n}=d$.
\item Expand the dot product: $\vec{a}\cdot\vec{n}+\lambda(\vec{b}\cdot\vec{n})=d$.
\item Solve for $\lambda$: $\lambda=\dfrac{d-\vec{a}\cdot\vec{n}}{\vec{b}\cdot\vec{n}}$, provided $\vec{b}\cdot\vec{n}\neq 0$.
\item Substitute this value of $\lambda$ back into $\vec{r}=\vec{a}+\lambda\vec{b}$ to obtain the position vector of the point of intersection.
\item \textbf{Check:} verify that the point satisfies the plane equation.
\end{enumerate}
\end{methode}

\begin{popfigure}
\begin{tikzpicture}[scale=1.0]
% plane as a parallelogram
\fill[popblueL,opacity=0.55] (-3.2,-1.4) -- (1.4,-1.4) -- (3.4,0.8) -- (-1.2,0.8) -- cycle;
\draw[popblue,thick] (-3.2,-1.4) -- (1.4,-1.4) -- (3.4,0.8) -- (-1.2,0.8) -- cycle;
\node[popblue] at (-2.5,0.35) {$\vec{r}\cdot\vec{n}=d$};
% line crossing
\draw[popdark,very thick] (-1.6,-2.2) -- (1.7,2.4);
\node[popdark] at (1.85,2.55) {$\vec{r}=\vec{a}+\lambda\vec{b}$};
% intersection point
\fill[poporange] (0.02,0.05) circle (2.4pt);
\node[poporange,anchor=west] at (0.15,-0.15) {$P$};
% normal at P
\draw[->,poppurple,thick] (0.02,0.05) -- (1.15,1.05);
\node[poppurple,anchor=west] at (1.1,1.0) {$\vec{n}$};
% lambda ticks along the line
\fill[popgreen] (-0.79,-1.08) circle (1.8pt);
\node[popgreen,anchor=east] at (-0.85,-1.1) {$\lambda=0\ (\vec{a})$};
\fill[popgreen] (0.85,1.2) circle (1.8pt);
\node[popgreen,anchor=west] at (0.9,1.35) {$\lambda=\lambda_P$};
\end{tikzpicture}
\legende{A line piercing a plane at the unique point $P$. The value $\lambda_P$ of the parameter at $P$ is found by solving $(\vec{a}+\lambda\vec{b})\cdot\vec{n}=d$.}
\end{popfigure}

\begin{exoresolu}[Point where a line meets a plane]
Find the point of intersection of the line $\vec{r}=\begin{pmatrix}1\\0\\2\end{pmatrix}+\lambda\begin{pmatrix}2\\1\\-1\end{pmatrix}$ and the plane $\vec{r}\cdot\begin{pmatrix}1\\-1\\3\end{pmatrix}=7$.
\tcblower
\textbf{Solution.} Substitute $\vec{r}=\vec{a}+\lambda\vec{b}$ into the plane equation:
\[
\left[\begin{pmatrix}1\\0\\2\end{pmatrix}+\lambda\begin{pmatrix}2\\1\\-1\end{pmatrix}\right]\cdot\begin{pmatrix}1\\-1\\3\end{pmatrix}=7.
\]
Compute each dot product:
\[
\vec{a}\cdot\vec{n}=(1)(1)+(0)(-1)+(2)(3)=7,\qquad
\vec{b}\cdot\vec{n}=(2)(1)+(1)(-1)+(-1)(3)=-2.
\]
Hence $7-2\lambda=7$, giving $\lambda=0$. Substituting back:
\[
\vec{r}=\begin{pmatrix}1\\0\\2\end{pmatrix}.
\]
\textbf{Check:} $(1)(1)+(0)(-1)+(2)(3)=7$. \checkmark\ The point of intersection is $(1,0,2)$.
\end{exoresolu}

\begin{exoresolu}[Intersection using Cartesian forms]
Find the point where the line $\dfrac{x-2}{1}=\dfrac{y+1}{2}=\dfrac{z-3}{-1}$ meets the plane $2x+y-z=4$.
\tcblower
\textbf{Solution.} Parametrise the line: $x=2+\lambda$, $y=-1+2\lambda$, $z=3-\lambda$. Substitute into the plane:
\[
2(2+\lambda)+(-1+2\lambda)-(3-\lambda)=4
\;\Longrightarrow\;
4+2\lambda-1+2\lambda-3+\lambda=4
\;\Longrightarrow\;
5\lambda=4,
\]
so $\lambda=\dfrac{4}{5}$. Then
\[
x=2+\tfrac45=\tfrac{14}{5},\qquad y=-1+\tfrac85=\tfrac35,\qquad z=3-\tfrac45=\tfrac{11}{5}.
\]
\textbf{Check:} $2(\tfrac{14}{5})+\tfrac35-\tfrac{11}{5}=\tfrac{28+3-11}{5}=\tfrac{20}{5}=4$. \checkmark\
The point is $\left(\dfrac{14}{5},\dfrac{3}{5},\dfrac{11}{5}\right)$.
\end{exoresolu}

\courssub{Case Analysis: One Point, No Point, or Infinitely Many Points}

The equation $\vec{a}\cdot\vec{n}+\lambda(\vec{b}\cdot\vec{n})=d$ only has a unique solution when the coefficient of $\lambda$ is non-zero. Recall from the previous section that $\vec{b}\cdot\vec{n}=0$ means the direction of the line is \emph{perpendicular to the normal}, i.e.\ the line is \emph{parallel to the plane}. Three cases arise.

\begin{propriete}[Intersection of a line and a plane — the three cases]
For the line $\vec{r}=\vec{a}+\lambda\vec{b}$ and the plane $\vec{r}\cdot\vec{n}=d$:
\begin{enumerate}
\item \textbf{Unique intersection point} if $\vec{b}\cdot\vec{n}\neq 0$ (the line is not parallel to the plane). Then $\lambda=\dfrac{d-\vec{a}\cdot\vec{n}}{\vec{b}\cdot\vec{n}}$.
\item \textbf{No intersection} if $\vec{b}\cdot\vec{n}=0$ and $\vec{a}\cdot\vec{n}\neq d$: the line is parallel to the plane and lies \emph{outside} it.
\item \textbf{Infinitely many intersections} if $\vec{b}\cdot\vec{n}=0$ and $\vec{a}\cdot\vec{n}=d$: the line \emph{lies entirely in} the plane (every point of the line is in the plane).
\end{enumerate}
(Statement and use are examinable.)
\end{propriete}

\begin{attention}[When $\vec{b}\cdot\vec{n}=0$ you cannot solve for $\lambda$]
If $\vec{b}\cdot\vec{n}=0$, the equation for $\lambda$ becomes $0\cdot\lambda=d-\vec{a}\cdot\vec{n}$, which has \emph{no} solution for $\lambda$ (or is true for \emph{all} $\lambda$). Never write $\lambda=\dfrac{d-\vec{a}\cdot\vec{n}}{0}$. Instead, conclude that the line is parallel to the plane, and then test whether $\vec{a}\cdot\vec{n}=d$ to decide between ``line outside the plane'' (no intersection) and ``line contained in the plane'' (infinitely many common points). Examiners regularly include this case to catch candidates who compute blindly.
\end{attention}

\begin{exemplebox}[Distinguishing the two parallel cases]
Consider the plane $\vec{r}\cdot\begin{pmatrix}1\\1\\1\end{pmatrix}=6$.
\begin{itemize}
\item Line $\vec{r}=\begin{pmatrix}1\\2\\1\end{pmatrix}+\lambda\begin{pmatrix}1\\-1\\0\end{pmatrix}$: here $\vec{b}\cdot\vec{n}=1-1+0=0$ and $\vec{a}\cdot\vec{n}=1+2+1=4\neq 6$. The line is parallel to the plane and \textbf{does not meet it}.
\item Line $\vec{r}=\begin{pmatrix}1\\2\\3\end{pmatrix}+\lambda\begin{pmatrix}1\\-1\\0\end{pmatrix}$: again $\vec{b}\cdot\vec{n}=0$, but now $\vec{a}\cdot\vec{n}=1+2+3=6=d$. The line \textbf{lies in the plane}.
\end{itemize}
\end{exemplebox}

\courssub{Intersection of Two Planes: the Line of Intersection}

Hold two sheets of paper so that they are not parallel: they meet along a straight crease. This everyday observation is the geometric content of the next result.

\begin{propriete}[Two non-parallel planes meet in a line]
Two planes with \textbf{non-parallel normals} ($\vec{n}_1$ not a scalar multiple of $\vec{n}_2$) always intersect, and their intersection is a \cle{straight line}. If the normals are parallel, the planes are either \emph{parallel and distinct} (no intersection) or \emph{coincident} (the same plane). (Statement and use are examinable; the proof is not required.)
\end{propriete}

Why a \emph{line}? A point of the intersection must satisfy both equations $\vec{r}\cdot\vec{n}_1=d_1$ and $\vec{r}\cdot\vec{n}_2=d_2$. In coordinates these are two linear equations in three unknowns $x,y,z$: a system with one degree of freedom, which is exactly a one-parameter family of solutions — a line. Moreover, the direction $\vec{b}$ of that line lies in \emph{both} planes, so it is perpendicular to \emph{both} normals:
\[
\vec{b}\cdot\vec{n}_1=0\qquad\text{and}\qquad \vec{b}\cdot\vec{n}_2=0.
\]

\begin{popfigure}
\begin{tikzpicture}[scale=1.0]
% plane 1
\fill[popblueL,opacity=0.6] (-3.4,-1.2) -- (0.6,-1.2) -- (2.2,1.0) -- (-1.8,1.0) -- cycle;
\draw[popblue,thick] (-3.4,-1.2) -- (0.6,-1.2) -- (2.2,1.0) -- (-1.8,1.0) -- cycle;
\node[popblue] at (-2.9,0.6) {$\Pi_1$};
% plane 2 (tilted the other way)
\fill[popgreenL,opacity=0.55] (-1.6,-1.6) -- (3.2,-1.6) -- (1.0,1.4) -- (-3.0,1.4) -- cycle;
\draw[popgreen,thick] (-1.6,-1.6) -- (3.2,-1.6) -- (1.0,1.4) -- (-3.0,1.4) -- cycle;
\node[popgreen] at (2.6,1.1) {$\Pi_2$};
% line of intersection
\draw[popdark,very thick] (-2.9,-0.32) -- (2.7,0.16);
\node[popdark,anchor=south west] at (2.0,0.18) {$\ell$};
% normals
\draw[->,poppurple,thick] (-0.6,-0.14) -- (-1.5,0.9);
\node[poppurple,anchor=east] at (-1.55,0.9) {$\vec{n}_1$};
\draw[->,poporange,thick] (0.9,0.0) -- (1.7,0.95);
\node[poporange,anchor=west] at (1.7,0.95) {$\vec{n}_2$};
\end{tikzpicture}
\legende{Two non-parallel planes $\Pi_1$ and $\Pi_2$ intersect along a line $\ell$ whose direction is perpendicular to both normals $\vec{n}_1$ and $\vec{n}_2$.}
\end{popfigure}

\begin{methode}[Finding the line of intersection of two planes]
Given planes $\vec{r}\cdot\vec{n}_1=d_1$ and $\vec{r}\cdot\vec{n}_2=d_2$ with non-parallel normals:
\begin{enumerate}
\item \textbf{Direction:} find a vector $\vec{b}$ with $\vec{b}\cdot\vec{n}_1=0$ and $\vec{b}\cdot\vec{n}_2=0$ (two equations in the three components of $\vec{b}$; choose one component freely, or use the cross product $\vec{b}=\vec{n}_1\times\vec{n}_2$ if you know it).
\item \textbf{Point:} find one point on the line by setting one coordinate to a convenient value (often $z=0$) in the two Cartesian equations and solving the resulting $2\times 2$ system.
\item \textbf{Equation:} write $\vec{r}=\vec{a}+\lambda\vec{b}$, and convert to Cartesian form if required.
\end{enumerate}
\emph{Alternative (Method 1):} solve the two Cartesian equations simultaneously, letting one variable equal $\lambda$ and expressing the other two in terms of $\lambda$.
\end{methode}

\begin{exoresolu}[Line of intersection by simultaneous equations]
Find, in vector form, the line of intersection of the planes
\[
x+y+z=6 \qquad\text{and}\qquad 2x-y+z=3.
\]
\tcblower
\textbf{Solution (Method 1).} Let $z=\lambda$. The equations become
\[
x+y=6-\lambda,\qquad 2x-y=3-\lambda.
\]
Adding: $3x=9-2\lambda$, so $x=3-\dfrac{2\lambda}{3}$. Then $y=6-\lambda-x=3-\dfrac{\lambda}{3}$. Hence
\[
\vec{r}=\begin{pmatrix}x\\y\\z\end{pmatrix}
=\begin{pmatrix}3\\3\\0\end{pmatrix}+\lambda\begin{pmatrix}-2/3\\-1/3\\1\end{pmatrix}
=\begin{pmatrix}3\\3\\0\end{pmatrix}+\mu\begin{pmatrix}2\\1\\-3\end{pmatrix},
\]
where we rescaled the parameter ($\mu=-\lambda/3$) to clear fractions.
\textbf{Check:} the point $(3,3,0)$: $3+3+0=6$ \checkmark\ and $6-3+0=3$ \checkmark. Direction $(2,1,-3)$: $(2,1,-3)\cdot(1,1,1)=0$ \checkmark\ and $(2,1,-3)\cdot(2,-1,1)=4-1-3=0$ \checkmark.
\end{exoresolu}

\begin{exoresolu}[Line of intersection via the two normals]
Find the vector and Cartesian equations of the line of intersection of
\[
\Pi_1:\ \vec{r}\cdot\begin{pmatrix}1\\2\\-1\end{pmatrix}=3,
\qquad
\Pi_2:\ \vec{r}\cdot\begin{pmatrix}2\\1\\1\end{pmatrix}=4.
\]
\tcblower
\textbf{Solution (Method 2).} \emph{Direction.} Let $\vec{b}=\begin{pmatrix}p\\q\\s\end{pmatrix}$. We need
\[
p+2q-s=0,\qquad 2p+q+s=0.
\]
Adding: $3p+3q=0$, so $q=-p$. Then $s=p+2q=-p$. Choosing $p=1$: $\vec{b}=\begin{pmatrix}1\\-1\\-1\end{pmatrix}$.
\emph{Point.} Set $z=0$ in the Cartesian equations $x+2y-z=3$ and $2x+y+z=4$:
\[
x+2y=3,\qquad 2x+y=4.
\]
From the first, $x=3-2y$; substituting: $6-4y+y=4$, so $y=\dfrac{2}{3}$, $x=\dfrac{5}{3}$. Thus $\vec{a}=\begin{pmatrix}5/3\\2/3\\0\end{pmatrix}$.
\emph{Equations.} Vector form:
\[
\vec{r}=\begin{pmatrix}5/3\\2/3\\0\end{pmatrix}+\lambda\begin{pmatrix}1\\-1\\-1\end{pmatrix}.
\]
Cartesian form: $\dfrac{x-\frac53}{1}=\dfrac{y-\frac23}{-1}=\dfrac{z}{-1}$.
\textbf{Check:} $\vec{a}\cdot\vec{n}_1=\tfrac53+\tfrac43-0=3$ \checkmark; $\vec{a}\cdot\vec{n}_2=\tfrac{10}{3}+\tfrac23+0=4$ \checkmark; $\vec{b}\cdot\vec{n}_1=1-2+1=0$ \checkmark; $\vec{b}\cdot\vec{n}_2=2-1-1=0$ \checkmark.
\end{exoresolu}

\begin{attention}[Parallel or coincident planes: no line of intersection]
Before hunting for a line of intersection, compare the normals. If $\vec{n}_2=k\vec{n}_1$, the planes are parallel. Then compare the constants: the planes are \emph{coincident} only if the whole equations are proportional, i.e.\ $d_2=kd_1$ as well. For example, $x+2y-z=3$ and $2x+4y-2z=6$ are the \emph{same} plane (infinitely many common points — the whole plane), whereas $x+2y-z=3$ and $2x+4y-2z=5$ are \emph{distinct parallel planes} (no common point at all). In neither case is there a single line of intersection.
\end{attention}

\begin{methode}[Exam tip: always verify your line of intersection]
After finding $\vec{r}=\vec{a}+\lambda\vec{b}$ as the intersection of two planes, spend thirty seconds checking:
\begin{enumerate}
\item $\vec{a}$ satisfies \textbf{both} plane equations ($\vec{a}\cdot\vec{n}_1=d_1$ and $\vec{a}\cdot\vec{n}_2=d_2$);
\item $\vec{b}$ is perpendicular to \textbf{both} normals ($\vec{b}\cdot\vec{n}_1=0$ and $\vec{b}\cdot\vec{n}_2=0$).
\end{enumerate}
These four quick dot products catch almost every arithmetic slip and cost you nothing in the examination hall.
\end{methode}

\begin{exercice}[Practice]
\begin{enumerate}
\item Find the point of intersection of the line $\vec{r}=\begin{pmatrix}2\\1\\0\end{pmatrix}+\lambda\begin{pmatrix}1\\-1\\2\end{pmatrix}$ with the plane $\vec{r}\cdot\begin{pmatrix}3\\1\\-1\end{pmatrix}=5$.
\item Show that the line $\vec{r}=\begin{pmatrix}0\\1\\2\end{pmatrix}+\lambda\begin{pmatrix}2\\1\\-1\end{pmatrix}$ does not meet the plane $x-y+z=10$, and state the geometrical relationship between them.
\item Find the vector equation of the line of intersection of the planes $x+2y-z=4$ and $3x-y+2z=1$.
\end{enumerate}
\end{exercice}

\begin{corrige}[Exercise 1]
\textbf{1.} $\vec{a}\cdot\vec{n}=6+1-0=7$; $\vec{b}\cdot\vec{n}=3-1-2=0$. Since $\vec{b}\cdot\vec{n}=0$ and $\vec{a}\cdot\vec{n}=7\neq 5$, the line is parallel to the plane and outside it: \textbf{there is no point of intersection}.

\textbf{2.} With $\vec{n}=\begin{pmatrix}1\\-1\\1\end{pmatrix}$: $\vec{b}\cdot\vec{n}=(2)(1)+(1)(-1)+(-1)(1)=2-1-1=0$, so the line is \emph{parallel} to the plane. Moreover $\vec{a}\cdot\vec{n}=0-1+2=1\neq 10$, so the point $(0,1,2)$ of the line is not in the plane. Hence the line is parallel to the plane and lies \textbf{outside} it: the two never meet.

\textbf{3.} Direction: solve $p+2q-s=0$, $3p-q+2s=0$. From the first, $s=p+2q$; substituting, $3p-q+2p+4q=0$, i.e.\ $5p+3q=0$. Take $p=3$, $q=-5$, then $s=3-10=-7$: $\vec{b}=\begin{pmatrix}3\\-5\\-7\end{pmatrix}$. Point: set $z=0$: $x+2y=4$, $3x-y=1$. From the second, $y=3x-1$; then $x+6x-2=4$, $x=\dfrac67$, $y=\dfrac{11}{7}$. Hence
\[
\vec{r}=\begin{pmatrix}6/7\\11/7\\0\end{pmatrix}+\lambda\begin{pmatrix}3\\-5\\-7\end{pmatrix}.
\]
Check: $\tfrac67+\tfrac{22}{7}=4$ \checkmark; $\tfrac{18}{7}-\tfrac{11}{7}=1$ \checkmark; $3-10+7=0$ \checkmark; $9+5-14=0$ \checkmark.
\end{corrige}

\begin{aretenir}[Key points of this section]
\begin{itemize}
\item Line meets plane: substitute $\vec{r}=\vec{a}+\lambda\vec{b}$ into $\vec{r}\cdot\vec{n}=d$ and solve $\lambda=\dfrac{d-\vec{a}\cdot\vec{n}}{\vec{b}\cdot\vec{n}}$.
\item If $\vec{b}\cdot\vec{n}=0$: the line is parallel to the plane — outside it if $\vec{a}\cdot\vec{n}\neq d$, contained in it if $\vec{a}\cdot\vec{n}=d$.
\item Two planes with non-parallel normals intersect in a \textbf{line}, whose direction $\vec{b}$ satisfies $\vec{b}\cdot\vec{n}_1=\vec{b}\cdot\vec{n}_2=0$.
\item Find the line either by solving the two Cartesian equations with one variable as parameter, or by finding $\vec{b}$ from the normals and one point by fixing a coordinate.
\item Parallel normals mean parallel planes: coincident if the full equations are proportional, otherwise no intersection.
\item Always verify: point in both planes, direction perpendicular to both normals.
\end{itemize}
\end{aretenir}

In the next and final section we reverse the problem: instead of finding the line common to two given planes, we shall construct \emph{families of planes} passing through a given line of intersection or through a given point — a powerful technique that appears regularly in the GCE Advanced Level paper.

\coursec{Families of Planes Through an Intersection and Synthesis}

In the previous section we learned how to find the line of intersection of two non‑parallel planes. That line is common to both planes. A natural question arises: \emph{what are all the planes that contain this line?} This leads to the concept of a \textbf{family of planes} — a powerful technique that frequently appears in the final parts of GCE vector questions, often carrying the last decisive marks. Mastering it will give you a reliable tool to tackle synthesis problems with confidence.

\courssub{The Family of Planes Through the Line of Intersection}

Consider two intersecting planes given in Cartesian form:
\[
\Pi_1:\ a_1x+b_1y+c_1z = d_1 \qquad\text{and}\qquad \Pi_2:\ a_2x+b_2y+c_2z = d_2.
\]
Their line of intersection $L$ consists of all points $(x,y,z)$ that satisfy \emph{both} equations simultaneously. If we form a linear combination
\[
(a_1x+b_1y+c_1z-d_1) + \lambda\,(a_2x+b_2y+c_2z-d_2) = 0 \qquad (\lambda \in \mathbb{R}),
\]
then any point on $L$ makes both brackets zero, hence satisfies the combined equation for \emph{any} value of $\lambda$. Moreover, the combined equation is linear in $x,y,z$ (provided we do not choose $\lambda$ such that all coefficients vanish), so it represents a plane. This plane clearly contains $L$.

\begin{propriete}[Family of Planes Through the Intersection Line]
If $\Pi_1=0$ and $\Pi_2=0$ are the Cartesian equations of two intersecting planes, then for every real number $\lambda$ the equation
\[
\Pi_1 + \lambda\,\Pi_2 = 0
\]
represents a plane that passes through the line of intersection of $\Pi_1$ and $\Pi_2$. Conversely, every plane through that line (except $\Pi_2$ itself) can be obtained for a unique $\lambda$.
\end{propriete}

\begin{attention}[The Missing Plane $\Pi_2$]
The family $\Pi_1 + \lambda\Pi_2 = 0$ does \textbf{not} include the plane $\Pi_2=0$ itself, because that would require $\lambda\to\infty$. When an extra condition is given, you must \textbf{always check separately} whether $\Pi_2$ itself satisfies that condition. If it does, $\Pi_2$ is also a valid answer.
\end{attention}

\begin{popfigure}
\begin{tikzpicture}[scale=1.2, >=Stealth]
% Axes
\draw[->] (-2,0,0) -- (3,0,0) node[right] {$x$};
\draw[->] (0,-2,0) -- (0,3,0) node[above] {$y$};
\draw[->] (0,0,-1) -- (0,0,3) node[above] {$z$};
% Two fixed planes: x+y=2 (vertical) and y+z=2 (slanted)
% Plane Pi1: x+y=2 -> intercepts (2,0,0), (0,2,0), parallel to z-axis
\fill[popblueL, opacity=0.3] (2,0,0) -- (0,2,0) -- (0,2,3) -- (2,0,3) -- cycle;
\draw[popblue, thick] (2,0,0) -- (0,2,0) -- (0,2,3) -- (2,0,3) -- cycle;
\node[popblue] at (2.5,1,1.5) {$\Pi_1$};
% Plane Pi2: y+z=2 -> intercepts (0,2,0), (0,0,2), parallel to x-axis
\fill[poporangeL, opacity=0.3] (0,2,0) -- (0,0,2) -- (3,0,2) -- (3,2,0) -- cycle;
\draw[poporange, thick] (0,2,0) -- (0,0,2) -- (3,0,2) -- (3,2,0) -- cycle;
\node[poporange] at (1.5,2.5,1) {$\Pi_2$};
% Intersection line L: x+y=2, y+z=2 -> param: y=t, x=2-t, z=2-t
\draw[popdark, very thick, dashed] (2,0,0) -- (0,2,0) -- (0,0,2);
% Actually line from (2,0,0) to (0,2,0) to (0,0,2)? Wait: intersection of x+y=2 and y+z=2.
% Solve: x=2-y, z=2-y. So line passes through (2,0,0) when y=0, and (0,2,0) when y=2, and (0,0,2) when y=2? No, (0,0,2) gives x=0,y=0 -> x+y=0 not 2. So line is (2,0,0) to (0,2,0) to (-2,4, -2)? Let's just draw segment between (2,0,0) and (0,2,0) and extend.
\draw[popdark, very thick] (2,0,0) -- (0,2,0);
\draw[popdark, very thick, dashed] (0,2,0) -- (-1,3, -1); % extension
\node[popdark] at (1,1,0.3) {$L$};
% Family members: lambda = -0.5, 0.5, 2
% lambda = -0.5: (x+y-2) -0.5(y+z-2)=0 -> x+0.5y-0.5z-1=0 -> 2x+y-z=2
% intercepts: x=1,y=0,z=0; x=0,y=2,z=0; x=0,y=0,z=-2
\fill[popgreenL, opacity=0.4] (1,0,0) -- (0,2,0) -- (0,0,-2) -- cycle;
\draw[popgreen, thick] (1,0,0) -- (0,2,0) -- (0,0,-2) -- cycle;
\node[popgreen] at (1.2,1,-0.5) {$\lambda=-0.5$};
% lambda = 0.5: (x+y-2)+0.5(y+z-2)=0 -> x+1.5y+0.5z-3=0 -> 2x+3y+z=6
% intercepts: x=3,y=0,z=0; x=0,y=2,z=0; x=0,y=0,z=6
\fill[poppurpleL, opacity=0.3] (3,0,0) -- (0,2,0) -- (0,0,2.5) -- cycle; % approximate
\draw[poppurple, thick] (3,0,0) -- (0,2,0) -- (0,0,2.5) -- cycle;
\node[poppurple] at (2,1,1.5) {$\lambda=0.5$};
% lambda = 2: (x+y-2)+2(y+z-2)=0 -> x+3y+2z-6=0
% intercepts: x=6,y=0,z=0; x=0,y=2,z=0; x=0,y=0,z=3
\fill[poppinkL, opacity=0.3] (2.5,0,0) -- (0,2,0) -- (0,0,3) -- cycle;
\draw[poppink, thick] (2.5,0,0) -- (0,2,0) -- (0,0,3) -- cycle;
\node[poppink] at (1.5,1,2) {$\lambda=2$};
% Legend
\node[align=left, anchor=north west] at (-1.8,2.8) {
\textbf{Family of planes through $L$}\\
Each coloured plane contains the line $L$ (thick black).\\
$\Pi_1$ (blue) corresponds to $\lambda=0$.\\
$\Pi_2$ (orange) is the limit $\lambda\to\infty$.
};
\end{tikzpicture}
\legende{The family $\Pi_1+\lambda\Pi_2=0$ rotates about the common line $L$.}
\end{popfigure}

\courssub{Determining the Particular Plane — The Parameter $\lambda$}

The family contains infinitely many planes. To single out the one required by the problem, we use an \textbf{extra condition} to determine $\lambda$. The three most common conditions are:

\begin{enumerate}
\item \textbf{The plane passes through a given point $P(x_0,y_0,z_0)$ not on $L$.}  
  Substitute the coordinates of $P$ into $\Pi_1+\lambda\Pi_2=0$ and solve for $\lambda$.
\item \textbf{The plane is perpendicular to a given plane $\Pi_3:\ a_3x+b_3y+c_3z=d_3$.}  
  The normal vector of the family plane is $\mathbf{n}_1+\lambda\mathbf{n}_2$. Perpendicularity means $(\mathbf{n}_1+\lambda\mathbf{n}_2)\cdot\mathbf{n}_3=0$. Solve for $\lambda$.
\item \textbf{The plane is parallel to a given line $L'$ with direction vector $\mathbf{v}$.}  
  The normal of the family plane must be perpendicular to $\mathbf{v}$: $(\mathbf{n}_1+\lambda\mathbf{n}_2)\cdot\mathbf{v}=0$. Solve for $\lambda$.
\end{enumerate}

\begin{methode}[Finding the Required Plane in the Family]
\begin{enumerate}
\item Write the two given planes in the form $\Pi_1=0$, $\Pi_2=0$ (i.e. bring all terms to one side).
\item Form the family equation $\Pi_1+\lambda\Pi_2=0$.
\item Use the extra condition to obtain an equation for $\lambda$.
\item Solve for $\lambda$.
\item \textbf{Check separately whether $\Pi_2=0$ itself satisfies the extra condition.}
\item Substitute the found $\lambda$ back into the family equation, simplify (preferably to integer coefficients), and give the final Cartesian equation.
\end{enumerate}
\end{methode}

\begin{exoresolu}[Plane Through the Intersection Line and a Point]
Find the equation of the plane that passes through the line of intersection of the planes
$\Pi_1:\ 2x-y+3z=5$ and $\Pi_2:\ x+4y-z=2$, and also passes through the point $P(1,-2,3)$.
\tcblower
\textbf{Solution.}
\begin{enumerate}
\item Write in ``$=0$'' form:
$\Pi_1:\ 2x-y+3z-5=0$, $\quad \Pi_2:\ x+4y-z-2=0$.
\item Family: $(2x-y+3z-5)+\lambda(x+4y-z-2)=0$.
\item Substitute $P(1,-2,3)$:
$(2(1)-(-2)+3(3)-5)+\lambda(1+4(-2)-3-2)=0$
$\Rightarrow (2+2+9-5)+\lambda(1-8-3-2)=0$
$\Rightarrow 8 + \lambda(-12)=0 \quad\Rightarrow\quad \lambda=\frac{2}{3}$.
\item Check $\Pi_2$: does $P$ lie on $\Pi_2$? $1+4(-2)-3=1-8-3=-10\neq2$, so $\Pi_2$ is not the answer.
\item Substitute $\lambda=\frac{2}{3}$:
$(2x-y+3z-5)+\frac{2}{3}(x+4y-z-2)=0$
Multiply by $3$: $3(2x-y+3z-5)+2(x+4y-z-2)=0$
$6x-3y+9z-15+2x+8y-2z-4=0$
$8x+5y+7z-19=0$.
\end{enumerate}
\textbf{Answer:} $\boxed{8x+5y+7z=19}$.
\end{exoresolu}

\begin{exoresolu}[Plane Through the Intersection Line, Perpendicular to Another Plane]
Find the equation of the plane through the line of intersection of
$\Pi_1:\ x-2y+z=3$ and $\Pi_2:\ 2x+y-z=1$, which is perpendicular to the plane $\Pi_3:\ 3x-y+2z=4$.
\tcblower
\textbf{Solution.}
\begin{enumerate}
\item $\Pi_1:\ x-2y+z-3=0$, $\quad \Pi_2:\ 2x+y-z-1=0$.
\item Family: $(x-2y+z-3)+\lambda(2x+y-z-1)=0$.
  Normal vector: $\mathbf{n}(\lambda) = (1+2\lambda,\ -2+\lambda,\ 1-\lambda)$.
\item $\Pi_3$ has normal $\mathbf{n}_3=(3,-1,2)$. Perpendicular $\Rightarrow \mathbf{n}(\lambda)\cdot\mathbf{n}_3=0$:
$(1+2\lambda)(3)+(-2+\lambda)(-1)+(1-\lambda)(2)=0$
$3+6\lambda+2-\lambda+2-2\lambda=0$
$7+3\lambda=0 \quad\Rightarrow\quad \lambda=-\frac{7}{3}$.
\item Check $\Pi_2$: its normal is $(2,1,-1)$. Dot with $\mathbf{n}_3$: $6-1-2=3\neq0$, so $\Pi_2$ is not perpendicular to $\Pi_3$.
\item Substitute $\lambda=-\frac{7}{3}$:
$(x-2y+z-3)-\frac{7}{3}(2x+y-z-1)=0$
Multiply by $3$: $3(x-2y+z-3)-7(2x+y-z-1)=0$
$3x-6y+3z-9-14x-7y+7z+7=0$
$-11x-13y+10z-2=0$ or $11x+13y-10z+2=0$.
\end{enumerate}
\textbf{Answer:} $\boxed{11x+13y-10z=-2}$.
\end{exoresolu}

\begin{exoresolu}[Plane Through the Intersection Line, Parallel to a Given Line]
Find the plane through the intersection of $\Pi_1:\ x+y+z=6$ and $\Pi_2:\ 2x-y+3z=5$ that is parallel to the line
$L':\ \frac{x-1}{2}=\frac{y+2}{-1}=\frac{z}{3}$.
\tcblower
\textbf{Solution.}
\begin{enumerate}
\item $\Pi_1:\ x+y+z-6=0$, $\quad \Pi_2:\ 2x-y+3z-5=0$.
\item Family: $(x+y+z-6)+\lambda(2x-y+3z-5)=0$.
  Normal: $\mathbf{n}(\lambda)=(1+2\lambda,\ 1-\lambda,\ 1+3\lambda)$.
\item Direction of $L'$: $\mathbf{v}=(2,-1,3)$. Parallel $\Rightarrow \mathbf{n}(\lambda)\cdot\mathbf{v}=0$:
$(1+2\lambda)(2)+(1-\lambda)(-1)+(1+3\lambda)(3)=0$
$2+4\lambda-1+\lambda+3+9\lambda=0$
$4+14\lambda=0 \quad\Rightarrow\quad \lambda=-\frac{2}{7}$.
\item Check $\Pi_2$: normal $(2,-1,3)$ dot $\mathbf{v}=4+1+9=14\neq0$, so $\Pi_2$ not parallel.
\item Substitute $\lambda=-\frac{2}{7}$:
$(x+y+z-6)-\frac{2}{7}(2x-y+3z-5)=0$
Multiply by $7$: $7(x+y+z-6)-2(2x-y+3z-5)=0$
$7x+7y+7z-42-4x+2y-6z+10=0$
$3x+9y+z-32=0$.
\end{enumerate}
\textbf{Answer:} $\boxed{3x+9y+z=32}$.
\end{exoresolu}

\courssub{Lesson 51 — Planes Through the Point of Intersection of Two Planes}

The syllabus mentions ``Equation of planes through the point of intersection of two planes''. Two planes intersect in a \emph{line}, not a point. However, if a \emph{third} plane is introduced, the three planes may intersect in a unique point. In that context, the family method can be adapted: the equation of any plane through the intersection point of three planes $\Pi_1=0$, $\Pi_2=0$, $\Pi_3=0$ (provided they meet at a single point) is a linear combination $\alpha\Pi_1+\beta\Pi_2+\gamma\Pi_3=0$. But the standard GCE technique for ``planes through the line of intersection of two planes and a given point'' is exactly what we have just done (Lessons 50 and 51 are essentially the same skill). 

\begin{exemplebox}[Plane Through the Line of Intersection and a Given Point (Lesson 50/51)]
Find the plane through the line of intersection of $x+2y-z=4$ and $2x-y+3z=1$ that also passes through $A(3,1,-1)$.
\tcblower
Family: $(x+2y-z-4)+\lambda(2x-y+3z-1)=0$.
Substitute $A$: $(3+2-(-1)-4)+\lambda(6-1-3-1)=0 \Rightarrow 2+\lambda(1)=0 \Rightarrow \lambda=-2$.
Check $\Pi_2$: $2(3)-1+3(-1)=6-1-3=2\neq1$, so $\Pi_2$ not valid.
Equation: $(x+2y-z-4)-2(2x-y+3z-1)=0 \Rightarrow -3x+4y-7z-2=0$ or $3x-4y+7z=-2$.
\end{exemplebox}

\courssub{Full GCE‑Style Synthesis}

A typical Upper Sixth question combines several skills: find the line of intersection of two planes, then find a plane through that line satisfying an extra condition. Let us work through a complete example.

\begin{exoresolu}[Synthesis: Line of Intersection + Family of Planes]
The planes $\Pi_1$ and $\Pi_2$ have equations
\[
\Pi_1:\ 2x-y+z=4,\qquad \Pi_2:\ x+3y-2z=1.
\]
\begin{enumerate}
\item Find a vector equation of the line $L$ of intersection of $\Pi_1$ and $\Pi_2$.
\item Find the Cartesian equation of the plane $\Pi$ that contains $L$ and is perpendicular to the plane $\Pi_3:\ 4x-2y+z=7$.
\end{enumerate}
\tcblower
\textbf{Solution.}
\begin{enumerate}
\item \textbf{Line of intersection.}
  Solve the system:
  \[
  \begin{cases}
  2x-y+z=4 & (1)\\
  x+3y-2z=1 & (2)
  \end{cases}
  \]
  Let $z=t$ (parameter). From (1): $2x-y=4-t$. From (2): $x+3y=1+2t$.
  Solve for $x,y$ in terms of $t$:
  Multiply (2) by 2: $2x+6y=2+4t$. Subtract (1): $7y = (2+4t)-(4-t) = -2+5t \Rightarrow y = \frac{5t-2}{7}$.
  Then $x = 1+2t-3y = 1+2t-\frac{15t-6}{7} = \frac{7+14t-15t+6}{7} = \frac{13-t}{7}$.
  So $L:\ \mathbf{r} = \begin{pmatrix}13/7 \\ -2/7 \\ 0\end{pmatrix} + t\begin{pmatrix}-1/7 \\ 5/7 \\ 1\end{pmatrix}$.
  A neater direction vector is $\mathbf{d}=(-1,5,7)$ (multiply by 7). A point on $L$ is $A\left(\frac{13}{7},-\frac{2}{7},0\right)$.
  Vector equation: $\mathbf{r} = \begin{pmatrix}13/7 \\ -2/7 \\ 0\end{pmatrix} + \mu\begin{pmatrix}-1 \\ 5 \\ 7\end{pmatrix}$.

\item \textbf{Plane through $L$ perpendicular to $\Pi_3$.}
  Family: $(2x-y+z-4)+\lambda(x+3y-2z-1)=0$.
  Normal of family: $\mathbf{n}(\lambda)=(2+\lambda,\ -1+3\lambda,\ 1-2\lambda)$.
  $\Pi_3$ has normal $\mathbf{n}_3=(4,-2,1)$. Perpendicular $\Rightarrow \mathbf{n}(\lambda)\cdot\mathbf{n}_3=0$:
  $(2+\lambda)4 + (-1+3\lambda)(-2) + (1-2\lambda)(1) = 0$
  $8+4\lambda + 2-6\lambda + 1-2\lambda = 0$
  $11 -4\lambda = 0 \quad\Rightarrow\quad \lambda = \frac{11}{4}$.
  Check $\Pi_2$: normal $(1,3,-2)$ dot $\mathbf{n}_3 = 4-6-2=-4\neq0$, so $\Pi_2$ not perpendicular.
  Substitute $\lambda=\frac{11}{4}$:
  $(2x-y+z-4)+\frac{11}{4}(x+3y-2z-1)=0$
  Multiply by 4: $4(2x-y+z-4)+11(x+3y-2z-1)=0$
  $8x-4y+4z-16+11x+33y-22z-11=0$
  $19x+29y-18z-27=0$.
\end{enumerate}
\textbf{Answers:}
\begin{enumerate}
\item $L:\ \mathbf{r} = \begin{pmatrix}13/7 \\ -2/7 \\ 0\end{pmatrix} + \mu\begin{pmatrix}-1 \\ 5 \\ 7\end{pmatrix}$ (or equivalent).
\item $\Pi:\ \boxed{19x+29y-18z=27}$.
\end{enumerate}
\end{exoresolu}

\courssub{Revision Summary — All Formulas of the Chapter}

The table below collects every key formula and result from the whole Vectors chapter (Sections 1–6). Keep it handy for revision.

\begin{center}
\begin{tabularx}{\linewidth}{|l|X|}\hline
\rowcolor{popblueL}
\textbf{Topic} & \textbf{Key Formulas / Results} \\ \hline
\textbf{1. Equations of a Plane} &
Vector: $\mathbf{r}\cdot\mathbf{n}=d$ or $\mathbf{r}=\mathbf{a}+\lambda\mathbf{b}+\mu\mathbf{c}$ \\
& Cartesian: $ax+by+cz=d$ where $\mathbf{n}=(a,b,c)$ \\
& Through point $\mathbf{a}$ with normal $\mathbf{n}$: $(\mathbf{r}-\mathbf{a})\cdot\mathbf{n}=0$ \\
& Through three non‑collinear points $A,B,C$: $(\mathbf{r}-\mathbf{a})\cdot((\mathbf{b}-\mathbf{a})\times(\mathbf{c}-\mathbf{a}))=0$ \\ \hline
\textbf{2. Parametric Form} &
$\mathbf{r}=\mathbf{a}+\lambda\mathbf{b}+\mu\mathbf{c}$ ($\mathbf{b},\mathbf{c}$ parallel to plane, not parallel to each other) \\
& Cartesian from parametric: eliminate $\lambda,\mu$ or use $\mathbf{n}=\mathbf{b}\times\mathbf{c}$ \\ \hline
\textbf{3. Perpendicular Distance} &
Point $P(\mathbf{p})$ to plane $\mathbf{r}\cdot\mathbf{n}=d$: $\displaystyle \frac{|\mathbf{p}\cdot\mathbf{n}-d|}{|\mathbf{n}|}$ \\
& Cartesian: $\displaystyle \frac{|ax_0+by_0+cz_0-d|}{\sqrt{a^2+b^2+c^2}}$ \\ \hline
\textbf{4. Angles} &
Angle between planes $\Pi_1,\Pi_2$: $\cos\theta=\dfrac{|\mathbf{n}_1\cdot\mathbf{n}_2|}{|\mathbf{n}_1||\mathbf{n}_2|}$ \\
& Angle between line $\mathbf{r}=\mathbf{a}+t\mathbf{v}$ and plane $\mathbf{r}\cdot\mathbf{n}=d$: $\sin\phi=\dfrac{|\mathbf{v}\cdot\mathbf{n}|}{|\mathbf{v}||\mathbf{n}|}$ \\ \hline
\textbf{5. Intersections} &
Line $\mathbf{r}=\mathbf{a}+t\mathbf{v}$ and plane $\mathbf{r}\cdot\mathbf{n}=d$: substitute, solve for $t$, find point. \\
& Two planes $\mathbf{r}\cdot\mathbf{n}_1=d_1$, $\mathbf{r}\cdot\mathbf{n}_2=d_2$: direction $\mathbf{d}=\mathbf{n}_1\times\mathbf{n}_2$; find one common point. \\
& Three planes: solve $3\times3$ system; unique point / line / no intersection (prism). \\ \hline
\textbf{6. Family of Planes} &
Through intersection of $\Pi_1=0$ and $\Pi_2=0$: $\Pi_1+\lambda\Pi_2=0$ ($\lambda\in\mathbb{R}$). \\
& \textbf{Extra conditions:} \\
& \quad – Through point $P$: substitute $P$ to find $\lambda$. \\
& \quad – Perpendicular to $\Pi_3$: $(\mathbf{n}_1+\lambda\mathbf{n}_2)\cdot\mathbf{n}_3=0$. \\
& \quad – Parallel to line $\mathbf{v}$: $(\mathbf{n}_1+\lambda\mathbf{n}_2)\cdot\mathbf{v}=0$. \\
& \textbf{Always check $\Pi_2$ separately!} \\ \hline
\end{tabularx}
\end{center}

\begin{popfigure}
\begin{tikzpicture}[node distance=1.2cm and 1.5cm, >=Stealth,
  decision/.style={diamond, draw=popdark, thick, fill=popblueL, text width=3.5cm, align=center, inner sep=2pt},
  process/.style={rectangle, draw=popdark, thick, fill=popgreenL, text width=3.5cm, align=center, inner sep=2pt, rounded corners},
  start/.style={ellipse, draw=popdark, thick, fill=poporangeL, text width=2.5cm, align=center, inner sep=2pt},
  io/.style={trapezium, draw=popdark, thick, fill=popyellowL, text width=3cm, align=center, inner sep=2pt, trapezium left angle=70, trapezium right angle=110}]
\node[start] (start) {Vector Problem Given};
\node[process, below of=start] (identify) {Identify objects: points, lines, planes, given equations};
\node[decision, below of=identify] (type) {What is asked?};
\node[process, below left=of type, align=center] (eqplane){Equation of plane:\\ through point + normal\\ through 3 points\\ parametric $\to$ Cartesian};
\node[process, below right=of type, align=center] (distance){Distance point–plane:\\ $\frac{|\mathbf{p}\cdot\mathbf{n}-d|}{|\mathbf{n}|}$};
\node[process, left=of eqplane, align=center] (angle){Angles:\\ plane–plane: $\cos\theta$\\ line–plane: $\sin\phi$};
\node[process, right=of distance, align=center] (intersect){Intersections:\\ line–plane: sub $\mathbf{r}$\\ two planes: $\mathbf{d}=\mathbf{n}_1\times\mathbf{n}_2$\\ three planes: solve system};
\node[process, below=2cm of type, align=center] (family){Family of planes through\\ intersection line:\\ $\Pi_1+\lambda\Pi_2=0$\\ Find $\lambda$ from extra condition\\ Check $\Pi_2$ separately!};
\draw[->] (start) -- (identify);
\draw[->] (identify) -- (type);
\draw[->] (type) -| node[near start, above] {Equation} (eqplane);
\draw[->] (type) -| node[near start, above] {Distance} (distance);
\draw[->] (type) -| node[near start, above] {Angle} (angle);
\draw[->] (type) -| node[near start, above] {Intersection} (intersect);
\draw[->] (type) -- node[near start, right, align=center]{Plane through\\ intersection line} (family);
\draw[->] (eqplane) -- (family);
\draw[->] (intersect) -- (family);
\end{tikzpicture}
\legende{Flowchart for solving GCE Vector problems — follow the arrows from the question type to the appropriate tool.}
\end{popfigure}

\courssub{Common Mistakes and Exam Tips}

\begin{attention}[Forgetting to Check $\Pi_2$]
The family $\Pi_1+\lambda\Pi_2=0$ never gives $\Pi_2$ itself. If the extra condition is satisfied by $\Pi_2$ (e.g. the given point lies on $\Pi_2$, or $\Pi_2$ is perpendicular to $\Pi_3$), then $\Pi_2$ is a valid answer and you will lose marks if you only give the $\lambda$-plane.
\end{attention}

\begin{attention}[Arithmetic Errors in $\lambda$]
Substituting the point or computing the dot product for perpendicularity/parallelism involves signed numbers. Work carefully, keep fractions exact, and double‑check your algebra before simplifying the final equation.
\end{attention}

\begin{attention}[Not Simplifying to Integer Coefficients]
GCE marking schemes usually expect the final Cartesian equation with integer coefficients (and preferably the constant term positive). After finding $\lambda$, multiply through by the denominator and simplify.
\end{attention}

\begin{aretenir}[Exam Tip — The Family Method is a ``Gift'' Question]
In many GCE papers, the last 4–6 marks of a vector question are awarded for correctly using the family of planes. The steps are routine:
\begin{enumerate}
\item Write $\Pi_1=0$, $\Pi_2=0$.
\item Form $\Pi_1+\lambda\Pi_2=0$.
\item Use the condition to find $\lambda$.
\item Check $\Pi_2$.
\item Substitute and simplify.
\end{enumerate}
Practise this pattern until it becomes automatic.
\end{aretenir}

\begin{savaistu}[Historical Note]
The idea of a ``pencil of planes'' (family through a common line) was studied by projective geometers like Poncelet and Möbius in the 19th century. It is a beautiful example of how a linear combination of equations preserves the common solutions — a principle that underlies much of linear algebra and algebraic geometry.
\end{savaistu}

\begin{exercice}[Consolidation Exercises]
\begin{enumerate}
\item Find the equation of the plane through the intersection of $x-2y+3z=4$ and $2x+y-z=5$ that passes through $(1,1,1)$.
\item Determine the plane through the line of intersection of $3x-y+2z=1$ and $x+4y-z=3$ which is perpendicular to $2x-y+3z=6$.
\item The planes $\Pi_1:\ x+y+z=3$ and $\Pi_2:\ 2x-y+3z=4$ intersect in line $L$. Find the plane containing $L$ that is parallel to the line $\frac{x}{1}=\frac{y}{-2}=\frac{z}{2}$.
\item (Synthesis) Given $\Pi_1:\ 2x-y+z=5$, $\Pi_2:\ x+2y-3z=1$, $\Pi_3:\ 3x-y+2z=4$.
  \begin{enumerate}
  \item Find the line of intersection of $\Pi_1$ and $\Pi_2$.
  \item Find the plane through this line that is perpendicular to $\Pi_3$.
  \end{enumerate}
\item A plane $\Pi$ passes through the intersection of $x+2y-z=2$ and $3x-y+2z=5$. If $\Pi$ also passes through the point $(0,0,0)$, find its equation. (Don't forget to check $\Pi_2$!)
\end{enumerate}
\end{exercice}

\begin{corrige}[Exercise 1]
Family: $(x-2y+3z-4)+\lambda(2x+y-z-5)=0$.
Substitute $(1,1,1)$: $(1-2+3-4)+\lambda(2+1-1-5)= -2+\lambda(-3)=0 \Rightarrow \lambda=-\frac{2}{3}$.
Check $\Pi_2$: $2+1-1=2\neq5$, so $\Pi_2$ not valid.
Equation: $(x-2y+3z-4)-\frac{2}{3}(2x+y-z-5)=0 \Rightarrow 3x-6y+9z-12-4x-2y+2z+10=0 \Rightarrow -x-8y+11z-2=0$ or $x+8y-11z=-2$.
\end{corrige}

\begin{corrige}[Exercise 2]
Family: $(3x-y+2z-1)+\lambda(x+4y-z-3)=0$.
Normal: $(3+\lambda,\ -1+4\lambda,\ 2-\lambda)$.
Perpendicular to $2x-y+3z=6$ $\Rightarrow$ $(3+\lambda)2+(-1+4\lambda)(-1)+(2-\lambda)3=0$
$6+2\lambda+1-4\lambda+6-3\lambda=0 \Rightarrow 13-5\lambda=0 \Rightarrow \lambda=13/5$.
Check $\Pi_2$: normal $(1,4,-1)$ dot $(2,-1,3)=2-4-3=-5\neq0$.
Equation: $(3x-y+2z-1)+\frac{13}{5}(x+4y-z-3)=0 \Rightarrow 5(3x-y+2z-1)+13(x+4y-z-3)=0$
$15x-5y+10z-5+13x+52y-13z-39=0 \Rightarrow 28x+47y-3z-44=0$.
\end{corrige}

\begin{corrige}[Exercise 3]
Family: $(x+y+z-3)+\lambda(2x-y+3z-4)=0$.
Normal: $(1+2\lambda,\ 1-\lambda,\ 1+3\lambda)$.
Parallel to line with direction $(1,-2,2)$ $\Rightarrow$ $(1+2\lambda)(1)+(1-\lambda)(-2)+(1+3\lambda)(2)=0$
$1+2\lambda-2+2\lambda+2+6\lambda=0 \Rightarrow 1+10\lambda=0 \Rightarrow \lambda=-0.1$.
Check $\Pi_2$: normal $(2,-1,3)$ dot $(1,-2,2)=2+2+6=10\neq0$.
Equation: $(x+y+z-3)-0.1(2x-y+3z-4)=0 \Rightarrow 10(x+y+z-3)-(2x-y+3z-4)=0$
$10x+10y+10z-30-2x+y-3z+4=0 \Rightarrow 8x+11y+7z-26=0$.
\end{corrige}

\begin{corrige}[Exercise 4]
(a) Solve $2x-y+z=5$, $x+2y-3z=1$. Let $z=t$: $2x-y=5-t$, $x+2y=1+3t$.
$x=2+?$ Solve: multiply first by 2: $4x-2y=10-2t$. Add second: $5x=11+t \Rightarrow x=(11+t)/5$.
$y=(1+3t-x)/2 = (1+3t-(11+t)/5)/2 = (5+15t-11-t)/10 = (14t-6)/10 = (7t-3)/5$.
Line: $\mathbf{r}=\begin{pmatrix}11/5\\-3/5\\0\end{pmatrix}+t\begin{pmatrix}1/5\\7/5\\1\end{pmatrix}$ or direction $(1,7,5)$.
(b) Family: $(2x-y+z-5)+\lambda(x+2y-3z-1)=0$.
Normal: $(2+\lambda,\ -1+2\lambda,\ 1-3\lambda)$.
Perpendicular to $\Pi_3$ (normal $(3,-1,2)$): $(2+\lambda)3+(-1+2\lambda)(-1)+(1-3\lambda)2=0$
$6+3\lambda+1-2\lambda+2-6\lambda=0 \Rightarrow 9-5\lambda=0 \Rightarrow \lambda=9/5$.
Check $\Pi_2$: normal $(1,2,-3)$ dot $(3,-1,2)=3-2-6=-5\neq0$.
Equation: $(2x-y+z-5)+\frac{9}{5}(x+2y-3z-1)=0 \Rightarrow 5(2x-y+z-5)+9(x+2y-3z-1)=0$
$10x-5y+5z-25+9x+18y-27z-9=0 \Rightarrow 19x+13y-22z-34=0$.
\end{corrige}

\begin{corrige}[Exercise 5]
Family: $(x+2y-z-2)+\lambda(3x-y+2z-5)=0$.
Passes through $(0,0,0)$: $-2+\lambda(-5)=0 \Rightarrow \lambda=-2/5$.
Check $\Pi_2$: $3(0)-0+2(0)=0\neq5$, so $\Pi_2$ not valid.
Equation: $(x+2y-z-2)-\frac{2}{5}(3x-y+2z-5)=0 \Rightarrow 5x+10y-5z-10-6x+2y-4z+10=0 \Rightarrow -x+12y-9z=0$ or $x-12y+9z=0$.
\end{corrige}

This concludes the Vectors chapter. You now have a complete toolkit: equations of planes in all forms, distances, angles, intersections, and the family of planes technique. Regular practice with past GCE papers will cement these skills. Good luck in your preparation!

\newpage

\coursec{Integration activity}

\coursec{Vectors: Integration Activity}

\courssub{Aims of the activity}

This integration activity closes the chapter on vectors and planes by bringing together, in one realistic engineering situation, every major skill you have developed:

\begin{itemize}
\item finding the \cle{Cartesian equation of a plane} through three given points;
\item checking whether a point lies on a plane;
\item computing the \cle{angle between two planes} using their normal vectors;
\item finding the \cle{point of intersection of a line and a plane};
\item computing the \cle{perpendicular distance from a point to a plane} and interpreting it in a practical decision;
\item determining the \cle{line of intersection of two planes} and writing its vector equation.
\end{itemize}

Beyond the calculations, the activity trains you to \emph{model} a real object (a workshop roof) with mathematical objects (planes and lines), to organise group work, to present a solution clearly on a flipchart, and to check each step against the formulas of the chapter — exactly the kind of structured reasoning expected at the GCE Advanced Level.

\courssub{Situation}

A Government Technical College in Buea, on the slopes of Mount Cameroon, is constructing a rectangular workshop for its carpentry section. Because Buea receives heavy rainfall for much of the year, the roof must slope so that rainwater runs off quickly into a gutter.

The architect sets up a three-dimensional coordinate system, with the unit being the \emph{metre}. The concrete floor of the workshop is the horizontal plane $z = 0$. The four top corners of the walls are
\[
A(0,0,3), \qquad B(8,0,3), \qquad C(8,6,4), \qquad D(0,6,4).
\]
Notice that the front wall (through $A$ and $B$) is $3$ m high, while the back wall (through $C$ and $D$) is $4$ m high: the roof therefore rises from front to back, and the gutter will be fixed along the front edge $AB$.

The sloping \cle{roof panel} is the plane $\Pi_1$ passing through the three points $A$, $B$ and $C$.

During construction, the site engineer faces five questions:

\begin{enumerate}
\item What is the Cartesian equation of the roof plane $\Pi_1$? The corner $D$ was designed to lie on the same panel — does it really?
\item What angle does the roof make with the horizontal floor? (The carpenter needs this pitch angle to cut the wooden trusses.)
\item A vertical steel support cable is fixed to the floor at the point $(4,3,0)$ and rises straight up. It is modelled by the line
\[
\vec{r} = \begin{pmatrix} 4 \\ 3 \\ 0 \end{pmatrix} + t \begin{pmatrix} 0 \\ 0 \\ 1 \end{pmatrix}, \qquad t \in \mathbb{R}.
\]
At what point does the cable meet the underside of the roof, and what length of cable must be ordered?
\item A lamp hangs from a hook at the point $P(4,3,5)$. Safety rules require a clearance of at least $1.5$ m between the lamp and the roof surface. Is the rule respected?
\item A second roof panel, on the other side of the building, has equation $2x - 3z + 6 = 0$. The \cle{gutter} is installed along the line where the two panels meet. What is the vector equation of this gutter line?
\end{enumerate}

The class works in groups of four. Each group writes its full solution on a flipchart, then presents it; the other groups check every step against the formulas of the chapter (normal vector, scalar product, distance formula, line--plane intersection).

\courssub{Tasks}

\begin{enumerate}
\item \textbf{Equation of the roof.} Compute the vectors $\overrightarrow{AB}$ and $\overrightarrow{AC}$. Deduce a normal vector $\vec{n} = \overrightarrow{AB} \times \overrightarrow{AC}$ to the plane $\Pi_1$. Hence find the Cartesian equation of $\Pi_1$ in the form $ax + by + cz + d = 0$. Verify that the coordinates of $D$ satisfy this equation.
\item \textbf{Pitch of the roof.} Write down a normal vector to the floor plane $z = 0$. Using the scalar product formula for the angle between two planes, find the acute angle between the roof and the floor, giving your answer to the nearest $0.1^\circ$.
\item \textbf{The support cable.} Substitute the parametric coordinates $(4,\, 3,\, t)$ of a point of the cable into the equation of $\Pi_1$ to find the value of $t$ at the meeting point $M$. State the coordinates of $M$ and deduce the length of cable needed from the floor to the roof.
\item \textbf{The lamp clearance.} Use the perpendicular distance formula
\[
d = \frac{|a x_0 + b y_0 + c z_0 + d|}{\sqrt{a^2 + b^2 + c^2}}
\]
to find the distance from $P(4,3,5)$ to $\Pi_1$, correct to $2$ decimal places. Decide, with justification, whether the $1.5$ m clearance rule is respected.
\item \textbf{The gutter line.} The gutter lies on both $\Pi_1$ and the second panel $\Pi_2 : 2x - 3z + 6 = 0$. Solve the two equations simultaneously (introduce a parameter $\lambda$) to obtain the vector equation of the line of intersection. Give a direction vector of the gutter and the coordinates of one point on it.
\item \textbf{Synthesis.} On your flipchart, draw a labelled sketch of the workshop showing the floor, the roof plane, the cable, the lamp and the gutter, and summarise all five answers in a results table.
\end{enumerate}

\courssub{Model Answer}

\textbf{1. Cartesian equation of the roof plane.}\par

From the coordinates of $A$, $B$ and $C$:
\[
\overrightarrow{AB} = \begin{pmatrix} 8 \\ 0 \\ 0 \end{pmatrix}, \qquad
\overrightarrow{AC} = \begin{pmatrix} 8 \\ 6 \\ 1 \end{pmatrix}.
\]
A normal vector to $\Pi_1$ is
\[
\vec{n} = \overrightarrow{AB} \times \overrightarrow{AC}
= \begin{pmatrix} 8 \\ 0 \\ 0 \end{pmatrix} \times \begin{pmatrix} 8 \\ 6 \\ 1 \end{pmatrix}
= \begin{pmatrix} 0 \cdot 1 - 0 \cdot 6 \\ 0 \cdot 8 - 8 \cdot 1 \\ 8 \cdot 6 - 0 \cdot 8 \end{pmatrix}
= \begin{pmatrix} 0 \\ -8 \\ 48 \end{pmatrix},
\]
which we simplify to $\vec{n} = (0, 1, -6)$ after dividing by $-8$. The plane therefore has equation
\[
y - 6z + d = 0.
\]
Substituting $A(0,0,3)$: $\;0 - 18 + d = 0$, so $d = 18$. Hence
\[
\boxed{\Pi_1 : \; y - 6z + 18 = 0.}
\]
\emph{Check for $D(0,6,4)$:} $6 - 24 + 18 = 0$. \checkmark\; So $D$ lies on the roof plane, as the architect intended.

\textbf{2. Pitch of the roof.}\par

The floor $z = 0$ has normal vector $\vec{n}_0 = (0,0,1)$. The angle $\theta$ between the planes equals the acute angle between their normals:
\[
\cos\theta = \frac{|\vec{n} \cdot \vec{n}_0|}{|\vec{n}|\,|\vec{n}_0|}
= \frac{|{-6}|}{\sqrt{0^2 + 1^2 + (-6)^2}\,\sqrt{1}}
= \frac{6}{\sqrt{37}}.
\]
Therefore
\[
\theta = \cos^{-1}\!\left(\frac{6}{\sqrt{37}}\right) \approx \cos^{-1}(0.9864) \approx 9.6^\circ.
\]
The roof rises at about $9.6^\circ$ to the horizontal — a gentle slope, sufficient for rainwater run-off.

\begin{attention}[Angle between planes]
Always take the \emph{acute} angle (or right angle) between planes. The formula uses the absolute value of the scalar product of the normals. If the computed angle is obtuse ($>90^\circ$), subtract it from $180^\circ$.
\end{attention}

\textbf{3. The support cable.}\par

A point of the cable has coordinates $(4, 3, t)$. Substituting into $y - 6z + 18 = 0$:
\[
3 - 6t + 18 = 0 \quad\Longrightarrow\quad 6t = 21 \quad\Longrightarrow\quad t = \tfrac{7}{2} = 3.5.
\]
The cable meets the roof at
\[
M\left(4,\, 3,\, 3.5\right).
\]
Since the cable starts at the floor point $(4,3,0)$ and rises vertically, the length needed is simply the height of $M$:
\[
\text{length} = 3.5\ \mathrm{m}.
\]
The engineer should order a cable of at least $3.5\ \mathrm{m}$ (plus a small allowance for fastening).

\textbf{4. The lamp clearance.}\par

The distance from $P(4,3,5)$ to $\Pi_1 : y - 6z + 18 = 0$ is
\[
d = \frac{|3 - 6(5) + 18|}{\sqrt{0^2 + 1^2 + (-6)^2}}
= \frac{|3 - 30 + 18|}{\sqrt{37}}
= \frac{9}{\sqrt{37}}
\approx \frac{9}{6.0828} \approx 1.48\ \mathrm{m}.
\]
Since $1.48\ \mathrm{m} < 1.5\ \mathrm{m}$, the clearance rule is \textbf{not respected}: the lamp is about $2$ cm too close to the roof surface. The hook must be lowered slightly (or the lamp repositioned) before the workshop is approved.

\begin{attention}[Interpreting the distance]
The perpendicular distance formula gives the \emph{shortest} distance from the point to the plane. Here the vertical height of $P$ above the floor is $5\ \mathrm{m}$, but the roof above $P$ is at height $z = \tfrac{3+18}{6} = 3.5\ \mathrm{m}$, so $P$ is actually $1.5\ \mathrm{m}$ \emph{vertically} above the roof — yet only $1.48\ \mathrm{m}$ away \emph{perpendicularly}. Always use the perpendicular distance for clearance checks, and note that $P$ lies on the opposite side of the plane from the origin (the sign of $y - 6z + 18$ at $P$ is negative).
\end{attention}

\textbf{5. Vector equation of the gutter line.}\par

The gutter lies on both planes:
\[
\begin{cases}
y - 6z + 18 = 0 & (\Pi_1)\\
2x - 3z + 6 = 0 & (\Pi_2)
\end{cases}
\]
Let $z = \lambda$, where $\lambda \in \mathbb{R}$. Then from $\Pi_1$: $y = 6\lambda - 18$, and from $\Pi_2$: $2x = 3\lambda - 6$, so $x = \tfrac{3}{2}\lambda - 3$. Hence
\[
\vec{r} = \begin{pmatrix} -3 \\ -18 \\ 0 \end{pmatrix} + \lambda \begin{pmatrix} \tfrac{3}{2} \\ 6 \\ 1 \end{pmatrix},
\qquad\text{or, clearing fractions with } \lambda = 2\mu,\quad
\vec{r} = \begin{pmatrix} -3 \\ -18 \\ 0 \end{pmatrix} + \mu \begin{pmatrix} 3 \\ 12 \\ 2 \end{pmatrix},\ \mu \in \mathbb{R}.
\]
\emph{Check of the direction vector:} it must be perpendicular to both normals:
\[
(3,12,2)\cdot(0,1,-6) = 12 - 12 = 0, \qquad (3,12,2)\cdot(2,0,-3) = 6 - 6 = 0. \ \checkmark
\]
So the gutter runs along the line through $(-3,-18,0)$ with direction $(3,12,2)$; within the actual building ($0 \le x \le 8$, $0 \le y \le 6$), it is the segment where the two panels physically meet, joining the points $(1.5,0,3)$ on edge $AB$ and $(3,6,4)$ on edge $CD$.

\textbf{6. Summary and sketch.}\par

\begin{popfigure}
\begin{tikzpicture}[x={(1cm,0cm)}, y={(-0.5cm,0.5cm)}, z={(0cm,0.7cm)}, scale=0.5, >=stealth]
% Coordinates in metres
% Floor corners
\coordinate (A') at (0,0,0);
\coordinate (B') at (8,0,0);
\coordinate (C') at (8,6,0);
\coordinate (D') at (0,6,0);
% Roof corners
\coordinate (A) at (0,0,3);
\coordinate (B) at (8,0,3);
\coordinate (C) at (8,6,4);
\coordinate (D) at (0,6,4);
% Cable and lamp
\coordinate (CableBase) at (4,3,0);
\coordinate (M) at (4,3,3.5);
\coordinate (P) at (4,3,5);
% Gutter segment endpoints
\coordinate (G1) at (1.5,0,3);
\coordinate (G2) at (3,6,4);

% Draw floor
\draw[popmuted, thick] (A') -- (B') -- (C') -- (D') -- cycle;
% Draw walls (vertical edges)
\draw[popmuted] (A') -- (A) (B') -- (B) (C') -- (C) (D') -- (D);
% Draw roof
\draw[popblue, very thick] (A) -- (B) -- (C) -- (D) -- cycle;
% Draw roof diagonals for clarity
\draw[popblue, dashed] (A) -- (C) (B) -- (D);
% Cable
\draw[poporange, very thick] (CableBase) -- (M);
\draw[poporange, dashed] (M) -- (P);
% Lamp
\fill[poporange] (P) circle (2pt) node[above right] {$P(4,3,5)$};
% Gutter
\draw[popgreen, very thick] (G1) -- (G2);
% Labels
\node[below left] at (A') {$A'(0,0,0)$};
\node[below right] at (B') {$B'(8,0,0)$};
\node[above right] at (C') {$C'(8,6,0)$};
\node[above left] at (D') {$D'(0,6,0)$};
\node[below left] at (A) {$A(0,0,3)$};
\node[below right] at (B) {$B(8,0,3)$};
\node[above right] at (C) {$C(8,6,4)$};
\node[above left] at (D) {$D(0,6,4)$};
\node[below] at (CableBase) {$(4,3,0)$};
\node[left] at (M) {$M(4,3,3.5)$};
\node[below left] at (G1) {$G_1$};
\node[above right] at (G2) {$G_2$};
% Axes
\draw[->, thick] (0,0,0) -- (9,0,0) node[below] {$x$};
\draw[->, thick] (0,0,0) -- (0,7,0) node[left] {$y$};
\draw[->, thick] (0,0,0) -- (0,0,6) node[above] {$z$};
\end{tikzpicture}
\legende{Isometric sketch of the workshop: floor (grey), roof (blue), vertical cable (orange), lamp $P$, gutter (green). The gutter segment $G_1G_2$ lies on the intersection of the two roof panels.}
\end{popfigure}

\begin{center}
\begin{tabularx}{\linewidth}{|l|X|}
\hline
\rowcolor{popblueL} \textbf{Quantity} & \textbf{Result} \\
\hline
Roof plane $\Pi_1$ & $y - 6z + 18 = 0$; $D$ lies on it \\
\hline
Pitch angle roof/floor & $\theta \approx 9.6^\circ$ \\
\hline
Cable meets roof at & $M(4,\,3,\,3.5)$; cable length $= 3.5\ \mathrm{m}$ \\
\hline
Lamp clearance & $d = \dfrac{9}{\sqrt{37}} \approx 1.48\ \mathrm{m} < 1.5\ \mathrm{m}$: rule \emph{not} respected \\
\hline
Gutter line & $\vec{r} = (-3,-18,0) + \mu(3,12,2)$, $\mu \in \mathbb{R}$; segment $G_1(1.5,0,3)$ to $G_2(3,6,4)$ \\
\hline
\end{tabularx}
\end{center}

\begin{aretenir}[What this activity consolidates]
One authentic situation has mobilised the whole chapter: the cross product gives a normal vector and hence the Cartesian equation of a plane; the scalar product of normals gives the angle between planes; substitution of a parametric line into the plane equation gives the line--plane intersection; the formula $\frac{|ax_0+by_0+cz_0+d|}{\sqrt{a^2+b^2+c^2}}$ gives the point--plane distance; and solving two plane equations simultaneously gives the line of intersection of two planes. In the examination, present each step exactly as here: formula first, substitution second, conclusion with units and interpretation last.
\end{aretenir}

\newpage

\coursec{Methods and worked examples}

\courssub{Method 1: Finding the Cartesian Equation of a Plane}

A plane in space is completely determined as soon as we know one of its points and the direction \emph{perpendicular} to it. This perpendicular direction is carried by a \cle{normal vector}. Observe first a concrete fact: if $\vec{n}$ is normal to the plane $\Pi$ and $A$ is a fixed point of $\Pi$, then a point $P$ lies on $\Pi$ if and only if $\vec{AP} \perp \vec{n}$, that is, $\vec{AP}\cdot\vec{n} = 0$. Writing $\vec{r} = \vec{OP}$ and $\vec{a} = \vec{OA}$, this becomes $(\vec{r}-\vec{a})\cdot\vec{n}=0$, i.e.
\[ \vec{r}\cdot\vec{n} = \vec{a}\cdot\vec{n}. \]
This is the \cle{vector equation of the plane}. Expanding the scalar products with $\vec{n} = \begin{pmatrix} a \\ b \\ c \end{pmatrix}$ and $\vec{r} = \begin{pmatrix} x \\ y \\ z \end{pmatrix}$ gives the \cle{Cartesian equation} $ax+by+cz = d$ where $d = \vec{a}\cdot\vec{n}$.

\begin{definition}[Normal vector; equations of a plane]
A \cle{normal vector} to a plane $\Pi$ is any non-zero vector perpendicular to every line lying in $\Pi$. If $A(x_0,y_0,z_0)\in\Pi$ and $\vec{n} = \begin{pmatrix} a \\ b \\ c \end{pmatrix}$ is a normal vector, then:
\begin{itemize}
\item vector (scalar product) form: $\vec{r}\cdot\vec{n} = \vec{a}\cdot\vec{n}$;
\item Cartesian form: $ax + by + cz = d$ with $d = ax_0 + by_0 + cz_0$.
\end{itemize}
\end{definition}

\begin{propriete}[Reading information from the Cartesian equation]
For the plane $\Pi:\ ax+by+cz=d$:
\begin{itemize}
\item $\vec{n} = \begin{pmatrix} a \\ b \\ c \end{pmatrix}$ is a normal vector to $\Pi$;
\item two planes are \textbf{parallel} if and only if their normal vectors are parallel (one is a scalar multiple of the other); they are \textbf{coincident} if, in addition, the equations are proportional;
\item two planes are \textbf{perpendicular} if and only if $\vec{n}_1\cdot\vec{n}_2 = 0$.
\end{itemize}
(Proof not examinable; these follow directly from the definition of the scalar product.)
\end{propriete}

\begin{methode}[Cartesian equation of a plane from a point and a normal vector]
To find the Cartesian equation of a plane $\Pi$:
\begin{enumerate}
\item Identify a \cle{normal vector} $\vec{n} = \begin{pmatrix} a \\ b \\ c \end{pmatrix}$ to the plane (given, or obtained as a vector product of two directions parallel to the plane).
\item Write the provisional equation $ax + by + cz = d$.
\item Substitute the coordinates of a known point $A(x_0, y_0, z_0)$ of the plane to find $d = ax_0 + by_0 + cz_0$.
\item Conclude: $\Pi:\ ax + by + cz = d$. Equivalently, the vector form is $\vec{r}\cdot\vec{n} = \vec{a}\cdot\vec{n}$.
\end{enumerate}
\textbf{Reflex:} the coefficients of $x$, $y$, $z$ are always the components of a normal vector. Two planes are parallel if and only if their normal vectors are parallel.
\end{methode}

\begin{exoresolu}[Plane through a point with a given normal]
Find the Cartesian equation of the plane $\Pi$ which passes through the point $A(2, -1, 3)$ and is perpendicular to the vector $\vec{n} = \begin{pmatrix} 1 \\ 4 \\ -2 \end{pmatrix}$.
\tcblower
\textbf{Step 1.} The normal vector is $\vec{n} = \begin{pmatrix} 1 \\ 4 \\ -2 \end{pmatrix}$, so the equation has the form
\[ x + 4y - 2z = d. \]
\textbf{Step 2.} Since $A(2,-1,3)$ lies on $\Pi$:
\[ d = (2) + 4(-1) - 2(3) = 2 - 4 - 6 = -8. \]
\textbf{Step 3.} Hence
\[ \Pi:\ x + 4y - 2z = -8. \]
\textbf{Check:} substituting $A$: $2 + 4(-1) - 2(3) = -8$. \checkmark\ The equation is satisfied.
\end{exoresolu}

\begin{popfigure}
\begin{tikzpicture}[scale=1.0]
% plane as a parallelogram
\draw[fill=popblueL, draw=popblue, thick] (-2.6,-0.9) -- (2.2,-1.5) -- (3.4,0.9) -- (-1.4,1.5) -- cycle;
\node[popblue] at (2.6,0.9) {$\Pi$};
% point A and normal
\fill[popdark] (0.3,0) circle (1.6pt) node[below left, black] {$A$};
\draw[->, very thick, poporange] (0.3,0) -- (0.3,1.9) node[right] {$\vec{n}$};
% a point P in the plane
\fill[popdark] (1.6,-0.55) circle (1.6pt) node[below right, black] {$P$};
\draw[->, thick, popgreen] (0.3,0) -- (1.6,-0.55) node[midway, below, popgreen] {$\vec{AP}$};
% right angle marker
\draw[poporange] (0.3,0.28) -- (0.58,0.28) -- (0.58,0.02);
\end{tikzpicture}
\legende{A plane $\Pi$ through $A$ with normal vector $\vec{n}$: $P\in\Pi \iff \vec{AP}\cdot\vec{n}=0$.}
\end{popfigure}

\courssub{Method 2: Plane Through Three Non-Collinear Points}

Three points \emph{always} lie in some plane; but they determine a \textbf{unique} plane only when they are not on the same straight line. This is why the collinearity test must come first.

\begin{methode}[Cartesian equation of the plane through three points]
Given three \textbf{non-collinear} points $A$, $B$, $C$:
\begin{enumerate}
\item \textbf{First check the points are not collinear:} verify that $\vec{AB}$ and $\vec{AC}$ are not parallel (no scalar $k$ with $\vec{AB} = k\vec{AC}$). If they are parallel, the three points lie on a line and do \emph{not} define a unique plane.
\item Form the two direction vectors $\vec{AB}$ and $\vec{AC}$.
\item Compute the normal vector $\vec{n} = \vec{AB} \times \vec{AC}$ (vector product): it is perpendicular to both, hence to the plane.
\item Write $ax+by+cz = d$ and find $d$ using any one of the three points.
\item \textbf{Verify} that all three points satisfy the final equation.
\end{enumerate}
\textbf{Reflex:} the vector product $\begin{pmatrix} a_1 \\ a_2 \\ a_3 \end{pmatrix} \times \begin{pmatrix} b_1 \\ b_2 \\ b_3 \end{pmatrix} = \begin{pmatrix} a_2b_3 - a_3b_2 \\ a_3b_1 - a_1b_3 \\ a_1b_2 - a_2b_1 \end{pmatrix}$.
\end{methode}

\begin{exoresolu}[Plane through three points]
Find the Cartesian equation of the plane containing the points $A(1, 0, 2)$, $B(2, 1, 1)$ and $C(-1, 2, 1)$.
\tcblower
\textbf{Step 1. Direction vectors.}
\[ \vec{AB} = \begin{pmatrix} 2-1 \\ 1-0 \\ 1-2 \end{pmatrix} = \begin{pmatrix} 1 \\ 1 \\ -1 \end{pmatrix}, \qquad \vec{AC} = \begin{pmatrix} -1-1 \\ 2-0 \\ 1-2 \end{pmatrix} = \begin{pmatrix} -2 \\ 2 \\ -1 \end{pmatrix}. \]
These are not parallel (no $k$ satisfies $1 = -2k$ and $1 = 2k$ simultaneously), so $A, B, C$ are not collinear and define a unique plane.

\textbf{Step 2. Normal vector.}
\[ \vec{n} = \vec{AB} \times \vec{AC} = \begin{pmatrix} (1)(-1) - (-1)(2) \\ (-1)(-2) - (1)(-1) \\ (1)(2) - (1)(-2) \end{pmatrix} = \begin{pmatrix} -1 + 2 \\ 2 + 1 \\ 2 + 2 \end{pmatrix} = \begin{pmatrix} 1 \\ 3 \\ 4 \end{pmatrix}. \]

\textbf{Step 3. Find $d$.} The equation is $x + 3y + 4z = d$. Using $A(1,0,2)$:
\[ d = 1 + 3(0) + 4(2) = 9. \]

\textbf{Step 4. Conclusion and verification.}
\[ \Pi:\ x + 3y + 4z = 9. \]
Check $B(2,1,1)$: $2 + 3 + 4 = 9$ \checkmark. Check $C(-1,2,1)$: $-1 + 6 + 4 = 9$ \checkmark.
\end{exoresolu}

\begin{attention}[Three collinear points]
If the three given points are collinear, $\vec{AB} \times \vec{AC} = \vec{0}$ and the method collapses: infinitely many planes contain a given line. Always test collinearity first. At the GCE, a question may deliberately give collinear points and ask you to comment.
\end{attention}

\begin{exoresolu}[The collinear case]
Show that the points $A(1, 2, 0)$, $B(3, 4, 2)$ and $C(5, 6, 4)$ do not determine a unique plane.
\tcblower
$\vec{AB} = \begin{pmatrix} 2 \\ 2 \\ 2 \end{pmatrix}$ and $\vec{AC} = \begin{pmatrix} 4 \\ 4 \\ 4 \end{pmatrix} = 2\vec{AB}$. The vectors are parallel, so $A$, $B$, $C$ are collinear: they all lie on the line through $A$ with direction $\begin{pmatrix} 1 \\ 1 \\ 1 \end{pmatrix}$. Infinitely many planes contain this line, so no unique plane is determined. (Note that $\vec{AB}\times\vec{AC} = \vec{0}$, confirming the failure of the cross-product method.)
\end{exoresolu}

\courssub{Method 3: Parametric (Vector) Form of the Equation of a Plane}

Just as a line is described by one point and one direction, a plane is described by one point and \textbf{two} independent directions. Moving from $A$ along any combination of $\vec{u}$ and $\vec{v}$ keeps us inside the plane: this is the geometric content of the parametric form.

\begin{definition}[Parametric vector equation of a plane]
The plane through $A$ parallel to the non-parallel vectors $\vec{u}$ and $\vec{v}$ has equation
\[ \vec{r} = \vec{a} + \lambda\vec{u} + \mu\vec{v}, \qquad \lambda, \mu \in \mathbb{R}. \]
Each pair $(\lambda, \mu)$ gives exactly one point of the plane.
\end{definition}

\begin{methode}[Parametric form of a plane; converting between forms]
A plane through $A$ with direction vectors $\vec{u}$ and $\vec{v}$ (non-parallel) has \cle{parametric vector equation}
\[ \vec{r} = \vec{a} + \lambda \vec{u} + \mu \vec{v}, \qquad \lambda, \mu \in \mathbb{R}. \]
\begin{enumerate}
\item \textbf{To obtain the parametric form} from three points $A, B, C$: take $\vec{u} = \vec{AB}$, $\vec{v} = \vec{AC}$.
\item \textbf{To pass from parametric to Cartesian form:} compute $\vec{n} = \vec{u} \times \vec{v}$, then $d = \vec{a}\cdot\vec{n}$, giving $\vec{r}\cdot\vec{n} = d$, i.e. $ax+by+cz=d$.
\item \textbf{To pass from Cartesian to parametric form:} find one point of the plane (set two coordinates to convenient values) and two independent vectors perpendicular to $\vec{n}$.
\end{enumerate}
\textbf{Reflex:} a point $P$ lies on the plane if and only if there exist $\lambda, \mu$ giving its position vector; solve the resulting system of three equations in two unknowns and check consistency.
\end{methode}

\begin{exoresolu}[Parametric form and conversion]
The plane $\Pi$ contains $A(1, 1, 0)$, $B(0, 2, 1)$ and $C(2, 0, 3)$.
\begin{enumerate}
\item[(a)] Write a vector equation of $\Pi$ in parametric form.
\item[(b)] Deduce the Cartesian equation of $\Pi$.
\item[(c)] Determine whether the point $P(3, -1, 5)$ lies on $\Pi$.
\end{enumerate}
\tcblower
\textbf{(a)} $\vec{AB} = \begin{pmatrix} -1 \\ 1 \\ 1 \end{pmatrix}$, $\vec{AC} = \begin{pmatrix} 1 \\ -1 \\ 3 \end{pmatrix}$. Hence
\[ \vec{r} = \begin{pmatrix} 1 \\ 1 \\ 0 \end{pmatrix} + \lambda \begin{pmatrix} -1 \\ 1 \\ 1 \end{pmatrix} + \mu \begin{pmatrix} 1 \\ -1 \\ 3 \end{pmatrix}, \qquad \lambda, \mu \in \mathbb{R}. \]

\textbf{(b)} Normal vector:
\[ \vec{n} = \vec{AB} \times \vec{AC} = \begin{pmatrix} (1)(3) - (1)(-1) \\ (1)(1) - (-1)(3) \\ (-1)(-1) - (1)(1) \end{pmatrix} = \begin{pmatrix} 4 \\ 4 \\ 0 \end{pmatrix} = 4\begin{pmatrix} 1 \\ 1 \\ 0 \end{pmatrix}. \]
Take $\vec{n} = \begin{pmatrix} 1 \\ 1 \\ 0 \end{pmatrix}$. Then $d = 1(1) + 1(1) + 0(0) = 2$, so
\[ \Pi:\ x + y = 2. \]

\textbf{(c)} \emph{Using the parametric form:} we need
\[ 1 - \lambda + \mu = 3, \qquad 1 + \lambda - \mu = -1, \qquad \lambda + 3\mu = 5. \]
The first equation gives $\mu - \lambda = 2$; the second gives $\lambda - \mu = -2$, i.e. the same condition. From the third, $\lambda = 5 - 3\mu$. Substituting: $\mu - (5 - 3\mu) = 2 \Rightarrow 4\mu = 7 \Rightarrow \mu = \tfrac{7}{4}$, $\lambda = 5 - \tfrac{21}{4} = -\tfrac{1}{4}$. The system is consistent, so $P \in \Pi$.

\emph{Quicker check with (b):} $3 + (-1) = 2$ \checkmark, confirming $P \in \Pi$.
\end{exoresolu}

\begin{attention}[Two parameters, not one]
A plane has \textbf{two} parameters ($\lambda$ and $\mu$); a line has only \textbf{one}. Writing $\vec{r} = \vec{a} + \lambda\vec{u}$ for a plane is a classic error: that is the equation of a \emph{line}. Also, the two direction vectors must not be parallel, otherwise they span only a line.
\end{attention}

\courssub{Method 4: Perpendicular Distance from a Point to a Plane}

The \emph{shortest} distance from a point $P$ to a plane $\Pi$ is measured along the perpendicular from $P$ to $\Pi$. If $H$ is the foot of this perpendicular and $A$ is any point of $\Pi$, then $\vec{PH}$ is parallel to $\vec{n}$, and projecting $\vec{AP}$ onto the unit normal $\hat{\vec{n}}$ gives $|PH| = |\vec{AP}\cdot\hat{\vec{n}}|$. Expanding this with $\vec{n} = \begin{pmatrix} a \\ b \\ c \end{pmatrix}$ and $\Pi:\ ax+by+cz=d$ yields the formula below.

\begin{propriete}[Distance from a point to a plane]
The perpendicular distance from $P(x_1,y_1,z_1)$ to the plane $\Pi:\ ax+by+cz=d$ is
\[ \delta = \frac{|ax_1 + by_1 + cz_1 - d|}{\sqrt{a^2+b^2+c^2}}. \]
(The derivation sketched above — projection onto the unit normal — is a useful justification but only the formula and its use are required at the GCE.)
\end{propriete}

\begin{methode}[Distance from a point to a plane]
For the plane $\Pi:\ ax + by + cz = d$ and the point $P(x_1, y_1, z_1)$:
\begin{enumerate}
\item Identify $a, b, c, d$ (write the equation in the form $ax+by+cz-d=0$ if necessary).
\item Apply the formula
\[ \delta = \frac{|ax_1 + by_1 + cz_1 - d|}{\sqrt{a^2 + b^2 + c^2}}. \]
\item If the \textbf{foot of the perpendicular} $H$ is required: parametrise the line through $P$ in direction $\vec{n}$: $\vec{r} = \vec{p} + t\vec{n}$, substitute into the plane equation to find $t$, then $H$ corresponds to that $t$.
\end{enumerate}
\textbf{Reflex:} the numerator is the left-hand side of the plane equation evaluated at $P$, in absolute value; the denominator is the modulus of the normal vector.
\end{methode}

\begin{exoresolu}[Distance and foot of the perpendicular]
Find the perpendicular distance from the point $P(1, 2, 3)$ to the plane $\Pi:\ 2x - y + 2z = 5$, and determine the coordinates of the foot of the perpendicular from $P$ to $\Pi$.
\tcblower
\textbf{Step 1. Distance.} Here $a=2$, $b=-1$, $c=2$, $d=5$:
\[ \delta = \frac{|2(1) - (2) + 2(3) - 5|}{\sqrt{2^2 + (-1)^2 + 2^2}} = \frac{|2 - 2 + 6 - 5|}{\sqrt{9}} = \frac{1}{3}. \]
So the perpendicular distance is $\dfrac{1}{3}$ units.

\textbf{Step 2. Foot of the perpendicular.} The line through $P$ perpendicular to $\Pi$ has direction $\vec{n} = \begin{pmatrix} 2 \\ -1 \\ 2 \end{pmatrix}$:
\[ \vec{r} = \begin{pmatrix} 1 \\ 2 \\ 3 \end{pmatrix} + t \begin{pmatrix} 2 \\ -1 \\ 2 \end{pmatrix}, \quad \text{i.e. } x = 1+2t,\ y = 2-t,\ z = 3+2t. \]
Substitute into $2x - y + 2z = 5$:
\[ 2(1+2t) - (2-t) + 2(3+2t) = 5 \]
\[ 2 + 4t - 2 + t + 6 + 4t = 5 \implies 9t + 6 = 5 \implies t = -\tfrac{1}{9}. \]
Hence
\[ H = \left(1 - \tfrac{2}{9},\ 2 + \tfrac{1}{9},\ 3 - \tfrac{2}{9}\right) = \left(\tfrac{7}{9},\ \tfrac{19}{9},\ \tfrac{25}{9}\right). \]
\textbf{Check:} $2\left(\tfrac{7}{9}\right) - \tfrac{19}{9} + 2\left(\tfrac{25}{9}\right) = \tfrac{14 - 19 + 50}{9} = \tfrac{45}{9} = 5$ \checkmark. Also $|PH| = |t|\,|\vec{n}| = \tfrac{1}{9}\times 3 = \tfrac{1}{3}$, consistent with Step 1.
\end{exoresolu}

\begin{attention}[Signs and the absolute value]
\begin{itemize}
\item The formula contains $-\,d$: for $\Pi:\ 2x-y+2z=5$ you must compute $2x_1 - y_1 + 2z_1 \mathbf{- 5}$, not $+\,5$.
\item Do not forget the absolute value: a distance is never negative.
\item If $\delta = 0$, the point lies \emph{on} the plane.
\end{itemize}
\end{attention}

\courssub{Method 5: Angles Between Planes, and Between a Line and a Plane}

The angle between two planes is defined as the angle between two lines, one in each plane, both perpendicular to the line of intersection. A moment's thought (or a sketch) shows this equals the angle between the two \textbf{normal vectors} — or its supplement; taking the absolute value of the scalar product selects the acute angle. For a line and a plane, the angle is measured between the line and its \emph{projection} on the plane; this is the \textbf{complement} of the angle between the line and the normal, which is why sine replaces cosine.

\begin{propriete}[Angle formulas]
\begin{itemize}
\item Angle between two planes with normals $\vec{n}_1, \vec{n}_2$:
\[ \cos\theta = \frac{|\vec{n}_1\cdot\vec{n}_2|}{|\vec{n}_1|\,|\vec{n}_2|}, \qquad 0 \le \theta \le 90^{\circ}. \]
\item Angle between a line (direction $\vec{d}$) and a plane (normal $\vec{n}$):
\[ \sin\varphi = \frac{|\vec{d}\cdot\vec{n}|}{|\vec{d}|\,|\vec{n}|}, \qquad 0 \le \varphi \le 90^{\circ}. \]
\end{itemize}
(Proofs not examinable; the formulas must be known and used confidently.)
\end{propriete}

\begin{methode}[Angles involving planes and lines]
\begin{enumerate}
\item \textbf{Angle between two planes} with normals $\vec{n}_1$, $\vec{n}_2$: this equals the angle between the normals:
\[ \cos\theta = \frac{|\vec{n}_1 \cdot \vec{n}_2|}{|\vec{n}_1|\,|\vec{n}_2|}. \]
The modulus ensures the acute angle between the planes.
\item \textbf{Angle $\varphi$ between a line} (direction $\vec{d}$) \textbf{and a plane} (normal $\vec{n}$): the line makes angle $\theta$ with the normal where $\cos\theta = \dfrac{|\vec{d}\cdot\vec{n}|}{|\vec{d}||\vec{n}|}$, and the required angle is the complement:
\[ \sin\varphi = \frac{|\vec{d} \cdot \vec{n}|}{|\vec{d}|\,|\vec{n}|}, \qquad \varphi = 90^{\circ} - \theta. \]
\end{enumerate}
\textbf{Reflex:} plane--plane uses \textbf{cosine} of the normals; line--plane uses \textbf{sine} with the normal. A line is parallel to a plane iff $\vec{d}\cdot\vec{n} = 0$; perpendicular iff $\vec{d}$ is parallel to $\vec{n}$.
\end{methode}

\begin{exoresolu}[Two angle computations]
Consider the planes $\Pi_1:\ x + 2y - 2z = 3$ and $\Pi_2:\ 2x - y + z = 1$, and the line $L:\ \vec{r} = \begin{pmatrix} 1 \\ 0 \\ 2 \end{pmatrix} + t\begin{pmatrix} 2 \\ 1 \\ 2 \end{pmatrix}$.
\begin{enumerate}
\item[(a)] Find the acute angle between $\Pi_1$ and $\Pi_2$.
\item[(b)] Find the angle between $L$ and $\Pi_1$.
\end{enumerate}
\tcblower
\textbf{(a)} The normals are $\vec{n}_1 = \begin{pmatrix} 1 \\ 2 \\ -2 \end{pmatrix}$ and $\vec{n}_2 = \begin{pmatrix} 2 \\ -1 \\ 1 \end{pmatrix}$.
\[ \vec{n}_1 \cdot \vec{n}_2 = (1)(2) + (2)(-1) + (-2)(1) = 2 - 2 - 2 = -2. \]
\[ |\vec{n}_1| = \sqrt{1+4+4} = 3, \qquad |\vec{n}_2| = \sqrt{4+1+1} = \sqrt{6}. \]
\[ \cos\theta = \frac{|-2|}{3\sqrt{6}} = \frac{2}{3\sqrt{6}} = \frac{2\sqrt{6}}{18} = \frac{\sqrt{6}}{9} \approx 0.2722. \]
\[ \theta = \cos^{-1}(0.2722) \approx 74.2^{\circ}. \]

\textbf{(b)} Direction of $L$: $\vec{d} = \begin{pmatrix} 2 \\ 1 \\ 2 \end{pmatrix}$, $|\vec{d}| = \sqrt{4+1+4} = 3$.
\[ \vec{d}\cdot\vec{n}_1 = (2)(1) + (1)(2) + (2)(-2) = 2 + 2 - 4 = 0. \]
Since $\vec{d}\cdot\vec{n}_1 = 0$, the line is \textbf{parallel to the plane} $\Pi_1$: the angle between $L$ and $\Pi_1$ is $0^{\circ}$. (Indeed $\sin\varphi = 0$.)
\end{exoresolu}

\courssub{Method 6: Intersection of a Line and a Plane; Intersection of Two Planes}

\begin{methode}[Line--plane and plane--plane intersections]
\textbf{Line and plane.}
\begin{enumerate}
\item Write the line in parametric form: $x = x_0 + \lambda a$, $y = y_0 + \lambda b$, $z = z_0 + \lambda c$.
\item Substitute into the plane equation $ax+by+cz=d$ and solve the linear equation in $\lambda$.
\item One solution $\Rightarrow$ one intersection point (substitute back). No solution $\Rightarrow$ line parallel to plane, not contained in it. Infinitely many $\Rightarrow$ line lies in the plane.
\end{enumerate}
\textbf{Two planes.}
\begin{enumerate}
\item Solve the two Cartesian equations simultaneously: eliminate one variable, set one free variable $= \lambda$, express $x, y, z$ in terms of $\lambda$.
\item The result is the line of intersection in parametric form; read off a point and a direction vector.
\item Alternative for the direction: $\vec{d} = \vec{n}_1 \times \vec{n}_2$.
\item If $\vec{n}_1 \times \vec{n}_2 = \vec{0}$, the planes are parallel (distinct or coincident).
\end{enumerate}
\end{methode}

\begin{exoresolu}[Line meets plane; two planes meet in a line]
\begin{enumerate}
\item[(a)] Find the point of intersection of the line $\vec{r} = \begin{pmatrix} 2 \\ 1 \\ 0 \end{pmatrix} + \lambda \begin{pmatrix} 1 \\ -1 \\ 2 \end{pmatrix}$ with the plane $\Pi:\ x + 2y - z = 6$.
\item[(b)] Find, in vector form, the line of intersection of the planes $\Pi_1:\ x + y + z = 4$ and $\Pi_2:\ 2x - y + z = 5$.
\end{enumerate}
\tcblower
\textbf{(a)} Parametrically: $x = 2+\lambda$, $y = 1-\lambda$, $z = 2\lambda$. Substitute into $\Pi$:
\[ (2+\lambda) + 2(1-\lambda) - 2\lambda = 6 \implies 2 + \lambda + 2 - 2\lambda - 2\lambda = 6 \]
\[ 4 - 3\lambda = 6 \implies \lambda = -\tfrac{2}{3}. \]
The intersection point is
\[ \left(2 - \tfrac{2}{3},\ 1 + \tfrac{2}{3},\ -\tfrac{4}{3}\right) = \left(\tfrac{4}{3},\ \tfrac{5}{3},\ -\tfrac{4}{3}\right). \]
\textbf{Check:} $\tfrac{4}{3} + 2\left(\tfrac{5}{3}\right) - \left(-\tfrac{4}{3}\right) = \tfrac{4+10+4}{3} = \tfrac{18}{3} = 6$ \checkmark.

\textbf{(b)} Subtract the equations: $(2x - y + z) - (x + y + z) = 5 - 4$, i.e.
\[ x - 2y = 1 \implies x = 1 + 2y. \]
Let $y = \lambda$. Then $x = 1 + 2\lambda$, and from $\Pi_1$: $z = 4 - x - y = 4 - (1+2\lambda) - \lambda = 3 - 3\lambda$. Hence
\[ \vec{r} = \begin{pmatrix} 1 \\ 0 \\ 3 \end{pmatrix} + \lambda \begin{pmatrix} 2 \\ 1 \\ -3 \end{pmatrix}, \qquad \lambda \in \mathbb{R}. \]
\textbf{Check of the direction:} $\vec{n}_1 \times \vec{n}_2 = \begin{pmatrix} 1 \\ 1 \\ 1 \end{pmatrix} \times \begin{pmatrix} 2 \\ -1 \\ 1 \end{pmatrix} = \begin{pmatrix} (1)(1)-(1)(-1) \\ (1)(2)-(1)(1) \\ (1)(-1)-(1)(2) \end{pmatrix} = \begin{pmatrix} 2 \\ 1 \\ -3 \end{pmatrix}$ \checkmark. Check the point $(1,0,3)$: in $\Pi_1$: $1+0+3=4$ \checkmark; in $\Pi_2$: $2-0+3=5$ \checkmark.
\end{exoresolu}

\begin{attention}[The three possible positions of a line and a plane]
After substitution you obtain a linear equation in $\lambda$:
\begin{itemize}
\item $k\lambda = m$ with $k \neq 0$: unique intersection point;
\item $0\cdot\lambda = m$ with $m \neq 0$ (e.g. $0 = 5$): impossible — the line is parallel to the plane and does not touch it;
\item $0\cdot\lambda = 0$: true for every $\lambda$ — the line lies entirely in the plane.
\end{itemize}
Always interpret the algebraic outcome geometrically; the GCE regularly asks for this comment.
\end{attention}

\courssub{Method 7: Family of Planes Through the Line of Intersection of Two Planes}

Why does the combination $\Pi_1 + \lambda\Pi_2 = 0$ work? If a point lies on the line of intersection $\ell$, it satisfies \emph{both} equations, so both brackets vanish and the combination vanishes for \emph{every} $\lambda$: hence every plane of the family contains $\ell$. Conversely, the combination is linear in $x, y, z$, so it always represents a plane; varying $\lambda$ tilts the plane around the fixed hinge $\ell$.

\begin{propriete}[Family of planes through a line]
If $\Pi_1:\ a_1x+b_1y+c_1z = d_1$ and $\Pi_2:\ a_2x+b_2y+c_2z = d_2$ intersect in a line $\ell$, then every plane through $\ell$ (except $\Pi_2$ itself) has an equation of the form
\[ (a_1x+b_1y+c_1z - d_1) + \lambda(a_2x+b_2y+c_2z - d_2) = 0, \qquad \lambda \in \mathbb{R}. \]
(Proof not examinable; the justification above suffices.)
\end{propriete}

\begin{methode}[Planes through the intersection of two planes]
If $\Pi_1:\ a_1x+b_1y+c_1z = d_1$ and $\Pi_2:\ a_2x+b_2y+c_2z = d_2$ intersect in a line $\ell$, then \textbf{every plane through $\ell$} (except $\Pi_2$ itself) has equation
\[ (a_1x+b_1y+c_1z - d_1) + \lambda(a_2x+b_2y+c_2z - d_2) = 0, \qquad \lambda \in \mathbb{R}. \]
\begin{enumerate}
\item Write the family $\Pi_1 + \lambda \Pi_2 = 0$.
\item Use the extra condition (passing through a given point, perpendicular to a given plane, at a given distance from a point) to form an equation in $\lambda$.
\item Solve for $\lambda$ and substitute back; simplify.
\item If the required plane turns out to be $\Pi_2$ itself, check it separately.
\end{enumerate}
\textbf{Reflex:} any point of $\ell$ automatically satisfies the family equation for every $\lambda$ — that is why the family consists exactly of the planes through $\ell$.
\end{methode}

\begin{exoresolu}[Plane through a line of intersection and an extra point]
Find the equation of the plane which passes through the line of intersection of the planes
\[ \Pi_1:\ x + y + z = 4 \qquad \text{and} \qquad \Pi_2:\ 2x - y + z = 5, \]
and also passes through the point $Q(1, 2, 3)$.
\tcblower
\textbf{Step 1. The family.} Every plane through the line of intersection has equation
\[ (x + y + z - 4) + \lambda(2x - y + z - 5) = 0. \]

\textbf{Step 2. Impose passage through $Q(1,2,3)$.}
\[ (1 + 2 + 3 - 4) + \lambda(2 - 2 + 3 - 5) = 0 \]
\[ 2 + \lambda(-2) = 0 \implies \lambda = 1. \]

\textbf{Step 3. Substitute back.}
\[ (x + y + z - 4) + (2x - y + z - 5) = 0 \]
\[ 3x + 2z - 9 = 0, \qquad \text{i.e. } \Pi:\ 3x + 2z = 9. \]

\textbf{Step 4. Verification.}
\begin{itemize}
\item $Q(1,2,3)$: $3(1) + 2(3) = 9$ \checkmark.
\item A point of the line of intersection, e.g. $(1, 0, 3)$ found earlier: $3(1) + 2(3) = 9$ \checkmark. Another point of the line, with $\lambda = 1$: $(3, 1, 0)$: $3(3) + 2(0) = 9$ \checkmark.
\end{itemize}
Hence the required plane is $3x + 2z = 9$.
\end{exoresolu}

\begin{exoresolu}[Plane of the family perpendicular to a given plane]
Find the equation of the plane through the line of intersection of $\Pi_1:\ x+y+z=4$ and $\Pi_2:\ 2x-y+z=5$ which is perpendicular to the plane $\Pi_3:\ x - y + z = 0$.
\tcblower
\textbf{Step 1.} The family is $(x+y+z-4) + \lambda(2x-y+z-5) = 0$, i.e.
\[ (1+2\lambda)x + (1-\lambda)y + (1+\lambda)z - (4+5\lambda) = 0. \]
A normal vector of a plane of the family is $\vec{n}_\lambda = \begin{pmatrix} 1+2\lambda \\ 1-\lambda \\ 1+\lambda \end{pmatrix}$.

\textbf{Step 2.} Perpendicularity with $\Pi_3$ (normal $\vec{n}_3 = \begin{pmatrix} 1 \\ -1 \\ 1 \end{pmatrix}$) requires $\vec{n}_\lambda\cdot\vec{n}_3 = 0$:
\[ (1+2\lambda)(1) + (1-\lambda)(-1) + (1+\lambda)(1) = 0 \]
\[ 1 + 2\lambda - 1 + \lambda + 1 + \lambda = 0 \implies 4\lambda + 1 = 0 \implies \lambda = -\tfrac{1}{4}. \]

\textbf{Step 3.} Substitute $\lambda = -\tfrac14$:
\[ \left(1 - \tfrac12\right)x + \left(1 + \tfrac14\right)y + \left(1 - \tfrac14\right)z - \left(4 - \tfrac54\right) = 0 \]
\[ \tfrac12 x + \tfrac54 y + \tfrac34 z - \tfrac{11}{4} = 0. \]
Multiplying by $4$:
\[ \Pi:\ 2x + 5y + 3z = 11. \]
\textbf{Check:} $\vec{n}\cdot\vec{n}_3 = 2 - 5 + 3 = 0$ \checkmark; the point $(1,0,3)$ of the line of intersection: $2 + 0 + 9 = 11$ \checkmark.
\end{exoresolu}

\courssub{Method 8: Synthesis — a Complete GCE-Style Problem}

\begin{methode}[Tackling a multi-part vector geometry question]
\begin{enumerate}
\item Read the whole question first: later parts often reuse earlier results (normal vectors, intersection lines).
\item Keep exact values (fractions, surds) throughout; round only at the final answer.
\item After each computation, verify by substitution — examiners award marks for correct method even after a slip, but a check can save you.
\item Structure your answer with clear part labels (a), (b), (c)\dots\ as in the question.
\end{enumerate}
\end{methode}

\begin{exoresolu}[GCE-style synthesis]
The points $A(1, 1, 1)$, $B(2, 0, 3)$ and $C(0, 2, 2)$ define a plane $\Pi$, and the line $L$ has equation $\vec{r} = \begin{pmatrix} 3 \\ 1 \\ 0 \end{pmatrix} + t\begin{pmatrix} 1 \\ 1 \\ -1 \end{pmatrix}$.
\begin{enumerate}
\item[(a)] Find the Cartesian equation of $\Pi$.
\item[(b)] Find the position vector of the point $T$ where $L$ meets $\Pi$.
\item[(c)] Find the angle between $L$ and $\Pi$.
\item[(d)] Find the perpendicular distance from $D(4, 0, 1)$ to $\Pi$.
\end{enumerate}
\tcblower
\textbf{(a)} $\vec{AB} = \begin{pmatrix} 1 \\ -1 \\ 2 \end{pmatrix}$, $\vec{AC} = \begin{pmatrix} -1 \\ 1 \\ 1 \end{pmatrix}$.
\[ \vec{n} = \vec{AB} \times \vec{AC} = \begin{pmatrix} (-1)(1) - (2)(1) \\ (2)(-1) - (1)(1) \\ (1)(1) - (-1)(-1) \end{pmatrix} = \begin{pmatrix} -3 \\ -3 \\ 0 \end{pmatrix} = -3\begin{pmatrix} 1 \\ 1 \\ 0 \end{pmatrix}. \]
Take $\vec{n} = \begin{pmatrix} 1 \\ 1 \\ 0 \end{pmatrix}$; then $d = 1 + 1 + 0 = 2$, so
\[ \Pi:\ x + y = 2. \]
Check $B$: $2 + 0 = 2$ \checkmark; check $C$: $0 + 2 = 2$ \checkmark.

\textbf{(b)} On $L$: $x = 3+t$, $y = 1+t$, $z = -t$. Substitute into $x + y = 2$:
\[ (3+t) + (1+t) = 2 \implies 4 + 2t = 2 \implies t = -1. \]
Hence $T = (3-1,\ 1-1,\ 1) = (2, 0, 1)$, i.e. $\vec{OT} = \begin{pmatrix} 2 \\ 0 \\ 1 \end{pmatrix}$.

\textbf{(c)} With $\vec{d} = \begin{pmatrix} 1 \\ 1 \\ -1 \end{pmatrix}$ and $\vec{n} = \begin{pmatrix} 1 \\ 1 \\ 0 \end{pmatrix}$:
\[ \vec{d}\cdot\vec{n} = 1 + 1 + 0 = 2, \quad |\vec{d}| = \sqrt{3}, \quad |\vec{n}| = \sqrt{2}. \]
\[ \sin\varphi = \frac{|2|}{\sqrt{3}\sqrt{2}} = \frac{2}{\sqrt{6}} = \frac{\sqrt{6}}{3} \approx 0.8165 \implies \varphi \approx 54.7^{\circ}. \]

\textbf{(d)} For $D(4, 0, 1)$:
\[ \delta = \frac{|4 + 0 - 2|}{\sqrt{1^2 + 1^2 + 0^2}} = \frac{2}{\sqrt{2}} = \sqrt{2} \approx 1.41 \text{ units}. \]
\end{exoresolu}

\begin{aretenir}[Summary of key formulas]
\begin{itemize}
\item Plane through $A$ with normal $\vec{n}$: $\vec{r}\cdot\vec{n} = \vec{a}\cdot\vec{n}$, i.e. $ax+by+cz=d$.
\item Plane through $A$ with directions $\vec{u}, \vec{v}$: $\vec{r} = \vec{a} + \lambda\vec{u} + \mu\vec{v}$; normal $\vec{n} = \vec{u}\times\vec{v}$.
\item Distance from $P(x_1,y_1,z_1)$ to $ax+by+cz=d$: $\delta = \dfrac{|ax_1+by_1+cz_1-d|}{\sqrt{a^2+b^2+c^2}}$.
\item Angle between planes: $\cos\theta = \dfrac{|\vec{n}_1\cdot\vec{n}_2|}{|\vec{n}_1||\vec{n}_2|}$; angle line--plane: $\sin\varphi = \dfrac{|\vec{d}\cdot\vec{n}|}{|\vec{d}||\vec{n}|}$.
\item Line of intersection of two planes: solve simultaneously; direction $\vec{n}_1\times\vec{n}_2$.
\item Family through the line of intersection: $\Pi_1 + \lambda\Pi_2 = 0$.
\end{itemize}
\end{aretenir}

\begin{attention}[Final reminders for the examination]
\begin{itemize}
\item Never confuse the roles: for a plane, the \textbf{normal} gives the Cartesian coefficients; for a line, the \textbf{direction vector} gives the parametric coefficients.
\item Angle line--plane: use $\sin\varphi$ with the normal — using $\cos$ gives the angle with the normal instead, a classic error.
\item In the distance formula, do not forget the absolute value and the subtraction of $d$.
\item When using the family $\Pi_1 + \lambda\Pi_2 = 0$, the plane $\Pi_2$ itself is not included; verify separately if needed.
\end{itemize}
\end{attention}

\begin{exercice}[Practice 1]
Find the Cartesian equation of the plane through $A(1, -2, 3)$ with normal vector $\vec{n} = \begin{pmatrix} 2 \\ -1 \\ 3 \end{pmatrix}$.
\end{exercice}

\begin{corrige}[Practice 1]
The equation is $2x - y + 3z = d$ with $d = 2(1) - (-2) + 3(3) = 2 + 2 + 9 = 13$. Hence $\Pi:\ 2x - y + 3z = 13$.
\end{corrige}

\begin{exercice}[Practice 2]
Find the Cartesian equation of the plane through $A(1, 1, 1)$, $B(2, 3, 0)$ and $C(-1, 0, 2)$.
\end{exercice}

\begin{corrige}[Practice 2]
$\vec{AB} = \begin{pmatrix} 1 \\ 2 \\ -1 \end{pmatrix}$, $\vec{AC} = \begin{pmatrix} -2 \\ -1 \\ 1 \end{pmatrix}$ (not parallel, so the points are not collinear).
\[ \vec{n} = \vec{AB}\times\vec{AC} = \begin{pmatrix} (2)(1)-(-1)(-1) \\ (-1)(-2)-(1)(1) \\ (1)(-1)-(2)(-2) \end{pmatrix} = \begin{pmatrix} 1 \\ 1 \\ 3 \end{pmatrix}. \]
Then $d = 1 + 1 + 3 = 5$, so $\Pi:\ x + y + 3z = 5$. Check $B$: $2+3+0=5$ \checkmark; check $C$: $-1+0+6=5$ \checkmark.
\end{corrige}

\begin{exercice}[Practice 3]
Find the perpendicular distance from $P(2, 0, -1)$ to the plane $\Pi:\ x - 2y + 2z = 8$.
\end{exercice}

\begin{corrige}[Practice 3]
$\delta = \dfrac{|(2) - 2(0) + 2(-1) - 8|}{\sqrt{1+4+4}} = \dfrac{|2 - 2 - 8|}{3} = \dfrac{8}{3}$ units.
\end{corrige}

\begin{exercice}[Practice 4]
Find the acute angle between the planes $x + y + z = 1$ and $x - y = 0$.
\end{exercice}

\begin{corrige}[Practice 4]
$\vec{n}_1\cdot\vec{n}_2 = 1 - 1 + 0 = 0$: the normals are perpendicular, so the planes are perpendicular; the angle is $90^{\circ}$.
\end{corrige}

\begin{exercice}[Practice 5]
Find the line of intersection of the planes $x + y - z = 2$ and $2x - y + z = 1$, giving your answer in vector form.
\end{exercice}

\begin{corrige}[Practice 5]
Add the equations: $3x = 3$, so $x = 1$. Then from the first: $z = x + y - 2 = y - 1$. Let $y = \lambda$: $z = \lambda - 1$. Hence
\[ \vec{r} = \begin{pmatrix} 1 \\ 0 \\ -1 \end{pmatrix} + \lambda \begin{pmatrix} 0 \\ 1 \\ 1 \end{pmatrix}. \]
Check direction: $\vec{n}_1\times\vec{n}_2 = \begin{pmatrix} 1 \\ 1 \\ -1 \end{pmatrix}\times\begin{pmatrix} 2 \\ -1 \\ 1 \end{pmatrix} = \begin{pmatrix} 1-1 \\ -2-1 \\ -1-2 \end{pmatrix} = \begin{pmatrix} 0 \\ -3 \\ -3 \end{pmatrix} = -3\begin{pmatrix} 0 \\ 1 \\ 1 \end{pmatrix}$ \checkmark.
\end{corrige}

\begin{exercice}[Practice 6]
Find the equation of the plane through the line of intersection of $x + 2y - z = 3$ and $2x - y + z = 4$ which passes through the origin.
\end{exercice}

\begin{corrige}[Practice 6]
Family: $(x + 2y - z - 3) + \lambda(2x - y + z - 4) = 0$. At $(0,0,0)$: $-3 - 4\lambda = 0$, so $\lambda = -\tfrac34$. Substituting:
\[ (x+2y-z-3) - \tfrac34(2x-y+z-4) = 0 \implies -\tfrac12 x + \tfrac{11}{4}y - \tfrac74 z = 0. \]
Multiplying by $-4$: $2x - 11y + 7z = 0$. Check: the origin satisfies it \checkmark, and it contains the line of intersection by construction.
\end{corrige}

\begin{savaistu}[Why planes matter beyond the classroom]
Vector equations of planes are the working tool of 3D computer graphics: every flat face of an object on a screen is stored as a plane, and the rendering engine computes line--plane intersections and angles between planes thousands of times per second to decide what is visible and how light reflects. The same mathematics is used by civil engineers — for instance when checking the inclination of a roof or the angle between two road surfaces — and by geologists studying the tilt of rock layers beneath the soil of the Cameroon Volcanic Line.
\end{savaistu}

\newpage

\coursec{Exercises}

\courssub{I check my knowledge}

\begin{exercice}[Multiple choice]
Choose the correct answer, justifying your choice briefly.
\begin{enumerate}
\item A Cartesian equation of the plane through $A(1,-2,3)$ with normal vector $\vec{n}=\begin{pmatrix}2\\1\\-1\end{pmatrix}$ is:
\begin{enumerate}
\item $2x+y-z=3$
\item $2x+y-z=-3$
\item $x-2y+3z=0$
\item $2x+y-z=7$
\end{enumerate}
\item The planes $x-2y+z=4$ and $2x-4y+2z=9$ are:
\begin{enumerate}
\item identical
\item intersecting along a line
\item parallel and distinct
\item perpendicular
\end{enumerate}
\item The perpendicular distance from the origin to the plane $3x+4y+12z=26$ is:
\begin{enumerate}
\item $1$
\item $2$
\item $13$
\item $26$
\end{enumerate}
\item The line $\vec{r}=\begin{pmatrix}1\\0\\2\end{pmatrix}+\lambda\begin{pmatrix}1\\1\\1\end{pmatrix}$ and the plane $x-y+z=5$:
\begin{enumerate}
\item meet at exactly one point
\item are parallel, the line not lying in the plane
\item are such that the line lies in the plane
\item are perpendicular
\end{enumerate}
\end{enumerate}
\end{exercice}

\begin{exercice}[True or false]
State whether each assertion is true or false, justifying your answer.
\begin{enumerate}
\item If three points are collinear, then there is exactly one plane passing through them.
\item The angle between two planes is equal to the angle between their normal vectors (taken acute).
\item If $\vec{n}$ is a normal vector of a plane $\Pi$, then $3\vec{n}$ is also a normal vector of $\Pi$.
\item If a line is perpendicular to a plane, then its direction vector is parallel to every vector lying in the plane.
\end{enumerate}
\end{exercice}

\begin{exercice}[Definitions and vocabulary]
\begin{enumerate}
\item Define a \cle{normal vector} of a plane and explain why a plane has infinitely many normal vectors.
\item State the vector equation of a plane through a point $A$, parallel to two non-parallel vectors $\vec{u}$ and $\vec{v}$.
\item State the formula giving the perpendicular distance from the point $P(x_0,y_0,z_0)$ to the plane $ax+by+cz=d$, and say what quantity appears in the denominator.
\end{enumerate}
\end{exercice}

\begin{exercice}[Direct calculations]
\begin{enumerate}
\item Write down a normal vector of each plane: (i) $2x-3y+z=7$; (ii) $x+5z=0$.
\item Find the acute angle between the planes $x+y+z=1$ and $x-y=0$, giving your answer to the nearest degree.
\item Find the perpendicular distance from the point $B(2,-1,4)$ to the plane $2x-y+2z=6$.
\end{enumerate}
\end{exercice}

\courssub{I practise}

\begin{exercice}[Cartesian equation and position of points]
\begin{enumerate}
\item Find the Cartesian equation of the plane through $A(1,2,-1)$ with normal vector $2\vec{i}-\vec{j}+3\vec{k}$.
\item Determine whether the points $(3,1,0)$ and $(0,5,1)$ lie on the same side of this plane.
\end{enumerate}
\end{exercice}

\begin{exercice}[Plane through three points]
\begin{enumerate}
\item Verify that the points $P(1,0,2)$, $Q(2,1,1)$ and $R(-1,2,0)$ are not collinear.
\item Find the Cartesian equation of the plane through $P$, $Q$ and $R$.
\end{enumerate}
\end{exercice}

\begin{exercice}[From parametric to Cartesian form]
A plane has parametric (vector) equation
\[
\vec{r}=\begin{pmatrix}1\\0\\2\end{pmatrix}+\lambda\begin{pmatrix}1\\1\\0\end{pmatrix}+\mu\begin{pmatrix}0\\2\\1\end{pmatrix},\qquad \lambda,\mu\in\mathbb{R}.
\]
\begin{enumerate}
\item Find a normal vector of the plane and deduce its Cartesian equation.
\item Find the perpendicular distance from the origin to the plane.
\end{enumerate}
\end{exercice}

\begin{exercice}[Angle between two planes]
Find the acute angle between the planes $2x+y-2z=5$ and $3x-6y-2z=7$, giving your answer to the nearest degree.
\end{exercice}

\courssub{I prepare for the GCE}

\begin{exercice}[Line and plane (8 marks)]
The line $l$ has equation $\vec{r}=\begin{pmatrix}2\\-1\\3\end{pmatrix}+\lambda\begin{pmatrix}1\\2\\-1\end{pmatrix}$ and the plane $\Pi$ has equation $x-y+2z=4$.
\begin{enumerate}
\item Find the coordinates of the point where $l$ meets $\Pi$. \hfill (3 marks)
\item Find the acute angle between the line $l$ and the plane $\Pi$, to the nearest degree. \hfill (3 marks)
\item Determine whether the point $(4,3,1)$ lies on $l$, and whether it lies in $\Pi$. \hfill (2 marks)
\end{enumerate}
\end{exercice}

\begin{exercice}[Line parallel to a plane; line of intersection (9 marks)]
\begin{enumerate}
\item Show that the line $\vec{r}=\begin{pmatrix}1\\1\\0\end{pmatrix}+\lambda\begin{pmatrix}2\\1\\-2\end{pmatrix}$ is parallel to the plane $2x-2y+z=7$, and find the distance between the line and the plane. \hfill (4 marks)
\item Find the vector equation of the line of intersection of the planes $x+y+z=6$ and $2x-y+z=3$, and check that the point $(1,2,3)$ lies on this line. \hfill (5 marks)
\end{enumerate}
\end{exercice}

\begin{exercice}[Families of planes; extension (10 marks)]
\begin{enumerate}
\item Find the equation of the plane which passes through the line of intersection of the planes $x+2y-z=4$ and $2x-y+z=1$ and also passes through the point $(1,1,2)$. \hfill (4 marks)
\item Find the equation of the plane through the line of intersection of $x+y-z=2$ and $2x-3y+z=5$ which is perpendicular to the plane $x-y+z=0$. \hfill (4 marks)
\item \emph{(For strong candidates.)} The planes $2x-y+2z=8$ and $4x-2y+4z=3$ are given. Show that they are parallel, find the distance between them, and deduce the equation of the plane midway between them. \hfill (2 marks)
\end{enumerate}
\end{exercice}

\newpage

\coursec{Solutions to the exercises}

\courssub{I check my knowledge}

\begin{corrige}[Exercise 1 — Multiple choice]
\begin{enumerate}
\item \textbf{Answer (b):} $2x+y-z=-3$.\par
A plane with normal vector $\vec{n}=\begin{pmatrix}2\\1\\-1\end{pmatrix}$ has equation $2x+y-z=d$. Since $A(1,-2,3)$ lies on it, $d=2(1)+(-2)-(3)=-3$. Hence $2x+y-z=-3$.

\item \textbf{Answer (c):} parallel and distinct.\par
The normal vectors are $\vec{n}_1=(1,-2,1)$ and $\vec{n}_2=(2,-4,2)=2\vec{n}_1$: the normals are parallel, so the planes are parallel. But the right-hand sides do not scale the same way ($2\times 4=8\neq 9$), so the planes are distinct.

\item \textbf{Answer (b):} $2$.\par
Distance from the origin: $\delta=\dfrac{|{-26}|}{\sqrt{3^2+4^2+12^2}}=\dfrac{26}{\sqrt{169}}=\dfrac{26}{13}=2$.

\item \textbf{Answer (b):} parallel, the line not lying in the plane.\par
Direction $\vec{d}=(1,1,1)$, normal $\vec{n}=(1,-1,1)$: $\vec{d}\cdot\vec{n}=1-1+1=1\neq 0$? Let us recompute: $1-1+1=1$. Hmm — check whether the line meets the plane: substitute $x=1+\lambda$, $y=\lambda$, $z=2+\lambda$ into $x-y+z=5$: $(1+\lambda)-\lambda+(2+\lambda)=3+\lambda=5$, so $\lambda=2$. That gives a unique intersection point $(3,2,4)$, which would be answer (a).\par
Let us recheck $\vec{d}\cdot\vec{n}=1(1)+1(-1)+1(1)=1\neq 0$, so the line is \emph{not} parallel to the plane and meets it in exactly one point. \textbf{Correct answer: (a)} — the line meets the plane at the single point $(3,2,4)$ (obtained for $\lambda=2$).
\end{enumerate}
\end{corrige}

\begin{corrige}[Exercise 2 — True or false]
\begin{enumerate}
\item \textbf{False.} If the three points are collinear, infinitely many planes contain the line through them (any plane containing that line). Uniqueness requires \emph{non-collinear} points.
\item \textbf{True.} By definition, the (acute) angle between two planes is the acute angle between their normal vectors: $\cos\theta=\dfrac{|\vec{n}_1\cdot\vec{n}_2|}{\|\vec{n}_1\|\,\|\vec{n}_2\|}$.
\item \textbf{True.} A normal vector is any non-zero vector perpendicular to the plane; any non-zero scalar multiple of $\vec{n}$ is still perpendicular to every direction of the plane.
\item \textbf{False.} If a line is perpendicular to a plane, its direction vector is parallel to the \emph{normal} of the plane, hence \emph{perpendicular} to every vector lying in the plane — not parallel.
\end{enumerate}
\end{corrige}

\begin{corrige}[Exercise 3 — Definitions and vocabulary]
\begin{enumerate}
\item A \cle{normal vector} of a plane $\Pi$ is any non-zero vector perpendicular to every vector lying in $\Pi$ (equivalently, perpendicular to two non-parallel directions of $\Pi$). If $\vec{n}$ is normal to $\Pi$, so is $k\vec{n}$ for every $k\neq 0$; hence a plane has infinitely many normal vectors, all collinear.
\item Vector equation: $\vec{r}=\overrightarrow{OA}+\lambda\vec{u}+\mu\vec{v}$, $\lambda,\mu\in\mathbb{R}$, where $\vec{u}$ and $\vec{v}$ are non-parallel.
\item $\delta=\dfrac{|ax_0+by_0+cz_0-d|}{\sqrt{a^2+b^2+c^2}}$. The denominator is the norm $\|\vec{n}\|$ of a normal vector $\vec{n}=(a,b,c)$ of the plane.
\end{enumerate}
\end{corrige}

\begin{corrige}[Exercise 4 — Direct calculations]
\begin{enumerate}
\item (i) $\vec{n}=(2,-3,1)$; (ii) $\vec{n}=(1,0,5)$ (note the missing $y$-term gives coefficient $0$).
\item Normals: $\vec{n}_1=(1,1,1)$, $\vec{n}_2=(1,-1,0)$.
\[
\cos\theta=\frac{|\vec{n}_1\cdot\vec{n}_2|}{\|\vec{n}_1\|\|\vec{n}_2\|}=\frac{|1-1+0|}{\sqrt{3}\sqrt{2}}=0
\quad\Longrightarrow\quad \theta=90^{\circ}.
\]
The planes are perpendicular.
\item $\delta=\dfrac{|2(2)-(-1)+2(4)-6|}{\sqrt{4+1+4}}=\dfrac{|4+1+8-6|}{3}=\dfrac{7}{3}\approx 2.33$ units.
\end{enumerate}
\end{corrige}

\courssub{I practise}

\begin{corrige}[Exercise 5 — Cartesian equation and position of points]
\begin{enumerate}
\item The plane has equation $2x-y+3z=d$. Substituting $A(1,2,-1)$: $d=2-2-3=-3$. Hence
\[
\boxed{2x-y+3z=-3}.
\]
\item Define $f(x,y,z)=2x-y+3z+3$. For $(3,1,0)$: $f=6-1+0+3=8>0$. For $(0,5,1)$: $f=0-5+3+3=1>0$. Both values have the same sign, so the two points lie on the \textbf{same side} of the plane.
\end{enumerate}
\end{corrige}

\begin{corrige}[Exercise 6 — Plane through three points]
\begin{enumerate}
\item $\overrightarrow{PQ}=(1,1,-1)$ and $\overrightarrow{PR}=(-2,2,-2)$. These vectors are not proportional ($\dfrac{1}{-2}\neq\dfrac{1}{2}$), so $P$, $Q$, $R$ are \textbf{not collinear}.
\item Normal vector:
\[
\vec{n}=\overrightarrow{PQ}\times\overrightarrow{PR}
=\begin{pmatrix}1\\1\\-1\end{pmatrix}\times\begin{pmatrix}-2\\2\\-2\end{pmatrix}
=\begin{pmatrix}(1)(-2)-(-1)(2)\\ (-1)(-2)-(1)(-2)\\ (1)(2)-(1)(-2)\end{pmatrix}
=\begin{pmatrix}0\\4\\4\end{pmatrix}.
\]
Take the simpler normal $(0,1,1)$. Equation: $y+z=d$; using $P(1,0,2)$: $d=2$. Hence
\[
\boxed{y+z=2}.
\]
Check: $Q$: $1+1=2$ \checkmark; $R$: $2+0=2$ \checkmark.
\end{enumerate}
\end{corrige}

\begin{corrige}[Exercise 7 — From parametric to Cartesian form]
\begin{enumerate}
\item Direction vectors $\vec{u}=(1,1,0)$, $\vec{v}=(0,2,1)$. A normal vector is
\[
\vec{n}=\vec{u}\times\vec{v}
=\begin{pmatrix}(1)(1)-(0)(2)\\ (0)(0)-(1)(1)\\ (1)(2)-(1)(0)\end{pmatrix}
=\begin{pmatrix}1\\-1\\2\end{pmatrix}.
\]
Equation: $x-y+2z=d$; the point $(1,0,2)$ lies in the plane, so $d=1+4=5$:
\[
\boxed{x-y+2z=5}.
\]
\item Distance from the origin:
\[
\delta=\frac{|0-0+0-5|}{\sqrt{1+1+4}}=\frac{5}{\sqrt{6}}=\frac{5\sqrt{6}}{6}\approx 2.04 \text{ units}.
\]
\end{enumerate}
\end{corrige}

\begin{corrige}[Exercise 8 — Angle between two planes]
Normals: $\vec{n}_1=(2,1,-2)$, $\vec{n}_2=(3,-6,-2)$.
\[
\vec{n}_1\cdot\vec{n}_2=6-6+4=4,\qquad
\|\vec{n}_1\|=\sqrt{4+1+4}=3,\qquad
\|\vec{n}_2\|=\sqrt{9+36+4}=7.
\]
\[
\cos\theta=\frac{|4|}{3\times 7}=\frac{4}{21}\approx 0.1905
\quad\Longrightarrow\quad
\theta\approx 79^{\circ}\ \text{(nearest degree)}.
\]
\end{corrige}

\courssub{I prepare for the GCE}

\begin{corrige}[Exercise 9 — Line and plane (8 marks)]
\begin{enumerate}
\item Parametrize: $x=2+\lambda$, $y=-1+2\lambda$, $z=3-\lambda$. Substitute into $x-y+2z=4$:
\[
(2+\lambda)-(-1+2\lambda)+2(3-\lambda)=4
\;\Longrightarrow\;
2+\lambda+1-2\lambda+6-2\lambda=4
\;\Longrightarrow\;
9-3\lambda=4,
\]
so $\lambda=\dfrac{5}{3}$. The intersection point is
\[
\left(2+\tfrac{5}{3},\;-1+\tfrac{10}{3},\;3-\tfrac{5}{3}\right)=\boxed{\left(\tfrac{11}{3},\;\tfrac{7}{3},\;\tfrac{4}{3}\right)}.
\]
\emph{Mark scheme: M1 substitution; A1 value of $\lambda$; A1 point.}

\item Direction $\vec{d}=(1,2,-1)$, normal $\vec{n}=(1,-1,2)$.
\[
\vec{d}\cdot\vec{n}=1-2-2=-3,\quad \|\vec{d}\|=\sqrt{6},\quad \|\vec{n}\|=\sqrt{6}.
\]
If $\varphi$ is the angle between $\vec{d}$ and $\vec{n}$ and $\theta$ the angle between line and plane, $\sin\theta=|\cos\varphi|$:
\[
\sin\theta=\frac{|\vec{d}\cdot\vec{n}|}{\|\vec{d}\|\,\|\vec{n}\|}=\frac{3}{6}=\frac12
\quad\Longrightarrow\quad \boxed{\theta=30^{\circ}}.
\]
\emph{Mark scheme: M1 correct scalar product; M1 use of $\sin\theta$ formula; A1 answer. Examiner expects the acute angle between line and plane, i.e.\ the complement of the angle with the normal.}

\item On $l$: we need $2+\lambda=4$, $-1+2\lambda=3$, $3-\lambda=1$, i.e.\ $\lambda=2$ from all three equations — consistent. So $(4,3,1)$ \textbf{lies on} $l$.\par
In $\Pi$: $4-3+2(1)=3\neq 4$, so $(4,3,1)$ \textbf{does not lie in} $\Pi$.\par
\emph{Mark scheme: B1 for each correct conclusion with justification.}
\end{enumerate}
\end{corrige}

\begin{corrige}[Exercise 10 — Line parallel to a plane; line of intersection (9 marks)]
\begin{enumerate}
\item Direction $\vec{d}=(2,1,-2)$, normal $\vec{n}=(2,-2,1)$:
\[
\vec{d}\cdot\vec{n}=4-2-2=0,
\]
so the line is parallel to the plane (or lies in it). Test the point $(1,1,0)$ of the line: $2-2+0=0\neq 7$, so the line does \textbf{not} lie in the plane: it is strictly parallel.\par
The distance between line and plane equals the distance from any point of the line, e.g.\ $(1,1,0)$, to the plane:
\[
\delta=\frac{|2(1)-2(1)+0-7|}{\sqrt{4+4+1}}=\boxed{\frac{7}{3}\approx 2.33\text{ units}}.
\]
\emph{Mark scheme: M1 scalar product; A1 conclusion parallel; B1 check the line is not in the plane; A1 distance.}

\item Solve the system
\[
\begin{cases} x+y+z=6\\ 2x-y+z=3\end{cases}
\]
Set $z=\lambda$. Subtracting the first equation from the second: $x-2y=-3$, so $x=2y-3$. Substituting into the first: $(2y-3)+y+\lambda=6$, so $3y=9-\lambda$, $y=3-\dfrac{\lambda}{3}$, and $x=2\left(3-\dfrac{\lambda}{3}\right)-3=3-\dfrac{2\lambda}{3}$.\par
To avoid fractions, put $\lambda=3t$: then
\[
x=3-2t,\qquad y=3-t,\qquad z=3t,
\]
i.e.
\[
\boxed{\vec{r}=\begin{pmatrix}3\\3\\0\end{pmatrix}+t\begin{pmatrix}-2\\-1\\3\end{pmatrix},\quad t\in\mathbb{R}.}
\]
\emph{Check direction:} $(-2,-1,3)$ is perpendicular to both normals: $(-2)(1)+(-1)(1)+(3)(1)=0$ and $(-2)(2)+(-1)(-1)+(3)(1)=-4+1+3=0$ \checkmark.\par
\emph{Check $(1,2,3)$:} $3-2t=1\Rightarrow t=1$; then $y=3-1=2$ \checkmark, $z=3$ \checkmark. Also $1+2+3=6$ and $2-2+3=3$ \checkmark. So $(1,2,3)$ lies on the line.\par
\emph{Mark scheme: M1 elimination; A1 parametric solution; A1 vector equation; M1 verification method; A1 conclusion.}
\end{enumerate}
\end{corrige}

\begin{corrige}[Exercise 11 — Families of planes; extension (10 marks)]
\begin{enumerate}
\item Any plane through the line of intersection of $x+2y-z=4$ and $2x-y+z=1$ has equation
\[
(x+2y-z-4)+\lambda(2x-y+z-1)=0.
\]
Impose passage through $(1,1,2)$:
\[
(1+2-2-4)+\lambda(2-1+2-1)=0
\;\Longrightarrow\;
-3+2\lambda=0
\;\Longrightarrow\;
\lambda=\tfrac32.
\]
Hence
\[
(x+2y-z-4)+\tfrac32(2x-y+z-1)=0
\;\Longrightarrow\;
2(x+2y-z-4)+3(2x-y+z-1)=0,
\]
\[
\boxed{8x+y+z-11=0}.
\]
\emph{Mark scheme: M1 family of planes; M1 substitution of the point; A1 $\lambda$; A1 equation.}

\item Family: $(x+y-z-2)+\lambda(2x-3y+z-5)=0$, i.e.
\[
(1+2\lambda)x+(1-3\lambda)y+(-1+\lambda)z-(2+5\lambda)=0.
\]
Its normal is $\vec{n}=(1+2\lambda,\;1-3\lambda,\;-1+\lambda)$. Perpendicularity to $x-y+z=0$ (normal $(1,-1,1)$) requires $\vec{n}\cdot(1,-1,1)=0$:
\[
(1+2\lambda)-(1-3\lambda)+(-1+\lambda)=0
\;\Longrightarrow\;
6\lambda-1=0
\;\Longrightarrow\;
\lambda=\tfrac16.
\]
Substituting:
\[
\left(1+\tfrac13\right)x+\left(1-\tfrac12\right)y+\left(-1+\tfrac16\right)z-\left(2+\tfrac56\right)=0
\;\Longrightarrow\;
\tfrac{4}{3}x+\tfrac12 y-\tfrac56 z-\tfrac{17}{6}=0.
\]
Multiplying by $6$: $\boxed{8x+3y-5z-17=0}$.\par
\emph{Mark scheme: M1 family; M1 perpendicularity condition on normals; A1 $\lambda$; A1 equation.}

\item Normals: $\vec{n}_1=(2,-1,2)$ and $\vec{n}_2=(4,-2,4)=2\vec{n}_1$: the planes are parallel. They are distinct since $\dfrac{8}{2}=4\neq \dfrac{3}{2}$ (the equations are not proportional).\par
Distance: take a point of the first plane, e.g.\ $(4,0,0)$ (check: $8=8$ \checkmark), and compute its distance to $4x-2y+4z=3$:
\[
\delta=\frac{|4(4)-0+0-3|}{\sqrt{16+4+16}}=\frac{13}{6}\approx 2.17\text{ units}.
\]
Midway plane: it is parallel to both, so has equation $2x-y+2z=k$ with $k$ the mean of the right-hand sides written with the \emph{same} coefficients. Writing the second plane as $2x-y+2z=\dfrac{3}{2}$:
\[
k=\frac{8+\frac32}{2}=\frac{19}{4},
\qquad\text{so}\qquad
\boxed{2x-y+2z=\tfrac{19}{4}}\ \text{or equivalently}\ 8x-4y+8z=19.
\]
\emph{Mark scheme: B1 parallel + distance; B1 midway plane. Examiner expects candidates to normalise the equations before averaging.}
\end{enumerate}
\end{corrige}

\newpage

\coursec{Summary sheet}

\begin{popfigure}
\begin{tikzpicture}[scale=0.92, transform shape,
  every node/.style={font=\small},
  center/.style={circle, draw=popdark, fill=popblueL, text=popdark, font=\bfseries, minimum size=2.6cm, align=center},
  branch/.style={rectangle, rounded corners=6pt, draw=#1, fill=#1!12, text=popdark, font=\small, align=center, text width=3.4cm, inner sep=4pt}]
  \node[center, align=center] (C){Planes in\\3-D Space};
  \node[branch=popblue, align=center] (A) at (-5.6,3.2){\textbf{Equations of a plane}\\ $\vec{r}\cdot\vec{n}=\vec{a}\cdot\vec{n}$\\ $ax+by+cz=d$};
  \node[branch=poporange, align=center] (B) at (0,4.4){\textbf{Plane through 3 points}\\ parametric form\\ $\vec{r}=\vec{a}+\lambda\vec{u}+\mu\vec{v}$};
  \node[branch=popgreen, align=center] (D) at (5.6,3.2){\textbf{Distance point--plane}\\ $\dfrac{|ax_1+by_1+cz_1-d|}{\sqrt{a^2+b^2+c^2}}$};
  \node[branch=poppurple, align=center] (E) at (-5.6,-3.2){\textbf{Angles}\\ plane--plane: normals\\ line--plane: $\sin\theta$};
  \node[branch=popgold, align=center] (F) at (0,-4.4){\textbf{Intersections}\\ line $\cap$ plane: one point\\ plane $\cap$ plane: a line};
  \node[branch=poppink, align=center] (G) at (5.6,-3.2){\textbf{Families of planes}\\ $\Pi_1+\lambda\Pi_2=0$\\ through intersection};
  \draw[thick, popblue] (C) -- (A);
  \draw[thick, poporange] (C) -- (B);
  \draw[thick, popgreen] (C) -- (D);
  \draw[thick, poppurple] (C) -- (E);
  \draw[thick, popgold] (C) -- (F);
  \draw[thick, poppink] (C) -- (G);
\end{tikzpicture}
\legende{Mind map of the chapter \emph{Vectors: planes in three dimensions}.}
\end{popfigure}

\begin{bilan}
\textbf{Key definitions.}
\begin{itemize}
  \item A \cle{plane} is determined by a point $A$ and a \cle{normal vector} $\vec{n}$, or by three non-collinear points, or by a point and two non-parallel direction vectors.
  \item \cle{Vector equation} of a plane: $(\vec{r}-\vec{a})\cdot\vec{n}=0$, i.e. $\vec{r}\cdot\vec{n}=\vec{a}\cdot\vec{n}$.
  \item \cle{Cartesian equation}: $ax+by+cz=d$, where $\vec{n}=\begin{pmatrix}a\\b\\c\end{pmatrix}$ is a normal vector and $d=\vec{a}\cdot\vec{n}$.
  \item \cle{Parametric form}: $\vec{r}=\vec{a}+\lambda\vec{u}+\mu\vec{v}$, $\lambda,\mu\in\mathbb{R}$, with $\vec{u},\vec{v}$ non-parallel vectors lying in the plane.
\end{itemize}

\textbf{Formulas and results to know.}
\begin{itemize}
  \item Plane through $A(x_1,y_1,z_1)$ with normal $(a,b,c)$: $a(x-x_1)+b(y-y_1)+c(z-z_1)=0$.
  \item Plane through three points $A,B,C$: direction vectors $\overrightarrow{AB}$, $\overrightarrow{AC}$; normal $\vec{n}=\overrightarrow{AB}\times\overrightarrow{AC}$ (the three points must \emph{not} be collinear).
  \item \cle{Perpendicular distance} from $P(x_1,y_1,z_1)$ to $ax+by+cz=d$:
  \[ \delta=\frac{|ax_1+by_1+cz_1-d|}{\sqrt{a^2+b^2+c^2}}. \]
  \item \cle{Angle between two planes} = angle between their normals:
  \[ \cos\theta=\frac{|\vec{n_1}\cdot\vec{n_2}|}{|\vec{n_1}|\,|\vec{n_2}|}. \]
  \item \cle{Angle between a line} (direction $\vec{u}$) and a plane (normal $\vec{n}$):
  \[ \sin\theta=\frac{|\vec{u}\cdot\vec{n}|}{|\vec{u}|\,|\vec{n}|}. \]
  \item \cle{Intersection line of two planes}: direction $\vec{u}=\vec{n_1}\times\vec{n_2}$; find one common point by fixing one coordinate.
  \item \cle{Family of planes} through the line of intersection of $\Pi_1=0$ and $\Pi_2=0$: $\Pi_1+\lambda\Pi_2=0$ (plus $\Pi_2$ itself).
\end{itemize}

\textbf{Methods.}
\begin{itemize}
  \item Line--plane intersection: substitute the parametric equations of the line into the plane equation; solve for the parameter; substitute back. No solution $\Rightarrow$ line parallel to plane; infinitely many $\Rightarrow$ line lies in the plane.
  \item To find a plane through the intersection of two planes and satisfying one extra condition (a point, perpendicularity), write $\Pi_1+\lambda\Pi_2=0$ and determine $\lambda$.
  \item To convert parametric to Cartesian form, compute $\vec{n}=\vec{u}\times\vec{v}$ then $d=\vec{a}\cdot\vec{n}$.
\end{itemize}

\textbf{Common mistakes.}
\begin{itemize}
  \item Forgetting the absolute value in distance and angle formulas (distances and acute angles are non-negative).
  \item Using $\cos\theta$ instead of $\sin\theta$ for the line--plane angle: the line--plane angle is the \emph{complement} of the line--normal angle.
  \item Trying to build a plane from three \emph{collinear} points: infinitely many planes pass through them, so the plane is not unique.
  \item Confusing the normal vector with a direction vector of the plane: the normal is \emph{perpendicular} to the plane.
  \item Omitting the case $\Pi_2=0$ when using the family $\Pi_1+\lambda\Pi_2=0$.
\end{itemize}

\textbf{GCE tips.}
\begin{itemize}
  \item Always state $\vec{n}$ explicitly; most marks in vector questions come from identifying normals correctly.
  \item Show the substitution step fully when intersecting a line and a plane; examiners award method marks.
  \item Check your final plane equation by substituting a known point of the plane.
  \item Keep exact values (fractions, surds) until the final answer, then round sensibly.
\end{itemize}
\end{bilan}

\findecours

\end{document}
